\medterm{G6PD Deficiency} An inherited condition in which red blood cells can break down after infections, fava beans, or certain medicines. Episodes may cause anemia, jaundice, fatigue, shortness of breath, or dark urine.

\medterm{Gabapentin} A medicine used for certain seizures and nerve-related pain. It can cause sleepiness, dizziness, swelling, or impaired coordination; the dose often needs adjustment in kidney disease and stopping suddenly can cause problems.

\medterm{Gaba Receptor} A brain-cell protein that responds to GABA, a chemical messenger that usually reduces nerve activity. GABA receptors help regulate sleep, anxiety, muscle control, and seizures and are targets of several medicines.

\medterm{GAD} An abbreviation for generalized anxiety disorder, involving excessive, difficult-to-control worry about several areas of life for at least six months, with symptoms such as restlessness, tension, fatigue, poor concentration, or sleep disturbance.

\synonyms
Generalized Anxiety Disorder

\medterm{Gadolinium} A metal used in a chemically bound form to improve some MRI images. Serious reactions are uncommon, but kidney function, pregnancy, prior reactions, and possible tissue retention are considered before use.

\medterm{Gadolinium-Based Contrast Agent} A substance injected during some MRI scans to make organs, blood vessels, or abnormal tissue easier to see. The agent and dose are selected according to kidney function, pregnancy, allergies, and the scan’s purpose.

\medterm{Gadolinium-Induced Nephrogenic Systemic Fibrosis} A rare illness causing thick, tight, hardened skin and sometimes damage to internal organs after certain MRI contrast exposure, mainly in severe kidney disease. Safer contrast agents have made it very uncommon.

\medterm{Gadopentetate Dimeglumine} An injectable gadolinium contrast agent used to improve selected MRI images. It can cause allergic reactions and is used cautiously or avoided in severe kidney disease because of rare serious complications.

\medterm{Gag Gene} A retroviral gene that encodes structural proteins forming the virus’s inner support, shell, and core. These proteins help assemble new virus particles inside infected cells.

\medterm{Gagging} Retching caused by stimulation of the throat or back of the tongue. It may occur with swallowing difficulty, reflux, infection, examination, aspiration, or distress.

\medterm{Gag Protein} A large retroviral protein that is cut into components forming the virus’s inner support, protective shell, and genetic-material core. These parts are required to assemble new virus particles.

\medterm{Gag Reflex} A protective response in which touching the back of the throat causes throat-muscle contraction and sometimes retching. It helps prevent food or objects from entering the airway.

\synonyms
Pharyngeal Reflex

\medterm{Gain} The amount by which a hearing aid or amplifier increases sound intensity, measured in decibels. Correct gain improves hearing while limiting excessive loudness, feedback, discomfort, and further hearing damage.

\medterm{Gain-Of-Function Mutation} A genetic change that makes a gene or its protein more active, gives it a new activity, or causes activity at the wrong time or place. Its effects range from harmless variation to inherited disease or cancer.

\medterm{Gaisbock Syndrome} A form of relative or stress polycythemia characterized by an elevated hematocrit resulting from decreased plasma volume rather than increased red cell mass, typically occurring in middle-aged, hypertensive, anxious, and overweight individuals.
\synonyms
Stress Polycythemia; Polycythemia Hypertonica

\medterm{Gait} A person’s pattern of walking, including speed, balance, step length, posture, arm movement, and turning. Gait reflects the combined function of the brain, nerves, muscles, joints, vision, and balance system.

\medterm{Gait Abnormality} A change from a person’s usual walking pattern caused by problems with muscles, joints, nerves, balance, vision, heart function, or general health. Assessment examines speed, stability, symmetry, strength, sensation, and pain.

\medterm{Gait Apraxia} Difficulty starting or carrying out walking even though basic muscle strength is adequate. It may occur with disease or injury affecting the front of the brain and can cause short steps, freezing, or instability.

\medterm{Gait Assessment} An examination of how a person walks, including speed, steps, balance, posture, turning, and use of aids. It can reveal problems with strength, sensation, coordination, joints, vision, or thinking.

\medterm{Gait Ataxia} Unsteady, poorly coordinated walking caused by problems in the cerebellum, sensation pathways, balance organs, or other parts of the nervous system.

\medterm{Gait Cycle} The complete sequence of walking movements from one heel strike to the next heel strike of the same foot. It includes the stance phase when the foot bears weight and the swing phase when it moves forward.

\medterm{Gait Disturbance} An abnormal walking pattern involving speed, balance, coordination, strength, or step movement. Causes include neurologic disease, joint or muscle problems, inner-ear disorders, pain, medicines, and systemic illness.

\medterm{Gaits} Patterns of walking. Different gaits may reflect pain, weakness, joint disease, balance problems, or neurologic conditions; assessment considers the person’s usual walking pattern and associated findings.

\medterm{Gait Speed} Walking speed measured over a set distance. It is a simple indicator of mobility and physical reserve, helps estimate fall or disability risk, and can track changes during illness or rehabilitation.

\medterm{Gait Training} Therapy that teaches or retrains safer walking after injury, illness, weakness, or balance problems. It may use exercises, parallel bars, walkers, canes, stairs, or other aids to improve independence.

\medterm{Galactocele} A milk-filled breast cyst caused by blockage of a milk duct, usually during breastfeeding or after weaning. It commonly feels like a painless lump; examination and ultrasound distinguish it from other breast masses.

\medterm{Galactocerebrosidase Deficiency} An inherited lack of an enzyme needed to break down certain fats in nerve tissue. Severe deficiency causes Krabbe disease, which progressively damages the brain and nerves.

\synonyms
Krabbe Disease; Galactosylceramidase Deficiency

\medterm{Galactogogue} A medicine, herb, food, or other substance claimed to increase breast-milk production. Evidence and safety vary, and effects may depend on feeding technique, milk removal, medical causes, and interactions.

\medterm{Galactokinase Deficiency} An autosomal recessive inborn error of galactose metabolism caused by mutations in the \textit{GALK1} gene, which encodes galactokinase. The resulting accumulation of galactose in tissues leads to reduction by aldose reductase into galactitol in the crystalline lens, causing the progressive development of bilateral infantile cataracts without the hepatic, renal, or cognitive damage seen in classic galactosemia.

\synonyms
GALK Deficiency; Type II Galactosemia

\medterm{Galactorrhea} The inappropriate or spontaneous secretion of milk or milk-like fluid from one or both breasts unrelated to normal pregnancy or breastfeeding. It is most commonly caused by hyperprolactinaemia secondary to pituitary prolactinomas, dopamine-antagonist medications (such as antipsychotics and antiemetics), primary hypothyroidism (via elevated thyrotropin-releasing hormone), chronic kidney disease, or local chest wall stimulation.

\synonyms
Galactorrhoea; Non-Puerperal Lactation; Hyperprolactinemic Discharge

\medterm{Galactorrhoea} Alternate name; see \textit{Galactorrhea}.

\medterm{Galactosaemia} A rare inherited disorder in which the body cannot properly process galactose, a sugar in milk. Untreated infants may develop feeding problems, jaundice, liver injury, cataracts, serious infection, or developmental problems.

\medterm{Galactose} A simple sugar that forms part of lactose in milk. The body normally processes it for energy, but inherited galactosemia can make milk and other galactose-containing foods dangerous for infants.

\medterm{Galactose-1-Phosphate Uridyltransferase Blood Test} A blood test measuring an enzyme needed to process galactose. It helps screen for or diagnose classic galactosemia, especially in newborns or infants with symptoms after milk feeding.

\medterm{Galactose Metabolism} The body’s process for changing galactose, part of milk sugar, into energy or other useful substances. Inherited enzyme problems can block this process and cause galactosemia.

\medterm{Galactosemia} An inherited disorder in which the body cannot properly process galactose. Classic disease can cause vomiting, jaundice, liver injury, cataracts, serious infection, and developmental problems after milk feeding.

\medterm{Galactosialidosis} A rare inherited storage disorder in which cells accumulate certain sugar-containing substances. It can cause developmental delay, muscle weakness, bone changes, kidney disease, and eye problems, with severity varying widely.

\medterm{Galactoside} A compound in which the sugar galactose is joined to another molecule, such as another sugar, a fat, or a protein. Galactosides occur in cells and in some foods.

\medterm{Galactosyltransferase} An enzyme that attaches galactose to proteins, fats, or other sugars while cells build complex molecules. These molecules support cell structure, recognition, signaling, and nerve insulation.

\medterm{Galant Reflex} A newborn reflex in which stroking one side of the back causes the infant’s trunk to curve toward that side. It normally fades during early infancy; marked asymmetry or persistence can occur with nervous-system problems.

\medterm{Galeazzi Fracture} A fracture of the middle to distal third of the radial shaft associated with dislocation or subluxation of the distal radioulnar joint (DRUJ), typically requiring anatomic open reduction and internal fixation to restore forearm pronation and supination.
\synonyms
Galeazzi Fracture-Dislocation; Piedmont Fracture

\medterm{Galeazzi Fracture-Dislocation} A forearm injury involving a break in the lower part of the radius and dislocation of the nearby joint between the radius and ulna. It can cause wrist pain, instability, and difficulty turning the forearm.

\synonyms
Galeazzi Fracture; Reverse Monteggia Fracture

\medterm{Galeazzi Sign} A finding in an infant lying with hips and knees bent, in which one knee is lower than the other. It may suggest hip dislocation, developmental hip dysplasia, or unequal thigh-bone length.

\synonyms
Allis Sign; Galeazzi Test

\medterm{Gallavardin Phenomenon} A finding in severe aortic-valve narrowing where the heart murmur sounds harsh near the upper chest but musical near the apex. It reflects different sound frequencies transmitted through the chest.

\medterm{Gall-Bladder} A small pear-shaped sac beneath the liver that stores and concentrates bile. Bile is released into the small intestine to help digest fats.

\medterm{Gallbladder Absence} Having no gallbladder because it did not develop or was surgically removed. Bile then flows directly from the liver into the intestine; most people can digest normally, though some develop digestive symptoms.

\medterm{Gallbladder Adenomyomatosis} A benign, non-inflammatory hyperplastic alteration of the gallbladder wall characterized by excessive epithelial proliferation with smooth muscle hypertrophy and deep, pouch-like invaginations of the luminal mucosa extending into the muscularis propria (termed Rokitansky-Aschoff sinuses). Often identified on ultrasound by intramural diverticula containing cholesterol crystals that produce characteristic V-shaped \"comet-tail\" reverberation artifacts.

\synonyms
Adenomyoma of Gallbladder; Diverticular Gallbladder Disease

\medterm{Gallbladder Cancer} A cancer arising in the gallbladder, sometimes associated with long-term inflammation or stones. It may cause abdominal pain, jaundice, weight loss, or blocked bile flow, but can also be found unexpectedly after surgery.

\synonyms
Gallbladder Carcinoma; Biliary Malignancy

\medterm{Gallbladder Carcinoma} A cancer of the gallbladder, most often arising from its lining cells. It may be found incidentally or cause right-upper-abdominal pain, jaundice, weight loss, or bile-duct blockage; treatment depends on its stage and spread.

\medterm{Gall Bladder Disease} Any disorder of the gallbladder, including stones, inflammation, infection, polyps, abnormal movement, or cancer. Symptoms may include right-upper-abdominal pain, nausea, vomiting, fever, or jaundice.

\medterm{Gallbladder Diseases} Conditions affecting the gallbladder, including stones, inflammation, infection, polyps, abnormal emptying, and cancer. Symptoms and treatment depend on the exact disorder and whether bile flow is blocked.

\medterm{Gallbladder Duplication} A rare birth variation in which two gallbladders develop, with separate or shared drainage ducts. It may cause no symptoms but can be associated with stones or inflammation and is important to identify before surgery.

\medterm{Gallbladder Hydrops} Severe, non-inflammatory distension and dilation of the gallbladder filled with clear, sterile, watery or mucoid transudate resulting from chronic, complete obstruction of the cystic duct (most often by an impacted gallstone in Hartmann's pouch) in the absence of acute infection.
\synonyms
Mucocele of Gallbladder; Hydrops Vesicae Felleae

\medterm{Gallbladder Infection} Inflammation or infection of the gallbladder, often after a stone blocks bile flow. Right-upper-abdominal pain, fever, nausea, vomiting, and tenderness may occur; treatment can include antibiotics, drainage, or surgery.

\medterm{Gallbladder Inflammation} Inflammation of the gallbladder, most often caused by a gallstone blocking its drainage. It can cause persistent right-upper-abdominal pain, fever, nausea, and vomiting.

\synonyms
Cholecystitis

\medterm{Gallbladder Mucocele} Enlargement of the gallbladder with mucus or watery fluid because its outlet is blocked, often by an impacted stone. It may cause pain or infection; drainage or removal may be used.

\medterm{Gallbladder Neck} The narrow, tapered S-shaped constriction of the gallbladder that connects the gallbladder body to the cystic duct, characterized internally by spiral mucosal folds (valves of Heister); an impacted gallstone here leads to acute cholecystitis or Mirizzi syndrome.

\synonyms
Collum Vesicae Biliaris

\medterm{Gallbladder Necrosis} Death of gallbladder tissue caused by severe inflammation or inadequate blood supply. It can lead to perforation, abdominal infection, and sepsis.

\medterm{Gallbladder Pain} Pain from temporary blockage of the gallbladder outlet, usually by a stone. It is often steady, severe, and felt under the right ribs after meals and may include nausea or vomiting.

\synonyms
Biliary Colic; Gallstone Pain

\medterm{Gallbladder Perforation} A hole through the gallbladder wall that allows bile or infected contents into the abdomen. It can cause severe inflammation, abscess, peritonitis, or sepsis.

\medterm{Gallbladder Polyp} A mucosal projection extending into the lumen of the gallbladder, categorized as pseudopolyps (cholesterol polyps, inflammatory polyps, adenomyomatosis) or true epithelial neoplasms (adenomas); polyps larger than 10 mm carry an increased malignancy risk warranting cholecystectomy.

\medterm{Gallbladder Radionuclide Scan} A nuclear medicine test that follows a tracer through the liver, gallbladder, and bile ducts. It helps assess suspected gallbladder inflammation, blockage, or abnormal bile flow.

\medterm{Gallbladder Volvulus} Twisting of the gallbladder around its supporting tissues and duct, cutting off blood flow. It causes sudden severe right-upper-abdominal pain, nausea, vomiting, and abdominal irritation.

\medterm{Gallium} A metallic element used in some radioactive tracers, medical compounds, and technology. Its medical effects depend on the isotope or compound and how it is administered and distributed in the body.

\medterm{Gallium-68} A radioactive isotope used to label tracers for PET scans. Linked to a suitable targeting molecule, it can help locate some neuroendocrine tumors or prostate-cancer lesions; it is used for diagnosis, not direct treatment.

\medterm{Gallium Nitrate} A gallium-containing compound formerly used in selected cancer settings to reduce high calcium levels. Its use is limited because it can damage the kidneys and has largely been replaced by other treatments.

\medterm{Gallium Scan} A nuclear medicine scan using radioactive gallium to locate some infections, inflammation, or tumors. It is used less often now because newer imaging tests may provide more specific or faster information.

\medterm{Gallop Rhythm} An extra heart sound creating a three- or four-beat rhythm. It may reflect abnormal heart filling, stiff heart muscle, or heart failure and is interpreted with symptoms and other examination findings.

\medterm{Gallstone Ileus} A blockage of the intestine caused by a large gallstone that has entered through an abnormal connection from the gallbladder. It can cause pain, vomiting, swelling, and inability to pass stool or gas.

\medterm{Gallstone Pancreatitis} Sudden inflammation of the pancreas caused by a gallstone blocking the shared opening of the bile and pancreatic ducts. It usually causes severe persistent upper-abdominal pain and may be associated with complications or remaining duct stones.

\medterm{Gallstones} Solid deposits that form in the gallbladder when bile contains too much cholesterol or bilirubin, too little bile salt, or remains in the gallbladder without emptying fully. They may cause no symptoms or block bile flow, producing severe upper-abdominal pain, inflammation, jaundice, infection, or pancreatitis.

\synonyms
Cholelithiasis; Biliary Calculi

\medterm{Galvanocautery} Use of heat generated by electric current to destroy selected tissue or control bleeding during a procedure. It can cause burns or damage nearby structures if not carefully controlled.

\medterm{Gambling Disorder} A persistent inability to control gambling despite financial, relationship, work, health, or emotional harm. The disorder can cause debt, secrecy, distress, and suicide risk and benefits from professional behavioral and mental-health support.

\synonyms
Pathological Gambling; Compulsive Gambling; Ludomania

\medterm{Gamete} A mature reproductive cell, such as a sperm or egg, containing one set of chromosomes. Two gametes combine during fertilization to form a cell with genetic material from both biological parents.

\medterm{Gamete Intrafallopian Transfer} An assisted-reproduction procedure in which eggs and sperm are placed into a fallopian tube so fertilization may occur inside the body. It is now used less often than in-vitro fertilization.

\medterm{Gametogenesis} The formation and maturation of sperm or eggs in the testes or ovaries. It includes cell division that reduces chromosome number so gametes can combine during fertilization.

\medterm{Gamma-Aminobutyric Acid} A chemical messenger that usually reduces nerve activity in the brain and spinal cord. It helps regulate sleep, anxiety, muscle activity, and seizures through two main receptor types.

\synonyms
GABA

\medterm{Gamma-Benzene Hexachloride} An insecticide, also called lindane, that can cause nervous-system toxicity and seizures after excessive exposure. Medical use is restricted in many places because safer treatments are available.

\medterm{Gamma Camera} A device used in some nuclear-medicine scans to detect radiation from a small tracer inside the body and create pictures showing an organ’s structure or activity.

\medterm{Gamma-Chemokines} Signaling proteins that attract or activate immune cells through specific receptors. They help coordinate inflammation, immune surveillance, and movement of defensive cells into affected tissue.

\medterm{Gamma Globulin} The part of blood plasma containing most immunoglobulins, or antibodies, which help defend against infections and can be measured or given as immune treatment.

\medterm{Gamma-Globulin} A group of blood proteins that includes antibodies. The term may refer to antibody-containing plasma fractions or immunoglobulin preparations used to provide temporary immune protection or treatment.

\medterm{Gamma-Glutamyl Transferase Blood Test} A blood test measuring an enzyme found mainly in the liver and bile ducts. It can support evaluation of liver or bile-duct disease and help explain a high alkaline-phosphatase result.

\medterm{Gammaherpesvirinae} A subgroup of herpesviruses that establish long-lasting infection in lymphoid or related cells. Some members, including Epstein-Barr virus, can cause human disease and contribute to certain cancers.

\medterm{Gamma Hydroxybutyrate} A central-nervous-system depressant that can cause drowsiness, confusion, slowed breathing, coma, or death in overdose. It is also misused recreationally and can cause dependence and withdrawal.

\medterm{Gamma-Interferon} A cytokine made by activated immune cells that strengthens macrophage activity and cell-mediated defense against certain infections. It also contributes to inflammation and immune regulation.

\medterm{Gamma Knife} A noninvasive radiation treatment that aims many focused beams at selected brain tumors or other lesions. It avoids an open operation but can cause swelling, bleeding, or injury to nearby brain tissue.

\medterm{Gamma Knife Radiosurgery} Gamma Knife is a focused radiation system that delivers multiple precisely aimed beams, mainly to selected brain tumors and vascular or other brain lesions.

\medterm{Gamma Rays} High-energy electromagnetic radiation produced by radioactive decay or nuclear reactions. It can damage cells and is used in controlled amounts for some cancer treatments and medical sterilization.

\medterm{Gamma-Ray Therapy} Treatment that uses controlled gamma radiation to damage the DNA of selected abnormal cells. It can treat some cancers, but nearby healthy tissue may also be affected and side effects depend on the treated area.

\medterm{Gamna-Gandy Bodies} Small, firm, yellow-brown or reddish-brown siderofibrotic nodules in the splenic red pulp containing deposits of hemosiderin, iron, and calcium encrusting connective tissue and elastic fibers, resulting from organized focal hemorrhages in chronic congestive splenomegaly and sickle cell disease.
\synonyms
Gamna-Gandy Nodules; Siderofibrotic Bodies

\medterm{Ganciclovir} An antiviral medicine that interferes with viral DNA production and is used mainly for serious cytomegalovirus infection. It can suppress bone marrow, affect the kidneys, and cause other serious toxicity.

\medterm{Ganglia} Groups of nerve-cell bodies outside the brain and spinal cord. They relay and organize signals between the central nervous system and the rest of the body.

\medterm{Ganglioglioma} A usually slow-growing tumor containing nerve cells and supporting brain cells. It often causes seizures and may occur in the brain or spinal cord; treatment commonly involves surgical removal.

\medterm{Ganglion} A cluster of nerve-cell bodies outside the brain and spinal cord. It acts as a relay point for signals between the central nervous system and body tissues.

\medterm{Ganglion Cell} A nerve cell located in a ganglion. Retinal ganglion cells collect visual signals and send them through the optic nerve to the brain.

\medterm{Ganglion Cyst} A noncancerous, fluid-filled lump near a joint or tendon, most often at the wrist or hand. It may be painless or cause pressure, pain, numbness, or limited movement.

\synonyms
Synovial Cyst; Bible Cyst

\medterm{Ganglionectomy} Surgical removal of a ganglion, a cluster of nerve-cell bodies, or occasionally a ganglion cyst. The purpose, benefits, and risks depend on the exact structure and body site being treated.

\medterm{Ganglioneuroblastoma} A tumor made of both mature and immature nerve-related cells, usually in the adrenal or tissue beside the spine. It mainly affects children and may cause a mass, pain, high blood pressure, or other symptoms.

\medterm{Ganglioneuroma} A usually noncancerous tumor arising from nerve tissue of the autonomic nervous system. It often causes no symptoms but may produce a mass, pain, or hormone-related effects depending on its location.

\medterm{Ganglioneuromatosis} A condition in which nerve-related ganglion tissue grows diffusely or in multiple areas. Effects depend on the organs involved and may include masses, pain, bowel symptoms, or hormone-related changes.

\medterm{Ganglionic Blockers} Medicines that interrupt nerve signals at autonomic ganglia, affecting blood pressure, heart rate, digestion, and other automatic functions. They are rarely used because of widespread side effects.

\medterm{Ganglioside} A fat-containing molecule in cell membranes, especially nerve cells. Gangliosides help cells communicate and recognize one another; abnormal accumulation occurs in several inherited storage disorders.

\medterm{Gangliosidosis} A group of inherited disorders in which gangliosides accumulate inside cells because a breakdown enzyme is missing. The buildup often damages the nervous system and may also affect the eyes, bones, organs, and development.

\medterm{Gangrene} Death of body tissue caused by severely reduced blood flow, infection, injury, or a combination of these factors. It may affect skin, muscle, or an internal organ. The area may become discolored, cold, swollen, painful, numb, or foul-smelling; fever, confusion, or feeling very ill may occur with spreading infection.

\medterm{Ganser Syndrome} A rare dissociative or psychiatric syndrome characterized by approximate answers, confusion, altered awareness, and sometimes hallucinations or neurologic-like symptoms. It may occur with severe stress, trauma, other mental disorders, or medical and substance-related conditions.

\medterm{Gap Junction} A tiny connection between neighboring cells that allows ions and small molecules to pass directly between them. Gap junctions help coordinate heartbeats, nerve activity, and tissue growth.

\medterm{Gardasil} A brand of vaccine that protects against infection with selected human papillomavirus (HPV) types. It helps prevent HPV-related cancers and genital warts and does not treat an existing HPV infection.

\medterm{Gardnerella} A group of bacteria that may live in the genital tract. Overgrowth of Gardnerella, especially G. vaginalis, can contribute to bacterial vaginosis with discharge, odor, and altered vaginal bacteria.

\medterm{Gardnerella Vaginalis} A bacterium associated with bacterial vaginosis, in which normal vaginal bacteria become imbalanced. It may be present without illness; when symptomatic, discharge, odor, irritation, or burning can occur.

\medterm{Gardner Syndrome} An inherited condition caused by an APC gene change, producing many colon polyps and increasing colorectal-cancer risk. It may also cause bone growths, dental abnormalities, skin cysts, and desmoid tumors.

\medterm{Gargle} To rinse the mouth and throat with liquid while making a bubbling sound. Gargling may soothe irritation or help deliver a prescribed treatment.

\medterm{Garré's Osteomyelitis} A chronic inflammatory reaction in the jawbone, often near an infected tooth, causing firm swelling. Dental infection, persistent swelling, or pain may be present.

\medterm{Gartner Duct} A remnant of an embryonic duct in the female reproductive tract. It can persist harmlessly or form a cyst along the side wall of the vagina or broad ligament.

\medterm{Gartner Duct Cyst} A usually harmless fluid-filled cyst arising from an embryonic duct remnant along the side wall of the vagina. It often causes no symptoms, but a large cyst may cause pressure, pain, or difficulty with intercourse.

\medterm{Gas} A state of matter without a fixed shape or volume. In medicine, gas may mean air or another gas in the digestive tract, blood, tissues, or a body cavity.

\medterm{Gas And Gas Pains} Fullness, bloating, cramps, or discomfort from swallowed air or gas made by intestinal bacteria. Foods, constipation, food intolerance, and digestive disorders can contribute.

\medterm{Gas Chromatography} A laboratory method that separates vaporized chemicals so they can be identified or measured. It is often combined with mass spectrometry to test medicines, poisons, pollutants, or biological samples.

\medterm{Gas Exchange} Movement of oxygen from inhaled air into blood and carbon dioxide from blood into exhaled air. Oxygen then moves to tissues while carbon dioxide returns to the lungs for removal.

\medterm{Gas –  Flatulence} Passage of gas from the digestive tract through the anus, caused by swallowed air and bacterial digestion of food. It is normal, but excess gas with bloating or pain may reflect diet, constipation, or digestive disease.

\synonyms
Flatulence; Flatus Expulsion

\medterm{Gas Gangrene} A life-threatening infection in which bacteria produce toxins and gas inside damaged tissue, rapidly destroying muscle. Severe pain, swelling, blisters, discoloration, or a crackling feeling may occur.

\medterm{Gaskin Maneuver} A maneuver for shoulder dystocia in which the birthing person moves onto hands and knees. The position can change pelvic dimensions and help release the baby’s shoulder.

\synonyms
All-Fours Position; Hands-and-Knees Posture

\medterm{Gasoline Poisoning} Harm from swallowing, inhaling, or getting gasoline on the skin. It can cause nausea, dizziness, drowsiness, seizures, abnormal heart rhythms, or chemical lung injury, especially if vomited and inhaled.

\synonyms
Petrol Poisoning; Hydrocarbon Toxicity

\medterm{Gas Poisoning} Illness caused by inhaling or otherwise being exposed to a harmful gas. Effects range from irritation and headache to oxygen deprivation, lung injury, confusion, or collapse.

\medterm{Gas Scavenger} A substance or device that removes unwanted gases from a fluid, container, or breathing system. In healthcare it may help prevent toxic exposure or maintain safe gas delivery.

\medterm{Gasserian Ganglion} A cluster of nerve cells near the base of the skull that carries sensation from the face and controls some chewing muscles. It is involved in trigeminal neuralgia and other facial-nerve disorders.

\medterm{Gastralgia} Pain in the upper abdomen or stomach area. Causes include indigestion, reflux, ulcers, gallbladder disease, pancreatic inflammation, infection, or other conditions; severe pain may occur with bleeding, vomiting, fever, or weight loss.

\medterm{Gas Transport} Movement of oxygen from the lungs to body tissues and carbon dioxide back to the lungs through the blood. Hemoglobin carries most oxygen and part of the carbon dioxide.

\medterm{Gastrectomy} Surgery removing part or all of the stomach, usually for cancer or severe disease. Food passage is reconstructed, and possible effects include early fullness, weight loss, nutritional deficiencies, reflux, or dumping symptoms.

\medterm{Gastric Acid} Strong acid made by the stomach lining that helps break down food, activate digestive enzymes, and kill many swallowed microbes. Excess can contribute to reflux or ulcers; too little can impair digestion and nutrient absorption.

\medterm{Gastric Adenocarcinoma} The most common stomach cancer, arising from cells lining the stomach. It may cause indigestion, pain, early fullness, vomiting, bleeding, or weight loss; tissue examination and staging identify its type and extent.

\medterm{Gastric Analysis} Examination of stomach contents for volume, acidity, or selected substances. It is used only in specific situations because endoscopy, imaging, and blood tests are more useful for many stomach disorders.

\medterm{Gastric Antral Vascular Ectasia} Enlarged fragile blood vessels in the lower stomach that can cause slow internal bleeding and iron-deficiency anemia. It may appear as red stripes or spots during endoscopy and is associated with some liver and autoimmune diseases.

\synonyms
GAVE; Watermelon Stomach

\medterm{Gastric Antrum} The distal, non-acid-secreting muscular portion of the stomach located between the incisura angularis and the pylorus, containing G cells that secrete gastrin and D cells that secrete somatostatin, serving as the site of mechanical food grinding and triturition.

\medterm{Gastric Balloon} A temporary balloon placed inside the stomach to create fullness and support weight loss. Nausea, vomiting, reflux, deflation, blockage, bleeding, or perforation can occur.

\medterm{Gastric Bezoar} A mass of material that cannot be digested and collects in the stomach. It may contain plant fiber, hair, or medicines and can cause fullness, pain, vomiting, blockage, or poor nutrition.

\medterm{Gastric Biopsy} Removal of a small sample of stomach lining, usually during endoscopy, for laboratory examination. It can help assess inflammation, infection, precancerous changes, or cancer.

\medterm{Gastric Bypass} An operation that makes the stomach smaller and connects it to a lower part of the small intestine. It limits food intake and nutrient absorption and can improve severe obesity and related diseases.

\medterm{Gastric Bypass Surgery} Surgery that creates a small stomach pouch and reroutes food to the small intestine. It produces substantial weight loss and may improve diabetes, high blood pressure, sleep apnea, and other obesity-related conditions.

\synonyms
Roux-en-Y Gastric Bypass; Bariatric Bypass

\medterm{Gastric Cancer} A cancer that begins in the stomach lining or other stomach tissue. It may cause indigestion, pain, early fullness, vomiting, weight loss, or bleeding, although early disease may cause few symptoms.

\medterm{Gastric Body} The main middle part of the stomach between the fundus and the antrum. It stores and mixes food and contains glands that release acid and digestive substances.

\medterm{Gastric Carcinoid Tumor} A neuroendocrine tumor arising in the stomach, often from cells that make digestive hormones. Some are linked to chronic gastritis and high gastrin levels; behavior ranges from slow-growing to aggressive.

\medterm{Gastric Cardia} The anatomical transition zone of the stomach immediately surrounding the gastroesophageal junction where the tubular esophagus empties into the stomach, lined by simple columnar mucous-secreting glands without parietal or chief cells.

\medterm{Gastric Culture} A laboratory test that grows bacteria from a stomach sample, usually collected during endoscopy. It can identify Helicobacter pylori and help select antibiotics, although other tests are more commonly used.

\synonyms
Helicobacter Pylori Culture; Gastric Biopsy Culture

\medterm{Gastric Dysplasia} Abnormal cells in the stomach lining that have some features of cancer but have not invaded deeper tissue. Low- or high-grade changes may increase the risk of stomach cancer.

\medterm{Gastric Emptying} The passage of food and liquid from the stomach into the small intestine through coordinated stomach contractions and opening of the pylorus. Abnormally slow or rapid emptying can cause digestive symptoms.

\medterm{Gastric Emptying Study} A test that measures how quickly food leaves the stomach. The patient eats a meal containing a small tracer, and a scanner records the remaining food over several hours to detect unusually slow or rapid emptying.

\medterm{Gastric Fistula} An abnormal passage between the stomach and another organ or the skin. Stomach contents may leak through it, causing drainage, infection, skin irritation, dehydration, or poor nutrition.

\medterm{Gastric Fundus} The upper, rounded part of the stomach, located above the main stomach body and near the diaphragm. It can expand to store food and contains glands that produce digestive secretions.

\medterm{Gastric Glands} Glands in the stomach lining that make acid, digestive enzymes, mucus, intrinsic factor, and hormones.

\medterm{Gastric Hemorrhage} Bleeding from the stomach, commonly caused by an ulcer, severe inflammation, a tumor, or blood-vessel damage. It may cause vomiting blood, black stools, dizziness, anemia, or shock.

\medterm{Gastric Juice} Fluid made by the stomach containing acid, digestive enzymes, mucus, water, salts, and intrinsic factor. It helps break down food, protect the stomach lining, and absorb vitamin B12 later in digestion.

\medterm{Gastric Lavage} A procedure that flushes the stomach through a tube passed through the mouth or nose. It is rarely used because it can cause aspiration or injury and is reserved for selected emergencies.

\medterm{Gastric Leiomyosarcoma} A rare cancer of smooth muscle in the stomach wall. It may cause pain, bleeding, vomiting, or a mass and is evaluated with imaging and biopsy; treatment usually involves surgery when possible.

\medterm{Gastric Lymphoma} A cancer of immune-system cells that starts in the stomach. Some types are linked to Helicobacter pylori infection and may improve after eradication treatment; others are treated with chemotherapy, immunotherapy, or surgery.

\medterm{Gastric Malignancy} Cancer that begins in the stomach, most often in the stomach lining. It may cause indigestion, early fullness, pain, vomiting, weight loss, anemia, or bleeding, though early disease may have no symptoms; endoscopy with biopsy confirms the type and scans assess spread.

\medterm{Gastric Motility} The coordinated muscle activity that mixes stomach contents and moves them toward the small intestine.

\medterm{Gastric Mucosa} The moist inner lining of the stomach. It contains glands that produce acid, digestive enzymes, mucus, and intrinsic factor and is protected by a mucus barrier; inflammation or injury can cause pain or bleeding.

\medterm{Gastric Muscularis} The muscle layers in the stomach wall that mix food and move it toward the pylorus. The stomach has an extra inner muscle layer that helps with grinding.

\medterm{Gastric Necrosis} Death of stomach tissue caused by severely reduced blood flow, corrosive injury, extreme stretching, infection, or another serious insult. It can lead to perforation, severe infection, and shock.

\medterm{Gastric Obstruction} Partial or complete blockage that prevents food and liquid from leaving the stomach normally. It may cause repeated vomiting, upper-abdominal pain, fullness, swelling, dehydration, and weight loss.

\medterm{Gastric Outlet Obstruction} Blockage where the stomach empties into the small intestine, usually at the pylorus or duodenum. Causes include ulcer scarring, tumors, inflammation, and infantile pyloric stenosis; vomiting can cause dehydration.

\medterm{Gastric Pneumatosis} Gas within the stomach wall. It may result from pressure-related injury and be relatively mild, or from severe infection, which can destroy tissue and cause sepsis.

\medterm{Gastric Polyps} Growths projecting from the stomach lining. Most are harmless, but some types are linked to chronic inflammation, inherited conditions, or long-term acid-suppressing treatment and may contain precancerous changes.

\medterm{Gastric Perforation} A hole through the stomach wall that allows stomach contents to leak into the abdomen. It can cause severe inflammation, infection, and shock.

\medterm{Gastric Rugae} Folds in the inner lining of the empty stomach that flatten as the stomach fills. They allow the stomach to expand and increase the lining's surface area.

\medterm{Gastric Secretion} The release of acid, enzymes, mucus, intrinsic factor, and hormones by the stomach lining.

\medterm{Gastric Serosa} The thin outer covering of the stomach, formed by the visceral peritoneum. It reduces friction as the stomach moves inside the abdomen.

\medterm{Gastric Stenosis} Abnormal narrowing of the stomach or its outlet that slows or blocks emptying. It may result from birth defects, ulcer scarring, inflammation, tumors, or thickening of the outlet muscle.

\medterm{Gastric Submucosa} The connective-tissue layer beneath the stomach lining. It contains blood vessels, lymph vessels, nerves, and supporting tissue.

\medterm{Gastric Suction} Removal of stomach contents through a tube passed through the nose or mouth. It may relieve severe distension or vomiting from obstruction and is used only when clinically indicated because complications can occur.

\synonyms
Nasogastric Aspiration; Stomach Pumping

\medterm{Gastric Surgery} An operation on the stomach to remove disease, repair damage, reduce its size, or bypass a blockage. The exact procedure and risks depend on the condition and may include bleeding, infection, leakage, or altered digestion.

\medterm{Gastric Tissue Biopsy And Culture} Removal of small stomach-lining samples during endoscopy for microscopic examination and, when needed, laboratory testing for infection. It can investigate inflammation, Helicobacter pylori, precancerous changes, or tumors.

\synonyms
Endoscopic Gastric Biopsy; Mucosal Culture

\medterm{Gastric Ulcer} An open sore in the stomach lining, commonly caused by Helicobacter pylori infection or anti-inflammatory medicines. It may cause burning upper-abdominal pain, bleeding, anemia, or perforation.

\medterm{Gastric Varices} Enlarged, fragile veins in the stomach caused by increased pressure in the portal circulation, often from advanced liver disease. Rupture can cause sudden, life-threatening vomiting of blood.

\medterm{Gastric Volvulus} Twisting of the stomach that can block food passage and cut off blood supply. It may cause sudden severe upper-abdominal pain, retching without vomiting, bloating, and inability to pass a tube.

\medterm{Gastric Wall} The layered wall of the stomach, made up of the mucosa, submucosa, muscle layers, and outer covering. It allows the stomach to store, mix, and move food.

\medterm{Gastrin} A hormone made mainly in the lower stomach and upper small intestine. It stimulates stomach acid production and helps maintain the stomach lining; levels rise after eating and may become excessive in certain diseases.

\medterm{Gastrin Blood Test} A blood test measuring the hormone gastrin. It helps investigate excessive stomach acid, recurrent ulcers, or loss of acid-producing cells. Acid-suppressing medicines and fasting can significantly affect the result.

\synonyms
Serum Gastrin Level; Fasting Gastrin Assay

\medterm{Gastrinoma} A tumor that makes too much gastrin, a hormone that causes the stomach to produce acid. It usually starts in the pancreas or duodenum and may cause Zollinger--Ellison syndrome, repeated peptic ulcers, diarrhea, and acid reflux. Some cases are linked to MEN1 and may spread to the liver or nearby lymph nodes.

\medterm{Gastrin-Releasing Peptide} A regulatory neuropeptide released by postganglionic parasympathetic vagal nerve endings in the gastric antrum that binds to receptors on G cells, stimulating gastrin release and subsequent gastric acid secretion.

\synonyms
GRP; Mammalian Bombesin

\medterm{Gastrin Secretion} Release of the hormone gastrin, mainly from cells in the lower stomach after food enters. Gastrin increases stomach acid and digestive-enzyme release and supports growth and function of the stomach lining.

\medterm{Gastritis} Inflammation or irritation of the stomach lining. Common causes include \textit{H. pylori} infection, anti-inflammatory medicines, alcohol, autoimmune disease, and severe illness. It may cause upper-abdominal pain, nausea, vomiting, or bleeding.

\synonyms
Type A Gastritis; Autoimmune Metaplastic Atrophic Gastritis

\medterm{Gastrocnemius} The large muscle at the back of the calf. It helps point the foot downward and bend the knee during standing, walking, running, and jumping; injury can cause calf pain, swelling, or weakness.

\medterm{Gastrocnemius Muscle} The large superficial calf muscle that joins the Achilles tendon. It points the foot downward and helps bend the knee and push the body forward during walking, running, and jumping.

\medterm{Gastroduodenal Artery} A major arterial branch arising from the common hepatic artery, descending posterior to the first part of the duodenum before dividing into the right gastro-omental and superior pancreaticoduodenal arteries; erosion by a posterior duodenal ulcer produces massive upper gastrointestinal hemorrhage.

\synonyms
GDA

\medterm{Gastroduodenal Crohn Disease} Crohn disease affecting the stomach or first part of the small intestine. It may cause upper-abdominal pain, nausea, vomiting, weight loss, ulcers, narrowing, or blockage; treatment varies with symptoms and disease severity.

\medterm{Gastroduodenostomy} An operation creating a new opening between the stomach and duodenum so food can pass directly into the small intestine. It may bypass a blockage or restore passage after part of the stomach is removed.

\medterm{Gastroenteritis} Infection or inflammation of the stomach and intestines, commonly caused by viruses and sometimes by bacteria, parasites, or toxins. It usually causes diarrhea, vomiting, cramps, and fever and can lead to dehydration.

\medterm{Gastroenteritis, Acute} Alternate name; see \textit{Stomach Flu}.

\medterm{Gastroenteritis Bacterial} Inflammation of the stomach or intestines caused by bacteria, usually after contaminated food or water. It can cause diarrhea, cramps, fever, vomiting, and dehydration; blood in stool or severe illness may occur.

\medterm{Gastroenteritis Viral} Alternate name; see \textit{Viral gastroenteritis}.

\medterm{Gastroenterologist} A physician specializing in disorders of the digestive tract, liver, gallbladder, bile ducts, and pancreas. They diagnose and treat digestive problems using medicines, endoscopy, procedures, lifestyle care, and surgery referrals.

\medterm{Gastroenterology} The medical specialty concerned with the digestive tract, liver, gallbladder, bile ducts, and pancreas. Gastroenterologists diagnose and treat these conditions with medicines, endoscopy, procedures, and referrals for surgery.

\medterm{Gastroenterostomy} Surgical creation of an opening between the stomach and a segment of the small intestine, usually the jejunum, to bypass or replace a blocked or diseased gastric outlet.

\medterm{Gastroepiploic Artery} The gastroepiploic arteries run along the greater curve of the stomach and supply blood to the stomach and nearby omentum.

\medterm{Gastroesophageal} Relating to both the stomach and esophagus, especially their joining point or the backward movement of stomach contents into the esophagus.

\medterm{Gastroesophageal Junction Adenocarcinoma} A gland-forming cancer arising where the esophagus joins the stomach. It may cause difficulty swallowing, reflux, weight loss, anemia, or bleeding; endoscopy, biopsy, and scans determine treatment and extent.

\medterm{Gastroesophageal Reflux} Backward flow of stomach contents into the esophagus. It may cause heartburn, sour fluid in the throat, cough, chest discomfort, or no symptoms; frequent reflux can injure the esophageal lining.

\medterm{Gastroesophageal Reflux Disease} A long-term condition in which stomach contents repeatedly flow back into the esophagus, causing symptoms such as heartburn, regurgitation, chest pain, cough, or difficulty swallowing and sometimes damaging the esophageal lining.

\synonyms
GERD; Acid Reflux Disease

\medterm{Gastrogastric Intussusception} A rare condition in which one part of the stomach folds into an adjoining part, sometimes because of a polyp or tumor. It may cause pain, vomiting, blockage, or bleeding.

\medterm{Gastrografin Study} An X-ray study using a water-soluble contrast liquid to show the digestive tract. It may be chosen when a leak or perforation is possible, because the contrast is safer outside the tract than barium.

\medterm{Gastrointestinal} Relating to the stomach and intestines, or to the complete digestive tract from the mouth to the anus. These organs digest food, absorb nutrients and water, move contents forward, and remove waste.

\medterm{Gastrointestinal Agents} Medicines or other substances that act on the digestive tract to relieve or treat problems such as nausea, vomiting, diarrhea, constipation, ulcers, reflux, or intestinal spasms. Their effects and risks depend on the specific agent.

\medterm{Gastrointestinal Amyloidosis} Buildup of abnormal protein deposits in the digestive tract. It can impair movement, absorption, or bleeding control and may occur with disease affecting other organs such as the heart, kidneys, or nerves.

\medterm{Gastrointestinal Bleeding} Bleeding anywhere in the digestive tract. It may be acute or chronic and may come from the upper digestive tract, small intestine, or lower digestive tract. It can cause vomiting blood, black stools, or red blood in the stool.

\medterm{Gastrointestinal Bleeding Scan} A nuclear-medicine scan that uses a small radioactive tracer in the blood to find active bleeding in the digestive tract. It can help localize bleeding when endoscopy does not identify the source.

\medterm{Gastrointestinal Cancer} Cancer arising in the digestive tract or related organs, including the esophagus, stomach, intestine, rectum, liver, pancreas, and biliary tract. Symptoms and treatment depend on the organ, tumor type, and stage.

\medterm{Gastrointestinal Decontamination} Methods used after some poisonings to reduce absorption, such as activated charcoal or whole-bowel irrigation. The choice depends on the substance, timing, symptoms, and risk of aspiration; stomach washing is rarely used.

\medterm{Gastrointestinal Disease} Any disorder affecting the digestive tract or related organs. Examples include reflux, ulcers, inflammatory bowel disease, infection, malabsorption, liver disease, and cancer.

\medterm{Gastrointestinal Disorder} A disease affecting any part of the digestive tract, from the esophagus to the anus. It may cause pain, nausea, vomiting, diarrhea, constipation, bleeding, blockage, or poor nutrient absorption, depending on the affected structure.

\medterm{Gastrointestinal Fistula} An abnormal connection between the digestive tract and another organ, the skin, or another section of bowel. Digestive fluid may leak through it, causing infection, dehydration, skin damage, or poor nutrition.

\medterm{Gastrointestinal Gas} Air or other gas within the stomach or intestines, usually from swallowed air or bacterial digestion of food. Excess gas may cause belching, bloating, cramps, or passing gas and can increase with constipation or food intolerance.

\medterm{Gastrointestinal Health} The normal structure and function of the digestive system, including digestion, absorption, bowel movement, and a balanced relationship with gut microorganisms. Problems may include bleeding, persistent pain, vomiting, or unexplained weight loss.

\medterm{Gastrointestinal Hemorrhage} Bleeding from any part of the digestive tract, from the esophagus to the anus. It may cause vomiting blood, black stools, red blood in stool, anemia, dizziness, or shock.

\medterm{Gastrointestinal Hormone} A chemical messenger made by cells in the digestive tract that helps regulate stomach acid, intestinal movement, digestion, appetite, gallbladder activity, or blood-glucose control.

\medterm{Gastrointestinal Motility Disorder} Abnormal movement of food, liquid, or stool through the digestive tract. It can cause difficulty swallowing, nausea, vomiting, bloating, abdominal pain, constipation, or diarrhea; causes may involve digestive muscles, nerves, hormones, or a blockage.

\medterm{Gastrointestinal Neoplasm} An abnormal growth arising in the digestive tract. It may be a harmless growth, a precancerous change, or cancer; symptoms and treatment depend on its tissue, location, and behavior.

\medterm{Gastrointestinal Obstruction} Partial or complete blockage of food, fluid, or stool through the digestive tract. It can cause cramping pain, vomiting, swelling, constipation, dehydration, loss of blood supply, perforation, or sepsis.

\medterm{Gastrointestinal Pain} Pain arising from the digestive tract or related organs. Causes include infection, ulcers, inflammation, blockage, gallstones, and diseases of the liver or pancreas. It may be severe, persistent, or worsening.

\medterm{Gastrointestinal Perforation} A hole through the wall of the digestive tract that allows its contents to leak into surrounding tissue. It can cause severe abdominal pain, peritonitis, sepsis, and shock.

\medterm{Gastrointestinal Stromal Tumor} A tumor that begins in supporting cells in the wall of the digestive tract, most often in the stomach or small intestine. Its behavior varies; treatment may include surgery or medicines designed to block specific tumor signals.

\medterm{Gastrointestinal System} The connected organs that move and digest food, absorb nutrients and water, and remove solid waste. It includes the mouth, esophagus, stomach, intestines, rectum, and anus, with help from the liver, gallbladder, and pancreas.

\medterm{Gastrointestinal Tract} The continuous digestive passage from the mouth through the esophagus, stomach, small and large intestines, rectum, and anus. It moves food, digests it, absorbs nutrients and water, and removes waste.

\medterm{Gastrointestinal Villi} Tiny finger-like projections lining the small intestine that greatly increase the area available to absorb nutrients and water. Each contains small blood vessels and a lymph vessel that carry absorbed substances away.

\medterm{Gastrojejunostomy} An operation creating a new connection between the stomach and jejunum, the middle part of the small intestine. It may bypass a blocked stomach outlet or replace a damaged section of digestive tract.

\medterm{Gastromalacia} Abnormal softening of the stomach wall. It is a rare descriptive finding that may result from severe inflammation, poor blood supply, or tissue injury and can increase the risk of weakness or perforation.

\medterm{Gastroparesis} A disorder in which the stomach empties too slowly, or sometimes stops emptying, even though there is no physical blockage. The stomach muscles or the nerves that control them do not move food normally. It may cause early fullness, nausea, vomiting, bloating, belching, upper-abdominal pain, heartburn, and poor appetite. Diabetes is the most common known cause; other causes include stomach surgery, nerve or muscle disorders, some autoimmune or viral illnesses, although the cause is often unknown. Some medicines can slow stomach emptying or cause similar symptoms without causing gastroparesis itself. Severe disease can lead to dehydration, poor nutrition, and difficult-to-control blood glucose.

\synonyms
Delayed Gastric Emptying; Gastric Hypomotility

\medterm{Gastropathy} Damage or disease of the stomach lining without the prominent inflammation seen in gastritis. Causes include medicines, alcohol, poor blood flow, bile reflux, and severe illness; symptoms may include pain, nausea, or bleeding.

\medterm{Gastroplasty} Surgery that changes the size or shape of the stomach, usually to support weight loss or repair a defect. Possible complications include bleeding, infection, leakage, reflux, and nutritional deficiencies.

\medterm{Gastroschisis} A birth defect in which the intestine, and sometimes other organs, protrudes through an opening beside the umbilicus without a protective covering. The exposed bowel may be swollen or damaged.

\medterm{Gastroschisis Repair} Surgical repair of gastroschisis, a congenital abdominal wall defect where intestines protrude through an opening near the umbilicus. The bowel is returned to the abdomen and the defect is closed, sometimes in stages using a temporary covering.

\synonyms
Abdominal Wall Defect Repair; Congenital Evisceration Closure

\medterm{Gastroscope} A thin, flexible instrument with a light and camera used to view the esophagus, stomach, and first part of the small intestine. It can also pass tools for biopsy or treatment.

\medterm{Gastroscopy} Examination of the esophagus, stomach, and duodenum using a flexible camera tube passed through the mouth. It can identify bleeding, ulcers, inflammation, narrowing, or tumors and allows biopsies or selected treatments.

\medterm{Gastrospenic Fistulization} An abnormal connection between the stomach and spleen, usually caused by a severe stomach ulcer or a splenic abscess. It can cause upper-abdominal pain, digestive-tract bleeding, severe infection, and sepsis.

\medterm{Gastrostomy} A surgically created opening through the abdominal wall into the stomach. A tube placed through it can provide food, fluids, or medicines or remove stomach contents when normal feeding is unsafe or impossible.

\medterm{Gastrostomy Tube} A tube passed through the abdominal wall into the stomach to provide food, fluids, or medicines when swallowing is unsafe or inadequate. Problems can include skin infection, leakage, blockage, bleeding, and movement of the tube.

\medterm{Gastrostomy Tube Insertion} A surgical, laparoscopic, or endoscopic procedure in which a feeding tube is placed directly through the abdominal wall into the stomach for long-term enteral nutrition.

\synonyms
G-Tube Placement; Percutaneous Endoscopic Gastrostomy; PEG Tube Insertion

\medterm{Gastrula} An early embryo stage formed after gastrulation, when cell layers reorganize into ectoderm, mesoderm, and endoderm that give rise to body tissues.

\medterm{Gastrulation} An early stage of embryo development in which cells move and reorganize into three layers that later form all major tissues and organs.

\medterm{Gated Blood Pool Scan} A radionuclide cardiac study, also called radionuclide ventriculography or MUGA, that uses ECG-gated images of labeled blood to measure ventricular ejection fraction and assess ventricular wall motion. It is used selectively when quantitative ventricular function is needed.

\medterm{Gatifloxacin} Gatifloxacin is a fluoroquinolone antibiotic that inhibits bacterial DNA gyrase and topoisomerase IV; systemic use is restricted in many settings because of dysglycaemia and other class toxicities, while ophthalmic use remains available in some regions.

\medterm{Gaucher Disease} An autosomal recessive lysosomal storage disorder caused by reduced activity of the enzyme glucocerebrosidase, usually because of disease-causing variants in the \textit{GBA1} gene. Glucosylceramide accumulates in macrophages and other cells, affecting the spleen, liver, bone marrow, bones, lungs, and sometimes the nervous system. It can cause enlargement of the spleen or liver, anemia, low platelets, easy bruising, bone pain or fractures, and neurologic problems. Type 1 is usually non-neuronopathic, while types 2 and 3 involve the nervous system.

\medterm{Gaucher's Disease} Alternate spelling; see \textit{Gaucher Disease}.

\medterm{Gaucher Splenomegaly} Alternate name; see \textit{Gaucher disease}.

\medterm{Gaultheria Oil} Oil of wintergreen, an essential oil containing concentrated methyl salicylate. It may be used in small amounts in topical products, but swallowing even a small amount can cause serious salicylate poisoning, especially in children.

\medterm{Gauss} A unit of magnetic flux density equal to $10^{-4}$ tesla, used in physics, imaging, and descriptions of magnetic fields.

\medterm{Gauze} A woven or nonwoven material used to clean, dress, pack, or absorb fluid from a wound. It may be sterile or nonsterile and comes in several forms.

\medterm{Gavage} Giving liquid food or medicine through a tube into the stomach. It is used when swallowing is unsafe or oral intake is inadequate, including for premature infants who cannot coordinate sucking, swallowing, and breathing.

\medterm{Gaze Palsy} An inability to move both eyes conjugately in a particular direction (horizontal, vertical, or upgaze), caused by lesions involving the cortical frontal eye fields, the paramedian pontine reticular formation (PPRF), or the midbrain rostral interstitial nucleus of the medial longitudinal fasciculus (riMLF).
\synonyms
Conjugate Gaze Palsy; Ocular Conjugate Palsy

\medterm{GBM} An abbreviation for glioblastoma, an aggressive cancer that starts in supporting cells of the brain. It can grow quickly and cause headache, seizures, weakness, speech or memory problems, or personality changes. Treatment may include surgery, radiation, and medicines.

\synonyms
Glioblastoma Multiforme; Grade IV Astrocytoma

\medterm{Gefitinib} An oral epidermal growth factor receptor tyrosine-kinase inhibitor used for selected advanced non-small-cell lung cancers with activating EGFR mutations.

\medterm{Gegenhalten} A velocity-independent, variable increase in resistance to passive limb movement in which the patient involuntarily opposes the examiner's manipulation with a force proportional to that applied by the examiner, commonly observed in diffuse frontal lobe disease, dementia, and metabolic encephalopathy.
\synonyms
Paratonia; Oppositional Paratonia

\medterm{Geiger Counter} An instrument that detects and measures ionizing radiation. It can identify radioactive contamination or exposure, but its readings do not by themselves show how much biological harm the radiation caused.

\medterm{Gel} A semisolid preparation in which liquid is held in a soft network. Gels may be applied to the skin, eyes, mouth, or other areas to soothe tissue, deliver medicine, or form a protective layer.

\medterm{Gelastic Seizure} A rare seizure that causes sudden, repeated bursts of laughter without feeling amused. It is often linked to a small growth called a hypothalamic hamartoma, but other brain conditions can also cause it.

\medterm{Gelatin} Gelatin is a protein made from collagen and used in foods, medicines, capsules, and some medical materials.

\medterm{Gel Electrophoresis} Gel electrophoresis separates DNA, RNA, or proteins by moving them through a gel in an electric field, usually according to size and charge.

\medterm{Gelfoam} An absorbable sponge made from purified animal gelatin that helps control bleeding during surgery. It provides a surface for clot formation and is gradually absorbed; swelling or infection can rarely occur.

\medterm{Gelineau Syndrome} A sleep disorder causing repeated, uncontrollable daytime sleep episodes. Some people also have sudden muscle weakness triggered by emotion, sleep paralysis, or vivid dream-like experiences while falling asleep or waking.

\medterm{Gemcitabine} A chemotherapy medicine that becomes active inside cells and interferes with DNA production, slowing or killing rapidly dividing cancer cells. It is used for several cancers and may cause low blood counts, nausea, rash, or organ toxicity.

\medterm{Gemcitabine Hydrochloride} The hydrochloride form of gemcitabine, a chemotherapy medicine that interferes with DNA production in dividing cells. It is used for several cancers; important risks include infection from low blood counts and treatment-related organ injury.

\medterm{Gemella} Gemella is a group of bacteria that normally may inhabit the mouth and upper airway but can rarely cause serious infections, including bloodstream infection.

\medterm{Gemfibrozil} A medicine that lowers very high blood triglycerides and may reduce the risk of pancreatitis. It changes how the liver processes fats but can interact with statins and increase the risk of serious muscle injury.

\medterm{Gemifloxacin} An antibiotic in the fluoroquinolone group used for selected bacterial respiratory infections. It can cause nausea and, rarely, tendon injury, nerve symptoms, abnormal heart rhythm, severe diarrhea, or dangerous allergic reactions.

\medterm{Gemination} An attempt by one developing tooth bud to divide, producing a large tooth with a split or partly divided crown. The total number of teeth is usually normal, and dental assessment guides treatment.

\medterm{Gemistocytic Astrocytoma} A former histologic variant of diffuse astrocytoma composed partly of gemistocytic astrocytes. The term is now used mainly descriptively because integrated molecular classification determines the modern tumor diagnosis and grade.

\medterm{Gemtuzumab Ozogamicin} An antibody-linked chemotherapy medicine used for selected acute myeloid leukemia. The antibody attaches to CD33-positive leukemia cells and delivers a toxic drug; liver injury, infusion reactions, and low blood counts are important risks.

\medterm{Gender} A social, cultural, and personal concept involving identity, roles, expression, and expectations. Gender is related to, but distinct from, sex characteristics and sexual orientation.

\medterm{Gender Dysphoria} Clinically significant distress or impairment associated with incongruence between a person’s experienced or expressed gender and sex assigned at birth. Having a gender identity different from sex assigned at birth does not by itself constitute gender dysphoria.

\medterm{Gender Identity} A person’s internal sense of being a woman, man, both, neither, or another gender, which may or may not correspond to sex assigned at birth.

\medterm{Gender Identity Disorder} Alternate historical term; see \textit{Gender Dysphoria}.

\medterm{Gender Incongruence} A difference between a person’s experienced or expressed gender and sex assigned at birth. Gender incongruence is a descriptive term and does not by itself imply distress, impairment, or a mental disorder.

\medterm{Gene} A section of DNA that contains instructions for making a functional product, usually a protein or RNA. Genes are arranged at specific locations on chromosomes and influence inherited traits and cell function.

\medterm{Gene Amplification} An increase in the number of copies of a gene or genomic region, which can raise gene expression and drive some inherited conditions or cancers.

\medterm{Gene Conversion} A nonreciprocal transfer of sequence information from one DNA region to a homologous region, potentially changing an allele during recombination or repair.

\medterm{Gene Deletion} Loss of part or all of a gene from DNA, potentially abolishing its product or altering dosage and causing disease when clinically significant.

\medterm{Gene Dosage} The number of functional copies of a gene and the amount of product they produce, influenced by deletions, duplications, chromosome number, and regulation.

\medterm{Gene Duplication} An extra copy of a gene created by replication or rearrangement, which may have no effect, alter dosage, or provide material for new gene function.

\medterm{Gene-Environment Interaction} A gene-environment interaction occurs when an environmental exposure changes disease risk differently according to a person's genotype, so inherited susceptibility and exposure jointly influence phenotype.

\medterm{Gene Expression} The process by which a cell uses information in a gene to make RNA or protein. Cells control when and how strongly genes are expressed, allowing different tissues to develop and perform specialized functions.

\medterm{Gene Frequency} The proportion of all gene copies in a population that carry a particular version, or allele, of that gene. It describes population genetics and does not by itself predict an individual’s disease risk.

\medterm{Gene Fusion} A hybrid gene formed when segments of two separate genes join, sometimes producing an abnormal protein that drives cancer or other disease.

\medterm{Gene Modification} An intentional or naturally occurring alteration of DNA sequence or gene regulation that changes cellular function or inherited traits.

\medterm{Gene-Modified} Describing cells or organisms whose genetic material has been deliberately altered using recombinant DNA or genome-editing methods.

\medterm{Gene Mutation} A permanent change in the DNA sequence of a gene, which may be harmless, beneficial, or pathogenic depending on its effect and inheritance.

\medterm{Gene Panel} A genetic test that examines a selected group of genes linked to a person’s symptoms or suspected condition. It can provide focused results but may miss changes in genes or DNA regions not included in the panel.

\medterm{General Anesthesia} A controlled, temporary state of unconsciousness produced by medicines so a person does not feel or remember a procedure. Breathing, blood pressure, heart rate, and other body functions are monitored throughout.

\medterm{General Anesthetic} A medicine that produces a controlled, reversible state of unconsciousness and reduced awareness of pain, often with muscle relaxation, so surgery or another procedure can be performed safely.

\medterm{Generalised} Widespread or involving multiple regions of the body rather than being confined to a single localized site.

\medterm{Generalized} Distributed widely or involving most or all of a body region, system, or population rather than being confined to one focus.

\medterm{Generalized Anxiety Disorder} A mental-health disorder involving excessive, difficult-to-control worry about several areas of life on most days for at least six months. Restlessness, tiredness, poor concentration, irritability, muscle tension, and sleep problems may occur.

\medterm{Generalized Anxiety Disorder In Children} Excessive, difficult-to-control worry in children about multiple topics (school performance, health, family safety) occurring more days than not for at least 6 months. Associated with restlessness, fatigue, difficulty concentrating, irritability, muscle tension, and sleep disturbance. Treatment includes cognitive behavioral therapy and possibly medication.

\synonyms
Pediatric GAD; Childhood Generalized Anxiety

\medterm{Generalized Granuloma Annulare} A widespread skin condition causing many small, skin-colored, red, or purple bumps over the trunk and limbs. It may be associated with diabetes and can last longer than localized granuloma annulare.

\medterm{Generalized Lipodystrophy} A disorder with near-total loss of adipose tissue, severe insulin resistance, diabetes, high triglycerides, fatty liver, and abnormal fat deposition in other tissues.

\medterm{Generalized Lymphadenopathy} Pathological enlargement of lymph nodes involving two or more non-contiguous anatomical lymph node chains. Unlike localized adenopathy, it reflects systemic, hematologic, infectious, or autoimmune disease. Major etiologies include lymphomas, leukemias (such as chronic lymphocytic leukemia), systemic viral infections (HIV, Epstein-Barr virus, cytomegalovirus), secondary syphilis, disseminated tuberculosis, systemic lupus erythematosus, sarcoidosis, and drug-induced hypersensitivity reactions (such as DRESS syndrome).

\synonyms
Diffuse Lymphadenopathy; Systemic Lymphadenopathy

\medterm{Generalized-Onset Seizure} A seizure that begins in networks on both sides of the brain. It may cause staring, brief muscle jerks, stiffening, loss of muscle tone, or whole-body stiffening and jerking, depending on the type.

\medterm{Generalized Seizure} A seizure arising from or rapidly involving both cerebral hemispheres, commonly producing impaired awareness, bilateral motor activity, or brief absence episodes.

\medterm{Generalized Tonic-Clonic Seizure} A seizure causing loss of consciousness, stiffening of the whole body, and then repeated jerking of the limbs. Confusion and sleepiness often follow; a seizure lasting five minutes or longer is prolonged.

\medterm{Generalized Weakness} Reduced strength affecting multiple body regions, caused by systemic illness, metabolic or electrolyte disturbance, medication, infection, muscle disease, nerve disease, neuromuscular-junction disease, or brain or spinal-cord disease. It may occur with speech, facial, breathing, or walking difficulty.

\medterm{General Paresis} A late form of syphilis affecting the brain and spinal cord. It can gradually cause memory loss, personality or mood changes, poor judgment, weakness, tremor, or difficulty walking. Blood and spinal-fluid tests help identify the infection.

\synonyms
General Paresis of the Insane; Paralytic Dementia; Neurosyphilis

\medterm{Generator} A device that produces a short-lived radioactive substance from a longer-lived one for nuclear-medicine tests. The product is separated, checked for purity and sterility, and then used to make a tracer.

\medterm{Gene Rearrangement} Reorganization of DNA segments within or between chromosomes that can alter gene structure, expression, immune-receptor diversity, or cancer biology.

\medterm{Generic} Describing a medicine by its active ingredient rather than a brand name; an approved generic drug must meet applicable standards for quality and therapeutic equivalence to its reference product.

\medterm{Genes} Sections of DNA that contain instructions for making proteins or functional RNA. They are arranged on chromosomes and influence inherited traits, cell activity, disease risk, and responses to some treatments.

\medterm{Gene Silencing} Reduction or stopping of a gene's activity through DNA changes, chromatin control, RNA interference, or other processes that limit protein production.

\medterm{Genesis} The origin, formation, or development of something, such as a disease, organ, symptom, or process. Understanding its genesis can help explain causes, risk factors, progression, prevention, and treatment.

\medterm{Gene Targeting} A molecular method that changes a chosen genomic locus through homologous recombination or engineered editing to study or alter gene function.

\medterm{Gene Therapy} Treatment that adds, removes, replaces, or edits genetic material in a patient’s cells to treat disease. It may correct a missing gene function or alter harmful activity; risks include immune reactions and unintended genetic effects.

\medterm{Genetic Anticipation} Earlier onset or increasing severity of a hereditary disorder in successive generations, often caused by expansion of unstable repeated DNA sequences.

\medterm{Genetic Carrier} A person with one pathogenic variant for a recessive disorder who is usually unaffected but can transmit the variant to children.

\medterm{Genetic Code} The rules cells use to read genetic instructions in groups of three DNA or RNA building blocks. Each group directs the cell to add a particular amino acid to a protein or to stop making the protein.

\medterm{Genetic Counseling} A professional discussion that explains how inherited factors may affect a person or family. It may cover risk, testing choices, results, reproductive options, emotional support, and implications for relatives.

\medterm{Genetic Counselling} A process in which a trained professional explains inherited conditions, test results, health implications, and reproductive choices to a person or family.

\medterm{Genetic Counselor} A trained health professional who explains inherited conditions, estimates genetic risk, discusses testing, and helps patients and families understand results and choices. They may coordinate testing of relatives when appropriate.

\medterm{Genetic Disease} Alternate name; see \textit{Genetic Disorder}.

\medterm{Genetic Disorder} A disease or disorder caused wholly or partly by a pathogenic change in genetic material. The change may involve one or more genes, chromosome number or structure, mitochondrial DNA, or gene regulation. It may be inherited through germline DNA, arise de novo, or occur only in certain tissues as a somatic change. Some conditions also result from interactions among multiple genetic variants and environmental factors; not every genetic change causes disease.

\medterm{Genetic Drift} Random changes in allele frequencies from one generation to the next, especially in small populations, independent of natural selection.

\medterm{Genetic Engineering} Deliberate changing, adding, removing, or transferring genetic material using laboratory methods. It supports research, diagnosis, treatment development, agriculture, and biotechnology but requires careful assessment of safety, ethics, and unintended effects.

\medterm{Genetic Heterogeneity} The situation in which similar features or the same disease can result from different genetic causes. Different genes or different changes within one gene may produce overlapping effects.

\medterm{Genetic Linkage} The tendency of nearby genes or DNA markers on the same chromosome to be inherited together because crossing over is less likely between them.

\medterm{Genetic Marker} A detectable DNA variation or feature used to track inheritance, identify a chromosome region, assess disease risk, or distinguish samples.

\medterm{Genetic Polymorphism} A common DNA variation with at least two alleles in a population, which may influence traits, drug response, or disease susceptibility.

\medterm{Genetic Profile} A set of genetic variants or DNA markers describing an individual, tumor, organism, or sample for diagnosis, risk assessment, identification, or research.

\medterm{Genetic Recombination} Exchange or rearrangement of DNA segments that creates new allele combinations during meiosis, repair, immune-cell development, or laboratory manipulation.

\medterm{Genetics} Genetics is the study of genes, heredity, genetic variation, and how DNA and gene regulation influence traits, disease risk, diagnosis, and treatment response.

\medterm{Genetic Screening} Testing people without symptoms to estimate the chance of a genetic condition. Examples include newborn screening, carrier screening before or during pregnancy, and prenatal screening; screening results are not always diagnostic.

\medterm{Genetic Sequencing} Analysis of the order of DNA bases to identify variants associated with genetic disorders. The method is selected according to the gene, variant, and clinical question.

\synonyms
DNA Sequencing; Next-Generation Sequencing (NGS)

\medterm{Genetic Sonogram} A detailed pregnancy ultrasound, usually performed around 18–22 weeks, that examines fetal anatomy and may look for findings associated with chromosome conditions. It cannot detect every genetic disorder.

\medterm{Genetic Susceptibility} Inherited tendency to develop a disease or respond to an exposure, which usually modifies risk rather than determining outcome alone.

\medterm{Genetic Technique} A laboratory method used to detect, amplify, sequence, edit, compare, or otherwise analyze DNA, RNA, chromosomes, or gene expression.

\medterm{Genetic Testing} Analysis of DNA or chromosomes to assess disease risk, explain symptoms, identify carriers, or guide treatment. Results may also provide information about inherited traits and conditions affecting relatives.

\medterm{Genetic Transcription} The first step in using a gene: an enzyme copies information from DNA into RNA. The RNA may then guide protein production or perform another function inside the cell.

\medterm{Genetic Variant Of Uncertain Significance} A DNA change whose effect on health is not yet known. It does not by itself show that a person has a disease or predict future disease and may later be classified as harmless or disease-causing as evidence improves.

\medterm{Gene Transfer} Movement of genetic material into a cell or organism, naturally or through laboratory vectors, potentially changing gene expression or correcting a defect.

\medterm{Geniculate Body} One of two paired relay structures in the brain. The lateral geniculate body processes visual signals, while the medial geniculate body processes hearing signals before sending them to the cerebral cortex.

\medterm{Geniculate Ganglion} An L-shaped sensory ganglion of the facial nerve (cranial nerve VII) located within the facial canal of the temporal bone, housing the primary pseudo-unipolar sensory neuronal cell bodies conveying special taste sensation from the anterior two-thirds of the tongue (via chorda tympani).

\medterm{Genital} Relating to the reproductive organs or external sexual structures. Genital symptoms may involve skin, mucosa, nerves, glands, or the urinary and reproductive systems.

\medterm{Genital Edema} Swelling of genital tissues from fluid accumulation caused by infection, inflammation, allergy, trauma, venous or lymphatic obstruction, or systemic disease.

\medterm{Genital Herpes} An infection of the genital or anal area caused by herpes simplex virus type 1 or type 2 (HSV-1 or HSV-2). It spreads mainly through sexual skin-to-skin contact and can be transmitted even when no sores are present. It may cause painful blisters or ulcers, burning urination, swollen glands, or no symptoms; outbreaks can recur. Infection during pregnancy can threaten the newborn.

\medterm{Genital Herpes in Women} Herpes simplex virus infection of the female genitalia, causing painful blisters or ulcers. It may have no symptoms. Antiviral treatment can reduce symptoms and transmission; infection during pregnancy can affect the newborn.

\synonyms
Female Genital Herpes; HSV-2 Infection in Women

\medterm{Genitalia} The external and internal reproductive organs. The term includes the vulva and vagina,
or the penis, scrotum, and associated structures, together with the gonads and
reproductive ducts when used broadly; anatomy varies with sex, age, and developmental
stage.

\medterm{Genital Injury} Damage to the external or internal genital tissues from childbirth, accidents, sexual contact, assault, medical procedures, or disease. It may cause pain, swelling, bleeding, or difficulty urinating.

\medterm{Genital Mycoplasmas} Bacteria that may live in or infect the genital or urinary tract. Some cause inflammation of the urethra or cervix, pelvic inflammatory disease, or few symptoms; testing helps distinguish infection from harmless colonization.

\medterm{Genital Prolapse} The dropping or bulging of one or more pelvic organs (such as the bladder, uterus, or rectum) into or through the vagina due to weakening of the pelvic floor muscles and supporting tissues. It commonly causes a feeling of a vaginal bulge or heaviness, lower back discomfort, and urinary or bowel difficulties; risks include vaginal childbirth, aging, menopause, and chronic straining.
\synonyms Pelvic organ prolapse; Uterovaginal prolapse; Urogenital prolapse.

\medterm{Genitals} The external and internal reproductive organs. Female genitals: external (vulva — mons
pubis, labia majora, labia minora, clitoris, vaginal opening, Bartholin glands) and
internal (vagina, cervix, uterus, fallopian tubes, ovaries).

\medterm{Genital System} The organs involved in reproduction and sexual function, including the gonads, reproductive ducts, accessory glands, and external genital structures. Anatomy differs between individuals and changes with development and age.

\medterm{Genital Ulcer Disease} A clinical syndrome characterized by single or multiple ulcerative breaches of the genital or perianal skin and mucosa, predominantly caused by sexually transmitted pathogens including herpes simplex virus (HSV-2), \textit{Treponema pallidum} (syphilis chancre), and \textit{Haemophilus ducreyi} (chancroid).

\synonyms
GUD

\medterm{Genital Wart} A wart in the genital or anal area caused by certain types of human papillomavirus. Warts may be flat or raised and can recur; treatment removes visible lesions but may not eliminate the virus immediately.

\medterm{Genital Warts} Warts of the genital or anal region caused most often by low-risk types of human papillomavirus. They spread through direct skin-to-skin sexual contact and may be visible or asymptomatic.

\synonyms
Condylomata Acuminata; Anogenital Warts

\medterm{Genitofemoral Nerve} A branch of the lumbar plexus (arising from L1 and L2 nerve roots) that pierces the anterior surface of the psoas major muscle and divides into a genital branch (innervating the cremaster muscle and scrotal/labial skin; mediating the cremasteric reflex) and a femoral branch (cutaneous sensation to the anterior upper thigh).

\medterm{Genito-Pelvic Pain/Penetration Disorder} Persistent or recurrent difficulty with vaginal penetration, including pain, fear or anxiety about pain, involuntary pelvic-floor tightening, or pain during attempted penetration, causing clinically significant distress. Causes may include pelvic-floor problems, infection, hormonal changes, injury, or emotional factors.

\medterm{Genitourinary} Relating to the urinary system and the reproductive organs, which share developmental and anatomical relationships.

\medterm{Genitourinary Fistula} Alternate name for urogenital fistula, an abnormal connection between the urinary and genital tracts.

\medterm{Genito-Urinary Medicine} The medical specialty concerned with urinary and reproductive health. It includes evaluation and treatment of urinary infections, sexual health problems, sexually transmitted infections, and some reproductive conditions.

\medterm{Genitourinary Neoplasm} An abnormal growth in the urinary or reproductive organs, such as the kidney, bladder, prostate, testis, ovary, or uterus. It may be benign or cancerous; examination, imaging, and sometimes biopsy identify its nature.

\medterm{Genitourinary Syndrome Of Menopause} A group of changes in the genital and urinary tissues after menopause, caused mainly by lower estrogen levels. It may include vaginal dryness, burning, pain during sex, urinary urgency, painful urination, or repeated urinary infections.

\medterm{Genitourinary Tract} The combined urinary and reproductive organs and their associated ducts.

\medterm{Genodermatoses} Inherited disorders that mainly affect the skin, hair, nails, sweat glands, or related organs. Examples include inherited scaling disorders, blistering diseases, and some neurocutaneous syndromes; severity varies widely.

\medterm{Genome} The complete genetic material of an organism, including genes and DNA that does not directly code proteins. It contains instructions influencing development, body function, inherited traits, disease risk, and treatment response.

\medterm{Genomic Imprinting} A pattern in which a gene’s activity depends on whether its copy came from the mother or father. Abnormal imprinting can cause conditions such as Prader–Willi or Angelman syndrome.

\medterm{Genomic Instability} An increased tendency for DNA changes, chromosome damage, and mutations to accumulate, contributing to cancer development, treatment resistance, and disease progression.

\medterm{Genomics} The study of an organism’s complete genetic material, including its DNA sequence, variation, activity, and interaction with disease and treatment. It supports research, diagnosis, risk assessment, and personalized care.

\medterm{Genotoxicity} The ability of a substance or physical agent to damage genetic material. Genotoxicity may cause mutations or chromosome damage, but a positive test does not by itself establish that an agent causes cancer in humans.

\medterm{Genotype} A person’s genetic makeup, meaning the versions of one or more genes they carry. The genotype helps influence observable traits and disease risk together with other genes and environmental factors.

\medterm{Genotype Determination} Genotype determination is the use of laboratory testing to identify an individual's genetic variants or alleles at one or more locations.

\medterm{Genotype-Phenotype Correlation} The relationship between a person’s genetic makeup and the observable features, severity, age of onset, or clinical course of a condition; correlations may be incomplete because other genetic and environmental factors contribute.

\medterm{Gentamicin} An aminoglycoside antibiotic that interferes with bacterial protein production. It can injure the kidneys or hearing and is used with dose and blood-level monitoring.

\medterm{Gentamicin Sulfate} The sulfate salt of gentamicin, an aminoglycoside antibiotic that interferes with bacterial protein production. It can injure the kidneys or hearing and is used with careful dosing and monitoring.

\medterm{Gentian} A plant of the Gentiana genus whose roots or other parts have been used in herbal preparations, traditionally as a bitter digestive; clinical evidence and product quality vary.

\medterm{Gentian Violet} A purple antiseptic and dye with antifungal and antibacterial activity, used in selected topical applications but capable of irritating tissue and staining skin.

\medterm{Genu} The knee or a knee-shaped part of an anatomical structure; the term also describes a bend, such as the genu of the corpus callosum.

\medterm{Genu Recurvatum} Hyperextension of the knee beyond the usual straight position, caused by laxity, deformity, muscle imbalance, neurologic disease, or injury.

\medterm{Genu Valgum} An alignment abnormality in which the knees angle inward and may touch while the ankles remain apart. It may be developmental or caused by bone, joint, or metabolic disease.

\synonyms
Knock-Knees; Tibial Valgus

\medterm{Genu Varum} Outward bowing of the legs, commonly called bowlegs. It is often normal in young children and improves with growth, but persistent, worsening, unequal, painful, or severe bowing may result from rickets, growth-plate disease, or another condition.

\synonyms
Bowlegs; Bandy Legs

\medterm{Geographic Atrophy} An advanced form of dry age-related macular degeneration in which areas of light-sensing and supporting cells in the center of the retina are permanently lost. It causes slowly worsening blind spots and loss of detailed central vision.

\medterm{Geographic Tongue} Benign inflammatory condition of the tongue characterized by irregular red patches with white or yellow borders, alternating with areas of depapillation. Often asymptomatic and of unknown cause. May cause burning or irritation with certain foods.

\synonyms
Benign Migratory Glossitis
Benign Migratory Glossitis

\medterm{Geophagia} Persistent eating of soil or other nonfood substances; it may cause poisoning, infection, intestinal obstruction, or iron-deficiency anemia.

\medterm{Geophagy} The deliberate ingestion of soil or clay; it may cause lead or parasite exposure, gastrointestinal obstruction, or iron-deficiency anemia and can occur with pica.

\medterm{Geotrichum} Geotrichum is a genus of fungi that may colonize humans and can rarely cause superficial or invasive infection, especially in people with weakened immunity.

\medterm{Geriatric Assessment} A broad review of an older person’s medical conditions, medicines, mobility, memory, mood, daily activities, nutrition, safety, home situation, and goals. It identifies treatable problems and guides an individualized care plan.

\medterm{Geriatric Care Manager} A professional who helps coordinate medical, social, housing, and support services for an older adult. They may assess needs, organize care, communicate with families, and help plan for changing abilities.

\synonyms
Aging-Life Care Manager; Eldercare Coordinator

\medterm{Geriatric Dentistry} Dental care adapted to the needs of older adults, including tooth decay, gum disease, dry mouth, dentures, swallowing concerns, medicine effects, and examination for oral cancer.

\medterm{Geriatric Depression} Depression in an older adult may cause low mood, loss of interest, fatigue, sleep or appetite changes, slowed thinking, physical complaints, or memory problems. Medical illness, medicines, isolation, and bereavement may contribute.

\synonyms
Late-Life Depression; Depression in the Elderly; Late-Onset Depressive Disorder

\medterm{Geriatric Depression Scale} A questionnaire used to screen older adults for symptoms of depression. It can identify people with possible depression, but a positive result does not by itself establish depression or explain its cause.

\medterm{Geriatric Dream Changes} Dreams in older adults may become more vivid or be affected by medicines, sleep disorders, depression, or changes in the sleep cycle. Repeated frightening dreams, acting them out, confusion, or sudden change may occur.

\synonyms
Geriatric Sleep Patterns; Dreaming in Older Adults

\medterm{Geriatric Dysphagia} Difficulty swallowing in an older adult. It may cause coughing or choking during meals, food remaining in the mouth, weight loss, dehydration, or aspiration and can result from stroke, dementia, muscle disease, blockage, or medicines.

\synonyms
Presbyphagia; Age-Related Swallowing Dysfunction

\medterm{Geriatric Hospitalization} Hospital care for an older adult, with attention to medical illness as well as mobility, medicines, nutrition, hearing, vision, and cognition. Preventing falls, delirium, pressure injuries, and loss of independence supports safer recovery.

\synonyms
Inpatient Geriatric Care; Acute Care for Elders (ACE)

\medterm{Geriatrician} A physician specializing in the health care of older adults, including complex illness, multiple conditions, medications, function, and goals of care.

\medterm{Geriatric Infection} Infection in an older adult may cause confusion, weakness, reduced appetite, or a fall without causing fever. Common sites include the urine, lungs, skin, and bloodstream; illness can worsen quickly.

\synonyms
Infection in the Elderly; Late-Life Infection

\medterm{Geriatric Informed Consent} The process of ensuring that an older patient understands a proposed intervention, its benefits, risks, and alternatives before choosing freely. If decision-making ability is impaired, an authorized decision-maker may help according to local law.

\synonyms
Geriatric Decision-Making Capacity; Surrogate Consent in Older Adults

\medterm{Geriatric Medical Ethics} Ethical decision-making in the care of older adults, including respect for the person’s wishes, decision-making capacity, likely benefits and harms, fairness, privacy, and appropriate involvement of family or caregivers.

\medterm{Geriatric Nursing} Nursing care focused on the health, function, medication safety, cognition, mobility, and goals of older adults.

\medterm{Geriatric Nutrition} Nutrition care for older adults that considers aging, illness, medicines, frailty, dental or swallowing problems, appetite, and access to food. Goals include maintaining strength, muscle, hydration, function, and quality of life.

\medterm{Geriatric Osteoarthritis} Osteoarthritis occurring in later life, commonly affecting the knees, hips, hands, or spine
and causing joint pain, stiffness, and reduced movement. Management may include exercise,
weight management, medicines, assistive measures, or joint replacement when appropriate.

\synonyms
Age-Related Osteoarthritis; Senile Osteoarthritis

\medterm{Geriatric Rehabilitation} Therapy that helps an older adult regain or maintain movement, strength, independence, and safety after illness, injury, surgery, or functional decline. The plan is adapted to health, abilities, goals, and support needs.

\medterm{Geriatrics} The medical specialty focused on health, function, independence, and quality of life in older adults. Care considers multiple illnesses, medicines, memory, mobility, nutrition, mood, safety, goals, and support.

\medterm{Geriatric Screening} Checking older adults who have no symptoms for diseases or risks that may benefit from early detection. Examples include falls, depression, osteoporosis, and selected cancers.

\synonyms
Geriatric Screening; Senior Health Assessment

\medterm{Geriatric Sleep Changes} Sleep patterns and sleep problems in older adults, including insomnia, sleep apnea, circadian changes,
and medication-related disturbances.

\synonyms
Geriatric Sleep Disorders; Elder Sleep Management

\medterm{Geriatric Syndrome} A common problem in older adults caused by several interacting factors rather than one disease. Examples include falls, frailty, delirium, incontinence, immobility, and pressure injuries; causes may be medical, social, or both.

\medterm{Geriatric Trauma} Injury in an older adult, often involving falls, fractures, head injury, or bleeding. Age-related frailty, osteoporosis, balance problems, and medicines that increase bleeding can make injuries more serious and recovery slower.

\medterm{Germ} Alternate informal term; see \textit{Microorganism}.

\medterm{Germanium} A chemical element sometimes included in unproven health supplements. Germanium compounds have no established medical benefit and can cause kidney damage, nerve damage, or other serious toxicity when taken internally.

\medterm{German Measles} Rubella: a contagious viral infection caused by rubella virus, characterized by fever, lymph-node enlargement, and a fine rash. Infection during pregnancy can harm the fetus, causing congenital rubella syndrome with cardiac, auditory, and developmental defects.

\synonyms
Rubella; Three-Day Measles

\medterm{Germ Cell} A reproductive cell or precursor that gives rise to sperm or eggs and carries genetic material to the next generation.

\medterm{Germ Cell Tumor} A tumor arising from cells that normally develop into sperm or eggs. It may occur in the testicles, ovaries, or other sites such as the chest or abdomen; some types are highly treatable with surgery, chemotherapy, or both.

\medterm{Germ Cell Tumors} Tumors arising from cells that normally develop into sperm or eggs. They can occur in the testicles, ovaries, or other body sites and may be benign or cancerous, depending on the type.

\medterm{Germinoma} A usually slow-growing germ-cell cancer that most often develops in the pineal or suprasellar area of the brain. Related tumors occur in the testicles or ovaries; treatment commonly uses radiation, chemotherapy, surgery, or combinations.

\medterm{Germ Layer} One of the three early layers of an embryo—ectoderm, mesoderm, or endoderm—from which all major tissues and organs develop. Each layer gives rise to characteristic body structures during development.

\medterm{Germline Mutation} A genetic change present in an egg or sperm and therefore usually present in nearly every cell of the person. Some changes can be passed to children and may increase the risk of inherited disorders or cancer.

\medterm{Germline Testing} Genetic testing of blood, saliva, or another sample to look for changes that may be inherited or passed to children. Results can affect the person’s care, relatives, cancer screening, and reproductive decisions.

\medterm{Germline Variant} A DNA change present in cells that can give rise to eggs or sperm and therefore may be passed to children. It is usually present in most cells of the person, unlike a change limited to a tumor or other tissue.

\medterm{Germ Theory Of Disease} The evidence-based principle that particular microorganisms can cause particular infectious diseases. It underlies infection control, vaccination, antimicrobial treatment, sterilization, and methods used to identify disease-causing organisms.

\medterm{Gerontologist} A specialist studying aging and its biological, psychological, social, or economic effects. Gerontologists may work in research, healthcare, policy, counseling, education, housing, or community services for older people.

\medterm{Gerstmann-Straussler-Scheinker Syndrome} A rare, autosomal dominant inherited prion disease caused by mutations in the PRNP gene (most commonly P102L), characterized pathologically by multicentric PrP-amyloid plaques throughout the cerebellum and cerebral cortex, presenting with slowly progressive cerebellar ataxia, spasticity, and late dementia.
\synonyms
GSS Syndrome; GSS Prion Disease

\medterm{Gerstmann Syndrome} A neurological cognitive syndrome caused by a lesion of the dominant angular gyrus in the inferior parietal lobule (typically left hemisphere), characterized by the classic tetrad of agraphia (dysgraphia), acalculia (dyscalculia), finger agnosia, and right-left disorientation.
\synonyms
Angular Gyrus Syndrome; Dominant Parietal Lobe Syndrome

\medterm{Gesell Developmental Schedules} Standardized developmental assessment schedules that describe expected motor, language, adaptive, and social milestones in children; they support screening but do not by themselves diagnose a disorder.

\medterm{Gestation} The period of pregnancy during which an embryo and fetus develop in the uterus. Pregnancy is usually dated as about 40 weeks from the first day of the last menstrual period, or about 38 weeks from fertilization.

\medterm{Gestational Age} The length of a pregnancy measured in weeks from the first day of the last menstrual period. It is usually about two weeks longer than the time since fertilization and helps estimate development, due date, and appropriate care.

\synonyms
Menstrual Age; Fetal Postmenstrual Age

\medterm{Gestational Diabetes} Carbohydrate intolerance of variable severity with onset or first recognition during pregnancy. It occurs when pancreatic insulin secretion cannot overcome pregnancy-induced insulin resistance driven by placental hormones, including human placental lactogen, progesterone, and cortisol. Maternal risks include preeclampsia and operative delivery; fetal and neonatal risks include macrosomia, shoulder dystocia, birth injury, neonatal hypoglycemia, hyperbilirubinemia, and respiratory distress. Screening is typically performed between 24 and 28 weeks of gestation using an oral glucose challenge or tolerance test. Management centers on medical nutrition therapy, physical activity, and blood glucose monitoring, with insulin or oral agents added when targets are not met. Glucose levels typically normalize after delivery, but affected individuals carry a significantly increased lifetime risk of developing type 2 diabetes mellitus.

\synonyms
GDM; Gestational Diabetes Mellitus

\medterm{Gestational Diabetes Mellitus} Alternate medical term; see \textit{Gestational Diabetes}.

\medterm{Gestational Hypertension} New-onset high blood pressure after 20 weeks of pregnancy without protein in the urine or other features of preeclampsia. It may later develop into preeclampsia.

\medterm{Gestational Pyelonephritis} A kidney infection occurring during pregnancy, usually causing fever, flank pain, nausea, or urinary symptoms. It can lead to sepsis, preterm labor, or breathing complications.

\synonyms
Pyelonephritis in Pregnancy; Gravid Kidney Infection

\medterm{Gestational Sac} The first fluid-filled structure usually seen inside the uterus on an early pregnancy ultrasound. Its size and appearance help estimate pregnancy age and assess early development, but an early scan may not yet show an embryo.

\medterm{Gestational Trophoblastic Disease} A group of uncommon conditions caused by abnormal growth of cells that normally help form the placenta. It includes molar pregnancy and some pregnancy-related cancers and may cause unusually high pregnancy hormone levels, bleeding, or a uterus larger than expected.

\medterm{Gestational Trophoblastic Neoplasia} Persistent or cancerous growth of placenta-related cells after a pregnancy, miscarriage, abortion, or molar pregnancy. It is often detected by pregnancy-hormone levels that remain high or rise and is usually very responsive to specialist cancer treatment.

\medterm{Gestation Period} The time from conception to birth or, clinically, from the last menstrual period to delivery; human pregnancy is usually about 40 weeks from the last menstrual period.

\medterm{GFR} Glomerular filtration rate, an estimate of how well the kidneys filter blood. It is usually calculated from blood creatinine and other information, and a persistently low value may indicate chronic kidney disease.

\synonyms
Glomerular Filtration Rate; Estimated GFR (eGFR)

\medterm{Ghon Complex} The combination of a primary tuberculous lung focus (Ghon focus) and caseous granulomatous lymphadenitis in the draining regional hilar or mediastinal lymph nodes, representing the hallmark of primary pulmonary tuberculosis that often undergoes fibrocalcification.
\synonyms
Primary Tuberculous Complex; Ranke Precursor Complex

\medterm{Ghon Focus} A primary granulomatous tuberculous pulmonary lesion consisting of a circumscribed area of caseous necrosis typically located subpleurally in the mid-to-lower lung zone, representing the initial site of Mycobacterium tuberculosis infection.
\synonyms
Primary Focus of Ghon; Subpleural Tuberculous Focus

\medterm{Ghrelin} A hormone produced mainly in the stomach that increases hunger and helps regulate energy balance. It also stimulates growth-hormone release and changes with eating, fasting, sleep, and body weight.

\medterm{Gianotti-Crosti Syndrome} A childhood skin condition characterized by a papular or papulovesicular rash on the face, buttocks, and extremities, often following viral infections (EBV, hepatitis B, CMV). Self-limited, lasting 2-8 weeks. Also called papular acrodermatitis of childhood.

\synonyms
Papular Acrodermatitis of Childhood; Infantile Lichenoid Acrodermatitis

\medterm{Giant Cell Arteritis} Inflammation of arteries, usually in people over 50, often affecting vessels near the temples and eyes. It can cause new headache, scalp tenderness, jaw pain while chewing, or vision changes and may permanently damage sight.

\medterm{Giant Cell Arteritis Emergency} A severe presentation of suspected giant cell arteritis with visual symptoms, jaw or tongue pain during chewing, or reduced blood flow to the nervous system. It carries a risk of permanent visual loss.

\medterm{Giant Cell Myocarditis} A rare, severe inflammation of the heart muscle caused by an abnormal immune reaction. It can weaken the heart and cause dangerous heart rhythms or heart failure.

\medterm{Giant Cells} Large cells formed by fusion of macrophages or abnormal cellular growth, seen in granulomatous inflammation, foreign-body reactions, and some tumors.

\medterm{Giant Cell Tumor} A usually noncancerous but locally destructive bone tumor containing many large abnormal cells. It commonly affects bones near major joints, may cause pain or swelling, and can return after treatment or rarely spread to the lungs.

\medterm{Giant Cell Tumor Of Bone} A locally aggressive osteoclast-rich bone tumor, usually arising near the end of a long bone in skeletally mature adults. It may cause pain, swelling, cortical destruction, and occasional pulmonary metastases.

\medterm{Giant Cell Tumor Of Soft Tissue} A usually low-grade soft-tissue neoplasm containing numerous osteoclast-like giant cells, most often arising in the extremities of adults and generally having a favorable outcome after excision.

\medterm{Giant Cell Tumor Of Tendon Sheath} A common benign localized tenosynovial tumor, usually presenting as a slow-growing firm nodule on a finger or hand; recurrence can occur after excision.

\medterm{Giant Congenital Nevus} A very large dark or brown skin mark present at birth. It may grow with the child and has a higher risk of melanoma than smaller birthmarks.

\synonyms
Giant Pigmented Nevus; Garment Nevus; Bathing Trunk Nevus

\medterm{Giant Coronary Aneurysm} An unusually large bulge in an artery supplying the heart. It can form a clot, reduce blood flow, or rarely rupture.

\medterm{Giant Hypertrophic Gastritis} Severe thickening of the stomach lining and its folds, which can cause stomach pain, nausea, swelling, and loss of blood proteins. It may increase stomach-cancer risk.

\medterm{Giant Lymph Node Hyperplasia} A rare condition in which lymph-node tissue grows excessively. It may affect one area or many parts of the body and can cause enlarged nodes, fever, fatigue, anemia, weight loss, or organ problems.

\synonyms
Castleman Disease; Angiofollicular Lymph Node Hyperplasia

\medterm{Giant Platelet Syndrome} An inherited platelet disorder caused by deficiency of glycoprotein Ib/IX/V complex, resulting in large platelets and impaired platelet adhesion. Presents with mucocutaneous bleeding, thrombocytopenia, and giant platelets on blood smear.

\synonyms
Bernard-Soulier Syndrome; Platelet Glycoprotein Ib Deficiency

\medterm{Giant Right Atrium} Marked enlargement of the heart’s right upper chamber, usually from long-standing valve disease, a congenital heart defect, or abnormal rhythm. It can disturb the heartbeat, compress nearby structures, or allow clots to form and is assessed with heart imaging.

\medterm{Giardia} A microscopic parasite that infects the small intestine and spreads through contaminated water, food, or hands. It can cause foul-smelling diarrhea, cramps, bloating, nausea, weight loss, or no symptoms.

\medterm{Giardia Infection} An intestinal infection caused by Giardia parasites, usually spread through contaminated water, food, or contact with an infected person. It can cause prolonged foul-smelling diarrhea, cramps, bloating, weight loss, and poor nutrient absorption.

\medterm{Giardia Lamblia} The parasite that causes giardiasis. It is swallowed in contaminated water, food, or material from an infected person and can cause foul-smelling diarrhea, abdominal cramps, gas, bloating, and weight loss.

\medterm{Giardiasis} An intestinal infection caused by the parasite Giardia duodenalis. It spreads when contaminated water, food, surfaces, or material from an infected person is swallowed. It may cause diarrhea, foul-smelling or greasy stools, gas, abdominal cramps, nausea, bloating, weight loss, and poor nutrient absorption.

\medterm{Gibbus Deformity} A sharp, angular, structural kyphosis of the spine produced by the anterior collapse or destruction of one or more adjacent vertebral bodies, classically seen in tuberculosis of the spine (Pott disease) or severe osteoporotic wedge fractures.
\synonyms
Pott Gibbus; Angular Kyphosis

\medterm{Gibraltar Fever} An old name for brucellosis, an infection caused by Brucella bacteria, usually acquired through unpasteurized dairy products or contact with infected animals.

\medterm{Giemsa Stain} A laboratory dye used to show blood cells and some bacteria, parasites, and other organisms under a microscope. It can help identify malaria parasites and other infections in blood or tissue samples.

\medterm{Gigabecquerel} A unit of radioactivity equal to one billion nuclear disintegrations per second. It measures the activity of a radioactive source, not directly the radiation dose absorbed by a person.

\medterm{Gigantism} Excessive growth caused by too much growth hormone during childhood, before the bones finish lengthening. It usually results from a pituitary growth and causes unusual height and enlargement of body tissues.

\medterm{Gigli Saw} A flexible wire saw used during some orthopedic, neurosurgical, or amputation procedures to cut bone. It is handled carefully to limit injury to nearby soft tissues.

\medterm{Gilbert's Syndrome} A common inherited condition in which the liver processes bilirubin slowly, causing occasional mild yellowing of the eyes, especially during illness, fasting, or stress.

\medterm{Gilbert Syndrome} A common inherited condition in which the liver processes bilirubin more slowly than usual. It can cause occasional mild yellowing of the eyes or skin, especially during fasting, illness, dehydration, or stress, and usually does not harm the liver.

\medterm{Gilles De La Tourette Syndrome} A neurological disorder causing repeated involuntary movements and sounds called motor and vocal tics. Symptoms begin in childhood and often change over time.

\medterm{Gingiva} The firm pink tissue, commonly called the gums, that surrounds and protects the teeth and the bone around their roots. Healthy gingiva usually does not bleed with gentle brushing; redness, swelling, or bleeding may indicate gum disease.

\synonyms
Gums; Gingival Tissue

\medterm{Gingiva Disease} A disorder of the gums, including plaque-induced gingivitis, periodontitis, infection, immune disease, recession, or tumor.

\medterm{Gingival Disease} Inflammation of the gums, usually caused by plaque buildup. The gums may become red, swollen, tender, or bleed during brushing. It usually does not damage the supporting bone and can often improve with daily cleaning and professional dental care.

\synonyms
Marginal Gingivitis; Periodontal Gum Inflammation

\medterm{Gingival Hemorrhage} Bleeding from gum tissue, commonly due to inflammation and plaque but also caused by trauma, periodontal disease, medicines, or bleeding disorders.

\medterm{Gingival Hyperplasia} An abnormal enlargement or overgrowth of gingival tissue around the teeth that can impair cleaning, cause bleeding, and increase periodontal risk. Causes include plaque-related inflammation and medicines such as phenytoin, cyclosporine, and some calcium-channel blockers; hereditary and systemic conditions can also contribute.

\medterm{Gingival Infections} Infections of the gums, usually involving bacteria in dental plaque. They can cause red, swollen, tender, or bleeding gums, bad breath, pus, and sometimes spreading facial swelling; dental assessment helps prevent tissue and bone loss.

\medterm{Gingival Margin} The edge of the gum tissue that surrounds and touches a tooth. It normally forms a shallow seal; inflammation or recession can make it bleed or expose the tooth root.

\medterm{Gingival Melanoma} A rare, aggressive cancer of pigment-producing cells in the gums. It may appear as a dark or non-pigmented lump, patch, or ulcer; biopsy can identify it.

\medterm{Gingival Neoplasm} An abnormal growth in the gums that may be reactive, benign, precancerous, or cancerous. It may cause persistent swelling, ulceration, bleeding, pain, or a lump; examination and sometimes biopsy identify its nature.

\medterm{Gingival Overgrowth} Gingival overgrowth is an abnormal increase in gum tissue, which may be caused by medicines, inflammation, or systemic disease.

\medterm{Gingival Pocket} The small space between the gum and tooth surface, measured during a dental examination. A persistently deep pocket usually indicates loss of attachment from periodontal disease and can collect plaque and bacteria.

\medterm{Gingival Recession} Pulling back of the gums so that part of a tooth root becomes exposed. Gum disease, forceful brushing, tooth position, thin gum tissue, or treatment can cause it and may lead to sensitivity or root decay.

\medterm{Gingivectomy} Surgical removal of gum tissue to treat overgrowth, reshape the gums, expose decay, or improve access. Possible complications include bleeding, sensitivity, infection, or gum recession.

\medterm{Gingivitis} Inflammation of the gingiva characterized by redness, swelling, and bleeding with brushing or probing, usually due to dental plaque. It is generally reversible with effective oral hygiene and professional cleaning; it can precede periodontitis, but does not inevitably progress to it.

\synonyms
ANUG; Trench Mouth; Vincent Stomatitis

\medterm{Gingivitis Transmission} Gingivitis itself is not directly contagious, but the bacteria that cause it (particularly Porphyromonas gingivalis) can be transmitted through saliva. Poor oral hygiene is the primary risk factor. Prevention focuses on personal oral care.

\medterm{Gingivoplasty} Surgical reshaping or contouring of the gums to improve their form, remove diseased tissue, or make cleaning easier. It may be combined with gingivectomy or other periodontal treatment.

\medterm{Gingivostomatitis} Common viral infection of the mouth and gums, most often caused by primary herpes simplex virus infection but also associated with other infections, trauma, or systemic disease. It may cause painful ulcers, swollen bleeding gums, fever, and difficulty eating or drinking.

\medterm{Ginkgo Biloba} An herbal product derived from Ginkgo biloba leaves and marketed for cognitive or circulatory benefits; evidence is variable and it may increase bleeding risk or interact with medicines.

\medterm{Ginseng} A group of plants whose roots are used in traditional medicine and dietary supplements; products vary in composition and may interact with medicines or affect blood glucose and clotting.

\medterm{Ginseng Preparation} A product made from ginseng root and sold as a supplement. Products vary in strength and purity; interactions and side effects can occur.

\medterm{Girdle Pain} Pain that wraps around the chest or abdomen in a band-like pattern. It may result from nerve irritation, spinal disease, shingles, muscle injury, or internal-organ disease.

\medterm{Gitelman Syndrome} A rare inherited kidney disorder that causes the body to lose too much salt and magnesium in urine. It may lead to muscle cramps, weakness, tiredness, thirst, or abnormal heart rhythms.

\medterm{Githism} An old term for poisoning by corn cockle, a weed whose seeds can contaminate grain and cause stomach upset, weakness, or other toxic effects.


\synonyms
Corn-Cockle Poisoning; Agrostemma Githago Intoxication

\medterm{Glabella} The smooth area of the forehead between the eyebrows and above the bridge of the nose.

\medterm{Gland} An organ or group of cells that produces and releases substances. Endocrine glands release hormones into blood, while exocrine glands release substances through ducts or onto a surface.

\medterm{Glanders} A serious infection caused by the bacterium \textit{Burkholderia mallei}, which mainly affects horses and can rarely infect people exposed to infected animals or tissues. It may cause skin sores, fever, pneumonia, or bloodstream infection.

\medterm{Glandular Fever} Alternate name; see \textit{Infectious Mononucleosis}.

\medterm{Glans Penis} The expanded, sensitive distal end of the penis containing the external urethral opening and covered by the foreskin when present.

\medterm{Glanzmann Thrombasthenia} An inherited platelet function disorder caused by deficiency or dysfunction of platelet glycoprotein IIb/IIIa, producing mucocutaneous bleeding despite a normal platelet count.

\medterm{Glasgow Coma Scale} A bedside score describing a person’s level of consciousness by assessing eye opening, speaking, and movement. Scores range from 3 to 15; lower scores indicate less response, but medicines and injuries can affect the result.

\medterm{Glass Eye} A glass eye is an ocular prosthesis made to replace the appearance of an eye that is missing or severely damaged.

\medterm{Glass Ionomer} A tooth-colored dental material that bonds to tooth structure and slowly releases fluoride. It is used for selected fillings, repairs, protective bases, and children’s teeth; durability depends on the location and amount of stress.

\medterm{Glassy Cell Carcinoma} A rare poorly differentiated adenosquamous carcinoma of the cervix with distinctive glassy cytoplasm, prominent nucleoli, and aggressive clinical behavior.

\medterm{Glaucoma} A group of eye diseases in which the optic nerve is damaged, often in association with high pressure inside the eye. Types include open-angle, angle-closure, normal-tension, secondary, and congenital glaucoma. It can cause permanent loss of side or central vision, while early disease may have no symptoms.

\medterm{Glaucoma Drainage Implant} A surgically implanted tube-and-plate device that diverts aqueous humor from the anterior chamber to a subconjunctival reservoir, lowering intraocular pressure in selected glaucomas.

\medterm{Glaze Poisoning} Toxic exposure to ceramic or pottery glaze, which may contain lead, cadmium, or other heavy metals. Symptoms depend on the metal involved and may include abdominal pain, neurological symptoms, and organ damage. Severe cases may be treated with chelation therapy.

\synonyms
Ceramic Glaze Poisoning; Lead or Cadmium Glaze Toxicity

\medterm{Gleason Grading System} A histological grading system for prostate cancer based on glandular pattern differentiation. Scores historically ranged from 2 to 10, but current reporting summarizes them as Grade Groups 1 to 5 (scores 6 to 10), with higher groups indicating more aggressive cancer. Used to guide treatment decisions and predict prognosis.

\synonyms
Gleason Score; Histological Prostate Grade

\medterm{Gleason Score} A grading system estimating prostate cancer aggressiveness by examining tissue patterns under a microscope and combining the two most common patterns.

\synonyms
Gleason Grade Group; Prostate Adenocarcinoma Score

\medterm{Gleason Tumor Grade} A histopathologic grade of prostate adenocarcinoma based on the architectural pattern of the tumor glands; it contributes to the Gleason score and Grade Group.

\medterm{Glenn Procedure} A staged heart operation for some children born with only one effective pumping chamber. It connects the large vein returning blood from the upper body directly to the lung arteries, allowing blood to pick up oxygen while reducing the workload on the heart.

\synonyms
Bidirectional Glenn (BDG) Procedure; Cavopulmonary Anastomosis

\medterm{Glenohumeral Joint} The ball-and-socket shoulder joint formed by the upper-arm bone and shoulder blade. It allows the arm to move in many directions but is relatively easy to dislocate because it sacrifices some stability for movement.

\medterm{Glenoid} The shallow socket on the shoulder blade that joins the head of the upper-arm bone to form the shoulder joint. Its shape and surrounding labrum help balance mobility with stability.
\synonyms
Glenoid Cavity; Scapular Glenoid Fossa

\medterm{Glenoid Cavity} The shallow socket on the shoulder blade that receives the rounded head of the upper-arm bone. A rim of tough tissue deepens the socket, but the joint’s wide range of movement also makes it prone to dislocation.

\medterm{Glenoid Labrum} A fibrocartilaginous rim attached to the peripheral margin of the glenoid fossa of the scapula, which deepens the shallow glenoid cavity to enhance glenohumeral joint stability and serves as the superior anchor for the long head of the biceps brachii tendon.

\synonyms
Glenoid Ligament

\medterm{Gliadin} A group of wheat storage proteins that, together with glutenin, form gluten; gliadin peptides trigger the immune response in celiac disease.

\medterm{Glial Cell} A non-neuronal cell that supports, nourishes, insulates, and regulates neurons; examples include astrocytes, oligodendrocytes, microglia, and Schwann cells.

\medterm{Glial Cells} Support cells of the brain and spinal cord that nourish and protect nerve cells, form insulation, maintain the chemical environment, and help repair tissue. Some brain tumors arise from glial cells.

\medterm{Glial Fibrillary Acidic Protein} A protein found mainly in support cells of the brain and nervous system. Laboratory tests can use it to help identify some brain tumors or detect injury to support cells after brain trauma.

\synonyms
GFAP

\medterm{Glibenclamide} An oral sulfonylurea medicine, also called glyburide, used for type 2 diabetes. It stimulates the pancreas to release insulin and can cause prolonged, sometimes severe low blood sugar and weight gain.

\medterm{Gliclazide} Gliclazide is a sulfonylurea for type 2 diabetes that stimulates pancreatic beta-cell insulin release; hypoglycaemia, weight gain, and dose accumulation in organ impairment are important considerations.

\medterm{GLIM Criteria} An international framework for diagnosing malnutrition. It considers body changes such as weight or muscle loss together with causes such as poor intake, poor absorption, or illness-related inflammation and grades severity according to the findings.

\medterm{Glimepiride} An oral sulfonylurea medicine for type 2 diabetes that increases insulin release from functioning pancreatic beta cells. Low blood sugar is the main serious risk, particularly with missed meals, kidney disease, or other glucose-lowering medicines.

\medterm{Glioblastoma} A fast-growing cancerous tumor of the central nervous system. It forms from glial tissue, which supports nerve cells, and usually develops in the brain of an adult; it can also occur in the spinal cord. The tumor can grow into nearby healthy tissue. Symptoms depend on its location and may include headache, seizures, weakness, changes in speech, memory, vision, balance, or behaviour. Also called glioblastoma multiforme (GBM).

\medterm{Glioblastoma Multiforme} An aggressive cancer arising from supporting cells of the brain or spinal cord. It may cause headache, seizures, personality or memory changes, weakness, or speech problems. Treatment often combines surgery, radiation, and medicines.

\medterm{Glioma} A tumor arising from supporting cells of the brain or spinal cord. Some gliomas grow slowly, while others are aggressive; symptoms and treatment depend on the tumor’s type, location, molecular features, and spread.

\medterm{Gliomatosis Cerebri} A diffuse infiltrative growth pattern of glioma involving multiple brain regions, rather than a separate tumor type in current classifications; it may cause progressive cognitive, motor, or seizure-related symptoms.

\medterm{Gliosarcoma} A rare aggressive brain cancer containing both glial and connective-tissue-like tumor features. Symptoms depend on location and may include headache, seizures, weakness, or cognitive changes; tissue examination identifies the tumor type.

\medterm{Gliosis} A reactive increase in glial cells after injury, infection, ischemia, or other damage to the central nervous system.

\medterm{Glipizide} Glipizide is a short-acting sulfonylurea that lowers blood glucose by stimulating insulin secretion from functioning pancreatic beta cells; hypoglycaemia is the principal adverse effect.

\medterm{Glisson's Capsule} A thin layer of connective tissue surrounding the liver and extending between its sections. It contains nerves and blood vessels and can cause pain when the liver becomes swollen or inflamed.

\medterm{Global Aphasia} A severe language disorder caused by extensive damage to the dominant cerebral hemisphere, impairing speech production, comprehension, reading, and writing.

\medterm{Global Developmental Delay} Noticeably delayed development in at least two areas, such as movement, speech, learning, social skills, or daily activities, in a child younger than five years. Causes vary and may include genetic, neurologic, hearing, vision, or environmental factors; assessment guides support and further testing.

\synonyms
GDD; Neurodevelopmental Delay

\medterm{Global Health} The study and practice of improving health across countries and populations. It includes problems that cross borders, such as outbreaks, unequal access to care, pollution, climate change, malnutrition, and differences in life expectancy.

\medterm{Global Longitudinal Strain} An echocardiogram measurement showing how much the heart muscle shortens lengthwise as it contracts. It can detect early weakening of the left ventricle before the usual pumping measurement changes and is interpreted with the complete heart scan and clinical findings.

\synonyms
GLS; 2D Speckle-Tracking GLS; Left Ventricular Longitudinal Strain

\medterm{Globin} A protein that joins with heme to form hemoglobin and related oxygen-carrying molecules. Different globin chains help determine the type of hemoglobin made at different stages of life.

\medterm{Globozoospermia} A rare sperm disorder in which sperm heads are round and lack a normal acrosome, causing impaired fertilization and male infertility.

\medterm{Globulin} A group of blood proteins that includes antibodies and proteins involved in transport, inflammation, and clotting. Blood tests may divide globulins into several groups; an abnormal level can occur with infection, immune disease, liver disease, kidney disease, or some cancers.

\medterm{Globus} A persistent or intermittent sensation of a lump or tightness in the throat without a true obstructing mass, often called globus pharyngeus.

\medterm{Globus Hystericus} A sensation of a lump or tightness in the throat without a structural obstruction, now called globus sensation; reflux, muscle tension, and anxiety may contribute.

\medterm{Globus Pallidus} The globus pallidus is a group of deep brain structures involved in regulating movement as part of the basal ganglia.

\medterm{Globus Sensation} The persistent or intermittent sensation of a lump or foreign body in the throat without
an anatomic lesion or dysphagia, also termed globus pharyngeus, typically noticed
between meals. There is no impedance to swallowing and no weight loss; the feeling
worsens with stress, anxiety, or laryngopharyngeal reflux.

\medterm{Glomangiomyoma} A benign glomus tumor showing both vascular and smooth-muscle differentiation, usually presenting as a small painful subcutaneous nodule.

\medterm{Glomangiosarcoma} A rare cancerous glomus tumor with abnormal cells and invasive growth. It may occur in the skin or deeper soft tissue.

\medterm{Glomerular Basement Membrane} The thick extracellular matrix layer of the renal filtration barrier interposed between glomerular endothelial cells and podocyte foot processes, composed of type IV collagen, laminin, and heparan sulfate proteoglycans that impart mechanical structural support and negative charge selectivity.

\synonyms
GBM

\medterm{Glomerular Capsule} A cup-shaped structure around a tiny group of blood vessels in the kidney. It collects the fluid filtered from blood and passes it into a kidney tubule, where useful water and substances are reclaimed and waste is removed.

\medterm{Glomerular Filtrate} The watery fluid filtered from blood into the kidney tubules by a glomerulus, a small group of kidney blood vessels. It normally contains water and small dissolved substances but very little protein or no blood cells.

\medterm{Glomerular Filtration} The first step in urine formation, in which pressure moves water and small dissolved substances from blood through kidney filters into the tubules. Blood cells and most proteins normally remain in the bloodstream.

\medterm{Glomerular Filtration Rate} An estimate of how much blood fluid the kidneys filter each minute. It is calculated mainly from blood creatinine and other factors; a persistently low result can indicate chronic kidney disease, but the result is interpreted with other findings.

\medterm{Glomerular Hyperfiltration} An abnormally elevated single-nephron or whole-kidney glomerular filtration rate exceeding physiological upper limits, representing an early hemodynamic adaptive response in early diabetic kidney disease, obesity, and after unilateral nephrectomy that eventually accelerates glomerulosclerosis.

\medterm{Glomeruli} Microscopic clusters of capillaries in the kidneys where blood filtration begins. Each glomerulus is surrounded by a capsule that collects filtered fluid and passes it into a kidney tubule.

\synonyms
Renal Glomeruli; Malpighian Corpuscles

\medterm{Glomerulomegaly} Abnormal enlargement of kidney glomeruli, seen with diabetes, obesity-related hyperfiltration, congenital syndromes, or other causes of glomerular stress.

\medterm{Glomerulonephritis} Inflammation or other immune-related injury of the kidney’s filtering units. It can cause blood or protein in urine, swelling, high blood pressure, and reduced kidney function and may follow an infection or occur with an autoimmune disease.

\synonyms
GN; Glomerular Nephritis

\medterm{Glomerulonephritis Acute} Sudden inflammation of the kidney glomeruli that can cause blood or protein in urine, edema, hypertension, and reduced kidney function.

\medterm{Glomerulopathy} Any disease affecting the kidney’s filtering units. It may cause blood or protein in urine, swelling, high blood pressure, reduced filtration, or kidney failure and may result from diabetes, infection, immune disease, inherited disorders, or other injury.

\medterm{Glomerulosclerosis} Scarring and hardening of the kidney’s filtering units. It can reduce filtration and cause protein in urine, swelling, high blood pressure, or kidney failure; causes include diabetes, inflammation, genetic disease, and other injury.

\medterm{Glomerulus} A tuft of tiny blood vessels inside a kidney capsule where water and small dissolved substances are filtered from blood to begin urine formation. Damage can allow blood or excess protein into urine and reduce kidney function.

\medterm{Glomus Jugulare Tumor} A usually noncancerous but locally invasive growth arising near the jugular vein at the skull base. It may cause pulsating noise in the ear, hearing loss, dizziness, or weakness of nearby nerves.

\synonyms
Jugular Paraganglioma; Tympanojugular Glomus

\medterm{Glomus Tumor} A usually harmless growth arising from cells that help control blood flow and temperature, most often beneath a fingernail or in a finger. It typically causes severe pinpoint pain and sensitivity to cold or pressure.

\medterm{Glomus Tympanicum} A benign vascular tumor arising from the glomus bodies on the promontory of the middle
ear, presenting with pulsatile tinnitus, hearing loss, and a reddish mass behind the
tympanic membrane.

\medterm{Glomus Tympanicum Tumor} A paraganglioma arising from the glomus bodies on the promontory of the middle ear. Benign tumor causing pulsatile tinnitus and conductive hearing loss. Diagnosis by otoscopy and imaging. Treatment includes surgery or observation.

\synonyms
Middle Ear Paraganglioma; Tympanic Glomus

\medterm{Glossal} Relating to the tongue, the muscular organ used for chewing, swallowing, tasting, and speaking.

\medterm{Glossalgia} Pain or burning of the tongue, with causes including trauma, infection, nutritional deficiency, neuropathy, medication, or burning-mouth syndrome.

\medterm{Glossectomy} Surgical removal of all or part of the tongue, performed for malignancies, severe trauma, or an unusually large tongue. Procedures range from partial to total removal. Speech and swallowing rehabilitation may follow.

\medterm{Glossinidae} The family of flies belonging to the genus Glossina, commonly known as tsetse flies. These insects are vectors of trypanosomes that cause African trypanosomiasis (sleeping sickness) in humans and nagana in animals.

\synonyms
Tsetse Flies; Glossina Vectors

\medterm{Glossitis} Inflammation of the tongue, which can cause pain, swelling, and changes in tongue appearance. Causes include nutritional deficiencies (iron, B12, folate), infections (bacterial, viral, fungal), allergic reactions, and autoimmune conditions. Treatment depends on the underlying cause.

\medterm{Glossodynia} Persistent burning, soreness, or tingling in the tongue, lips, palate, or other mouth tissues, often without visible changes. It may occur with dry mouth, altered taste, nutritional deficiency, medicines, nerve problems, or other illness.

\synonyms
Burning Mouth Syndrome; Orofacial Dysesthesia

\medterm{Glossopharyngeal Nerve} The ninth cranial nerve, carrying sensation and taste from the back of the tongue and throat, signals from blood-pressure and blood-chemistry sensors, and motor control for one throat muscle. It also stimulates saliva production.

\medterm{Glossopharyngeal Neuralgia} Recurrent, brief attacks of severe stabbing or electric-shock-like pain in the throat,
tonsillar region, base of the tongue, or ear, caused by irritation of the
glossopharyngeal nerve. Swallowing, coughing, or talking may trigger attacks.

\medterm{Glossoptosis} Backward displacement of the tongue that can narrow the upper airway, particularly in infants with craniofacial disorders.

\medterm{Glottis} The part of the voice box containing the vocal folds and the opening between them. It opens for breathing, closes to protect the airway during swallowing, and vibrates to produce the voice.

\medterm{GLP-1 Agonists} Medicines that mimic a natural gut hormone. They help lower blood glucose after meals, slow stomach emptying, and reduce appetite. They are used for selected people with diabetes or obesity and commonly cause nausea, vomiting, diarrhea, or constipation.

\medterm{GLP-1 Receptor Agonists} A class of medicines that activate the glucagon-like peptide-1 receptor. They increase glucose-dependent insulin release, reduce inappropriate glucagon release, slow gastric emptying, and increase fullness. They are used for selected people with type 2 diabetes or chronic weight management. Examples include semaglutide, liraglutide, and dulaglutide.

\synonyms
GLP-1 Agonists; Incretin Mimetics

\medterm{GLS} Alternate name; see \textit{Global Longitudinal Strain}.

\medterm{Glu} Commonly an abbreviation for glucose in laboratory or clinical notes; in biochemical notation, Glu can also denote the amino acid glutamic acid, so context matters.

\medterm{Glucagon} A hormone made by pancreatic alpha cells that raises blood glucose when it falls too low. It prompts the liver to release stored glucose and can be given for severe hypoglycemia.

\medterm{Glucagon Blood Test} Laboratory measurement of glucagon levels in the blood, used to evaluate suspected glucagonoma, hypoglycemia, or other pancreatic endocrine disorders. Levels are affected by fasting, meals, and medications.

\synonyms
Serum Glucagon Level; Plasma Glucagon Assay

\medterm{Glucagon-like Peptide-1} A hormone released mainly by the intestine after eating. It helps regulate blood glucose by enhancing glucose-dependent insulin secretion, reducing glucagon release, slowing stomach emptying, and influencing appetite.

\medterm{Glucagonoma} A rare pancreatic neuroendocrine tumor that secretes excess glucagon, causing diabetes mellitus (mild, new onset), dermatitis (necrolytic migratory erythema, characteristic painful, erythematous, blistering rash on groin, perineum, extremities), diarrhea, and deep vein thrombosis.

\medterm{Glucagon Signaling} The chain of cell signals started when glucagon binds its receptor. These signals increase cyclic AMP and activate proteins that help the liver release glucose and the body maintain energy balance.

\medterm{Glucan} A complex carbohydrate made of many linked glucose molecules. Different glucans occur in plants, fungi, bacteria, and foods and may act as dietary fiber or interact with the immune system.

\medterm{Glucarpidase} An enzyme that rapidly breaks down methotrexate in the bloodstream and is used for toxic methotrexate concentrations when renal clearance is inadequate.

\medterm{Glucocerebrosidase Deficiency} Gaucher Disease: an inherited lysosomal storage disorder caused by deficient glucocerebrosidase, leading to lipid accumulation in cells and possible enlargement of the spleen and liver, anemia, bone disease, or neurologic disease.

\synonyms
Gaucher Disease; Acid Beta-Glucosidase Deficiency

\medterm{Glucocorticoid} A corticosteroid hormone or medicine that regulates inflammation, immunity, metabolism, and stress responses; prolonged treatment has systemic adverse effects.

\medterm{Glucocorticoid-Induced Osteoporosis} Bone loss and increased fracture risk caused or accelerated by prolonged or high-dose glucocorticoid treatment; risk depends on dose, duration, baseline bone health, and other factors.

\medterm{Glucocorticoid Remediable Aldosteronism} An autosomal dominant cause of secondary endocrine hypertension caused by an unequal crossing-over mutation that fuses the ACTH-responsive promoter of the 11-beta-hydroxylase gene (\textit{CYP11B1}) to the aldosterone synthase coding sequence (\textit{CYP11B2}), rendering aldosterone secretion suppressible by dexamethasone.

\synonyms
GRA; Familial Hyperaldosteronism Type I

\medterm{Glucocorticoid-Remediable Aldosteronism} An autosomal dominant form of secondary hypertension caused by an unequal crossover genetic recombination fusing the promoter of the 11$\beta$-hydroxylase gene (\textit{CYP11B1}) to the coding region of the aldosterone synthase gene (\textit{CYP11B2}). This places aldosterone production under the ectopic regulatory control of adrenocorticotropic hormone (ACTH), resulting in hyperaldosteronism, hypokalemia, and early-onset hypertension that responds promptly to exogenous glucocorticoid suppression.

\synonyms
GRA; Familial Hyperaldosteronism Type I

\medterm{Glucocorticoid Replacement Therapy} Treatment of adrenal insufficiency with a glucocorticoid medicine to replace the cortisol the body lacks. Doses may vary with symptoms, illness, and physical stress.

\medterm{Glucocorticoids} Steroid hormones made by the adrenal glands, along with medicines that act in the same way. They reduce inflammation and immune activity but can cause infection, high blood glucose, bone loss, or adrenal suppression with prolonged use.

\medterm{Glucocorticoid-Sparing Therapy} Use of a non-glucocorticoid treatment strategy to control inflammation while reducing
cumulative corticosteroid exposure. It aims to limit osteoporosis, infection, diabetes,
hypertension, adrenal suppression, and other glucocorticoid complications.

\medterm{Glucometer} A handheld device that measures blood glucose, usually from a small capillary-blood sample obtained by finger prick.

\medterm{Gluconates} Salts or esters of gluconic acid used in some medicines and supplements. Examples include calcium gluconate and ferrous gluconate; their effects, dosing, and risks depend on the specific compound.

\medterm{Gluconeogenesis} A metabolic pathway that generates glucose from non-carbohydrate precursors such as lactate, glycerol, and glucogenic amino acids. It occurs mainly in the liver and, to a variable extent, in the kidneys and other tissues.

\medterm{Glucosamine} A naturally occurring compound used in cartilage formation and sometimes taken as a supplement for joint symptoms.

\medterm{Glucose} A simple sugar that is the main energy source for the body’s cells. It comes from digesting carbohydrates and can also be released from stored energy or made by the liver. Hormones, especially insulin, help keep blood glucose within a healthy range.

\medterm{Glucose-6-Phosphatase} An enzyme in the liver and kidneys that releases glucose from glucose-6-phosphate. Inherited deficiency causes glycogen-storage disease type I, which can produce low blood glucose, enlarged liver, and metabolic problems.

\medterm{Glucose-6-Phosphate Dehydrogenase Deficiency} An X-linked red-blood-cell enzyme disorder. Infections, fava beans, or certain medicines can cause oxidative damage and sudden destruction of red blood cells, leading to hemolytic anemia, jaundice, dark urine, fatigue, or shortness of breath.

\medterm{Glucose-6-Phosphate Dehydrogenase Test} Laboratory test to measure G6PD enzyme activity in red blood cells, used to diagnose G6PD deficiency. Important before prescribing certain medications (primaquine, dapsone) that can cause hemolytic anemia in deficient individuals.

\synonyms
G6PD Assay; Erythrocyte Enzyme Screen

\medterm{Glucose Challenge Test} A screening test for gestational diabetes mellitus, commonly using a nonfasting 50-g glucose drink followed by a one-hour plasma-glucose measurement.
An abnormal screening result is followed by a diagnostic oral glucose tolerance test; the threshold depends on the protocol used.

\synonyms
GCT; One-Hour Glucose Challenge; O'Sullivan Screening Test

\medterm{Glucose Intolerance} Impaired handling of glucose resulting in blood-glucose concentrations above normal but
below the diagnostic threshold for diabetes in a defined testing context. It often
reflects insulin resistance or inadequate insulin secretion and is associated with
increased risk of type 2 diabetes and cardiovascular disease.

\medterm{Glucose Management Indicator} An estimate of the hemoglobin A1c value calculated from average continuous-glucose-monitoring readings. It helps summarize recent glucose exposure but is not the same as a laboratory A1c and may differ from it.

\medterm{Glucose Screening Tests During Pregnancy} Laboratory tests to detect gestational diabetes mellitus, including the oral glucose tolerance test and glucose challenge test. They are usually performed between 24 and 28 weeks of gestation; abnormal results are followed by additional testing or management.

\synonyms
Gestational Diabetes Screening; Maternal Glucose Screen

\medterm{Glucose Sensor} A small sensor placed under the skin that measures glucose in the fluid between cells. Used in continuous glucose monitoring, it sends repeated readings to a receiver or compatible device and can show trends and alerts.

\medterm{Glucose Test} A laboratory or point-of-care test that measures glucose in blood or another sample. It is used to screen for, diagnose, or monitor diabetes and low or high blood sugar; interpretation depends on timing, fasting status, and clinical context.

\medterm{Glucose Tolerance Test} A test of how the body handles a measured glucose drink. Blood samples are taken before and at set times afterward to help diagnose diabetes or gestational diabetes; thresholds depend on the protocol.

\synonyms
Oral Glucose Tolerance Test (OGTT); 3-Hour Diagnostic OGTT

\medterm{Glucose Transport} Movement of glucose across cell membranes through specialized transport proteins. Insulin increases glucose entry into muscle and fat cells, while other transporters allow glucose to enter the brain, liver, kidneys, and intestines.

\medterm{Glucose Transporter} A membrane protein that moves glucose across a cell membrane; transporter defects or altered regulation can affect glucose metabolism.

\medterm{Glucose Urine Test} Laboratory test to detect glucose in the urine (glycosuria). Normally absent; presence indicates blood glucose exceeding renal threshold, suggesting diabetes mellitus or renal tubular dysfunction. May be detected by dipstick or laboratory analysis.

\synonyms
Urinary Glucose; Glycosuria Screen

\medterm{Glucosuria} The presence of glucose in urine, most commonly from blood glucose exceeding the renal reabsorption threshold, diabetes, pregnancy, or a proximal tubular reabsorption defect.

\medterm{Glucosylceramidase} A lysosomal enzyme that breaks down glucosylceramide, a fat-containing cell substance. Inherited deficiency causes Gaucher disease, allowing fatty material to accumulate in the spleen, liver, bone marrow, bones, or nervous system.

\medterm{Glucuronate} The anion or salt of glucuronic acid, a sugar-acid component of connective-tissue substances and a molecule used in hepatic conjugation reactions.

\medterm{Glucuronic Acid} A sugar acid used by the liver to conjugate bilirubin, drugs, hormones, and toxins, increasing their water solubility for excretion.

\medterm{Glucuronide} A compound formed when glucuronic acid is conjugated to a drug, hormone, bilirubin, or other molecule, usually increasing water solubility and facilitating excretion.

\medterm{Glucuronyl Transferase} UDP-glucuronosyltransferase: a liver enzyme involved in phase II drug metabolism, catalyzing glucuronidation of bilirubin, drugs, and toxins to make them water-soluble for excretion. Deficiency causes Gilbert syndrome or Crigler-Najjar syndrome.

\synonyms
UDP-Glucuronosyltransferase; UGT Enzyme

\medterm{Glue Ear} Persistent fluid in the middle ear without acute infection, often causing conductive hearing loss and a feeling of blockage. It is common in children.

\synonyms
Secretory Otitis Media; Otitis Media with Effusion (OME)

\medterm{GLUT4 Transporter} The major insulin- and exercise-regulated glucose transporter isoform expressed primarily in striated muscle (skeletal and cardiac) and adipose tissue, which translocates from intracellular vesicles to the plasma membrane upon insulin stimulation to facilitate cellular glucose uptake.

\medterm{Glutamate} The main excitatory chemical messenger in the brain and spinal cord. It helps nerve cells communicate, learn, and form memories; excessive glutamate activity can contribute to nerve-cell injury.

\medterm{Glutamate Decarboxylase} Glutamate decarboxylase is an enzyme that converts glutamate into the neurotransmitter GABA; antibodies against it are associated with some autoimmune disorders.

\medterm{Glutamate Receptor} A neuronal receptor activated by glutamate, the major excitatory neurotransmitter, including NMDA, AMPA, kainate, and metabotropic receptor classes.

\medterm{Glutamic Acid} An amino acid that acts as a metabolic intermediate and neurotransmitter precursor; its charged form, glutamate, mediates most fast excitatory signaling in the brain.

\medterm{Glutamic Acid Decarboxylase} An enzyme that converts glutamate to GABA; antibodies to it can occur in type 1 diabetes.

\synonyms
GAD; Glutamate Decarboxylase

\medterm{Glutamine} An amino acid used to build proteins, transport nitrogen, and provide energy for some cells, especially cells in the intestine and immune system. The body usually makes enough, but needs may change during severe illness.

\medterm{Glutaral} Glutaral is a chemical disinfectant and tissue fixative that can irritate the skin, eyes, and airways.

\medterm{Glutaric Acidemia Type 1} A rare inherited disorder in which the body cannot break down certain amino acids, allowing harmful substances to build up. Babies may have a large head and can develop sudden brain injury, abnormal movements, or stiffness during illness; newborn screening and dietary treatment can reduce harm.

\synonyms
Glutaric Aciduria Type 1 (GA-1); Glutaryl-CoA Dehydrogenase Deficiency

\medterm{Glutaric Aciduria} An inherited metabolic disorder in which the body cannot properly break down certain amino acids. Type I can cause brain injury, movement problems, and enlarged head size.

\medterm{Glutathione} A small antioxidant molecule made from glutamate, cysteine, and glycine. Reduced glutathione helps control oxidative stress and helps the liver neutralize reactive chemicals, including NAPQI produced from acetaminophen. N-acetylcysteine can restore glutathione after acetaminophen overdose.

\medterm{Glutathione Disulfide} The oxidized form of glutathione, produced after glutathione neutralizes reactive oxygen compounds. Cells normally convert it back to active glutathione to maintain antioxidant protection.

\medterm{Glutathione Peroxidase} A selenium-dependent antioxidant enzyme family that catalyzes the reduction of toxic hydrogen peroxide and organic lipid hydroperoxides to water and corresponding alcohols, utilizing reduced glutathione (GSH) as an essential electron donor to protect cell membranes from oxidative injury.

\synonyms
GPx

\medterm{Glutathione Transferase} A family of enzymes that attaches glutathione to selected medicines, toxins, and products of cell metabolism. This usually makes them less reactive and easier for the body to process or remove.

\medterm{Gluteal} Relating to the buttock or the gluteal muscles, which extend and rotate the hip and help stabilize the pelvis.

\medterm{Gluteal Gait} An abnormal walking gait caused by weakness or paralysis of the hip abductor muscles (gluteus medius and minimus) innervated by the superior gluteal nerve, manifested by a downward pelvic tilt toward the unsupported swing leg during single-leg stance, compensated by leaning the trunk toward the affected stance side.
\synonyms
Trendelenburg Gait; Gluteus Medius Gait

\medterm{Gluteal Injury} Damage to the muscles, tendons, nerves, or soft tissues of the buttock caused by trauma, overuse, injection, or falls. Pain, bruising, weakness, difficulty walking, severe swelling, numbness, inability to bear weight, or bowel or bladder changes may occur.

\medterm{Gluteal Muscle} One of the large muscles of the buttock, especially the gluteus maximus, medius, or minimus. These muscles extend, rotate, and move the hip and help stabilize the pelvis during standing and walking.

\medterm{Gluteal Tendinopathy} Tendon degeneration, microtrauma, and tearing affecting the insertion of the gluteus medius or gluteus minimus tendons onto the greater trochanter of the femur, representing the predominant cause of greater trochanteric pain syndrome in middle-aged and older adults.

\medterm{Gluten} A complex of storage proteins (prolamins and glutelins) found in wheat (gliadin), barley
(hordein), and rye (secalin). Gluten provides elasticity to dough. In susceptible
individuals, gluten triggers an autoimmune response (celiac disease) or non-autoimmune
sensitivity (non-celiac gluten sensitivity, NCGS).

\medterm{Gluten Enteropathy} Celiac disease, an immune reaction to gluten that damages the small intestine and can cause diarrhea, weight loss, anemia, fatigue, or poor nutrient absorption.

\synonyms
Celiac Disease; Celiac Sprue

\medterm{Gluten-Free Diet} A diet that avoids wheat, barley, rye, and foods contaminated with them. It is essential for celiac disease and may help selected people with non-celiac gluten sensitivity; poorly planned diets can lack fiber or nutrients.

\medterm{Gluten Intolerance} Alternate name; see \textit{Non-Celiac Gluten Sensitivity}.

\medterm{Gluten-Sensitive Enteropathy} Alternate name; see \textit{Celiac Disease}.

\medterm{Gluten Sensitivity} Alternate name; see \textit{Nonceliac gluten sensitivity}.

\medterm{Gluteus} Relating to the buttock muscles. The gluteus maximus, medius, and minimus move and stabilize the hip and pelvis during standing, walking, climbing, and other movements.

\medterm{Gluteus Maximus Muscle} Largest, most superficial gluteal muscle forming the buttock bulk, from posterior iliac
gluteal line, sacrum, coccyx, and sacrotuberous ligament to the iliotibial tract and
gluteal tuberosity. Innervated by inferior gluteal nerve (L5-S2). Primary hip extensor
and lateral rotator.

\medterm{Gluteus Medius Muscle} Thick fan-shaped muscle between maximus and minimus, from external ilium to the greater
trochanter. Innervated by superior gluteal nerve (L4-S1). Primary hip abductor essential
for pelvic stability during single-leg stance. Superior gluteal nerve injury causes
Trendelenburg sign with contralateral pelvic drop.

\medterm{Gluteus Minimus Muscle} The smallest and deepest gluteal muscle, extending from the outer ilium to the greater trochanter. It abducts and
medially rotates the hip and helps stabilize the pelvis during walking.

\medterm{Glyburide} An oral sulfonylurea that stimulates pancreatic insulin release for type 2 diabetes; prolonged or severe hypoglycemia is an important risk.

\medterm{Glycaemia} The amount of glucose in the blood. Insulin lowers blood glucose, while glucagon and other hormones raise it when needed. Abnormally high or low levels can cause symptoms and, if persistent, harm organs.

\medterm{Glycated Hemoglobin} A blood test showing the average blood glucose level over roughly the previous two to three months. It is used to diagnose diabetes and monitor treatment, but results may be misleading when red blood cell survival is abnormal.

\synonyms
Hemoglobin A1c; HbA1c; Glycohemoglobin

\medterm{Glycation} A chemical process in which sugar attaches to proteins, fats, or genetic material without the usual enzyme control. Excess glycation during prolonged high blood glucose can damage blood vessels, nerves, kidneys, eyes, and other tissues.

\medterm{Glycemia} The amount of glucose in the blood. It is affected by food, activity, illness, the liver, and hormones such as insulin and glucagon. Too little is hypoglycemia; too much is hyperglycemia.

\medterm{Glycemic Control} The level and pattern of blood-glucose management over time. It includes glucose values, symptoms, and the effects of food, activity, illness, medicines, pregnancy, and other circumstances.

\medterm{Glycemic Index} A ranking of carbohydrate-containing foods according to how quickly they raise blood glucose compared with a reference food. It does not account for the amount of carbohydrate eaten in a serving.

\medterm{Glycemic Index And Diabetes} The glycemic index ranks carbohydrate-containing foods by their effect on blood glucose levels. Low GI foods (\(\leq 55\)) cause gradual glucose rise; high GI foods (\(\geq 70\)) cause rapid spikes. Important for diabetes management and dietary planning.

\synonyms
GI; Glycemic Load Assessment

\medterm{Glycemic Load} A measure that combines a food’s glycemic index with the amount of available carbohydrate in a serving to estimate its effect on blood glucose.

\medterm{Glyceraldehyde} A small sugar with three carbon atoms that appears as the body breaks down glucose and other carbohydrates. It is quickly changed into other substances during the steps that release energy from food.

\medterm{Glyceride} A glyceride is a compound formed when glycerol is linked to one or more fatty acids, including many dietary fats.

\medterm{Glycerin} A colorless, odorless, sweet-tasting liquid also known as glycerol. It attracts water, dissolves some substances, and is used in medicines, skin products, foods, and selected treatments.

\medterm{Glycerol} A colorless, sweet-tasting alcohol that forms part of fats. It is used in medicines, foods, and skin products and can draw water into the bowel when used in some laxatives.
\synonyms
Glycerine; Trihydroxypropane

\medterm{Glycerol Kinase} Glycerol kinase is an enzyme that adds phosphate to glycerol, allowing it to enter pathways for energy production and fat metabolism.

\medterm{Glycerophosphates} Compounds containing glycerol and phosphate that help build cell membranes and participate in energy metabolism, chemical signaling, and other cellular processes.

\medterm{Glycerophospholipids} The major membrane phospholipids built from glycerol, fatty acids, phosphate, and a polar head group, including phosphatidylcholine and phosphatidylethanolamine.

\medterm{Glycerylphosphorylcholine} A choline-containing substance formed when certain membrane fats are broken down. It occurs in body fluids and cells and participates in membrane metabolism and the supply of choline.

\medterm{Glyceryl Trinitrate} A nitrate medicine, also called nitroglycerin, that releases nitric oxide and relaxes blood vessels. It can relieve or prevent angina, but may cause headache, dizziness, or low blood pressure and can interact dangerously with phosphodiesterase-5 inhibitors.

\medterm{Glycine} A small amino acid used to make proteins and heme. It also acts as an inhibitory chemical messenger, especially in the spinal cord and brainstem, helping regulate movement and reflexes.

\medterm{Glycine Receptor} A ligand-gated chloride channel activated by glycine that mediates inhibitory neurotransmission, especially in the spinal cord and brainstem.

\medterm{Glycocalyx} A sugar-rich coating on the outside of many cells, including cells lining blood vessels. It helps protect surfaces, control passage of substances, support cell attachment, and influence inflammation and clotting.

\medterm{Glycogen} The body’s stored form of glucose. It is kept mainly in the liver, which releases it between meals, and in muscles, which use it during activity.

\medterm{Glycogenesis} The process of converting glucose into glycogen for storage, mainly in the liver and muscles. It occurs after eating, when blood glucose and insulin levels rise.

\medterm{Glycogenolysis} The breakdown of stored glycogen. The liver uses it to release glucose into the blood between meals, while muscles use their glycogen locally for energy during exercise.

\medterm{Glycogen Storage Disease} Any of several inherited disorders in which the body cannot make or break down glycogen normally. Depending on the type, these disorders may cause low blood glucose, an enlarged liver, muscle weakness, exercise intolerance, or heart problems.

\medterm{Glycogen Storage Diseases} Inherited disorders that prevent cells from making or breaking down glycogen normally. They can cause low blood sugar, enlarged liver, muscle weakness, or heart disease.

\medterm{Glycohemoglobin} Hemoglobin with glucose chemically attached to it; glycated hemoglobin A1c reflects average blood glucose over roughly the preceding two to three months.

\medterm{Glycolipid} A lipid containing a carbohydrate group that contributes to cell membranes, recognition, signaling, and myelin structure.

\medterm{Glycolysis} A series of reactions in the cytosol that breaks one glucose molecule into two pyruvate molecules. It produces ATP and other molecules used in further energy production and can proceed when oxygen is limited.

\medterm{Glycopeptide} A molecule made of a short protein chain with one or more attached sugar groups. Glycopeptides occur in cell-surface proteins, antibodies, and some medicines, and the attached sugars can affect structure and function.

\medterm{Glycophorin} A sugar-coated protein in the red-blood-cell membrane that helps maintain the cell surface and carries markers used to distinguish blood-cell types.

\medterm{Glycoprotein} A protein with covalently attached carbohydrate chains, important in receptors, antibodies, enzymes, mucus, cell adhesion, and extracellular matrix.

\medterm{Glycopyrrolate} Glycopyrrolate is an anticholinergic medicine used to reduce secretions and, in some inhaled forms, to help control chronic obstructive lung disease; it may cause dry mouth, constipation, or urinary retention.

\medterm{Glycosaminoglycan} A long, unbranched carbohydrate chain found in the extracellular matrix and connective tissues. Abnormal breakdown or accumulation of GAGs occurs in mucopolysaccharidoses.

\medterm{Glycosaminoglycans} Glycosaminoglycans are long chains of sugars found in connective tissue and involved in its structure, hydration, and signaling.

\medterm{Glycoside} A compound in which a sugar is chemically linked to another molecule. Glycosides occur in plants, foods, and medicines; their effects depend on the attached molecule and how the body breaks the bond.

\medterm{Glycosphingolipid} A membrane lipid containing ceramide and one or more sugars, involved in cell recognition, signaling, and nervous-system structure.

\medterm{Glycosuria} The presence of glucose in the urine. It most often occurs when blood glucose is high, but it can also result from pregnancy, certain medicines, or a kidney-tubule problem that prevents normal glucose reabsorption.

\medterm{Glycosylated Hemoglobin} Hemoglobin with glucose attached; the A1C measurement reflects average blood sugar over roughly two to three months and helps diagnose or monitor diabetes.

\synonyms
Hemoglobin A1c; HbA1c; Glycohemoglobin

\medterm{Glycosylation} Attachment of sugar chains to proteins or fats inside cells. This process helps determine their shape, stability, location, and activity; abnormal glycosylation can contribute to inherited disease, cancer, and immune disorders.

\medterm{Glycosylphosphatidylinositol} A lipid-and-sugar anchor that attaches certain proteins to the outer surface of a cell membrane. Defects in making or attaching this anchor can cause rare blood, nerve, or developmental disorders.

\medterm{Glycosyltransferase} Glycosyltransferase is an enzyme that transfers a sugar to a protein, lipid, or other molecule during glycan synthesis.

\medterm{Glycyrrhetinic Acid} An active metabolite of licorice glycosides that inhibits 11-beta-hydroxysteroid dehydrogenase; excess exposure can cause hypertension, low potassium, and fluid retention.

\medterm{Glycyrrhizic Acid} The major saponin in licorice root that is converted to glycyrrhetinic acid and can produce mineralocorticoid-like toxicity when consumed excessively.

\medterm{Glyphosate} A broad-spectrum herbicide that inhibits plant EPSP synthase; concentrated exposure can irritate tissues and cause systemic poisoning, especially after ingestion.

\medterm{Gm} Abbreviation; see \textit{Gram}.

\medterm{GN} Abbreviation; see \textit{Glomerulonephritis}.

\medterm{Gnathic} Relating to the jaw, especially the upper jaw or lower jaw and their development, structure, or position.

\medterm{Gnathion} The lowest midline point on the outer surface of the lower jaw. It is used as a landmark in dental, orthodontic, facial, and skull measurements.

\medterm{Gnathitis} Inflammation of the jaw or jaw tissues. It may cause pain, swelling, tenderness, difficulty opening the mouth, or chewing problems and can result from infection, injury, or inflammatory disease.

\medterm{Gnathoplasty} Plastic or reconstructive surgery to reshape or repair the jaw. It may improve function, facial balance, or a structural defect; risks include infection, nerve injury, altered bite, and recurrence.

\medterm{Gnathostoma} A group of roundworms that can infect people who eat raw or undercooked freshwater fish, frogs, or other carriers. Larvae may migrate through skin, organs, or the nervous system and cause swelling, pain, or serious neurologic disease.

\medterm{Gnathostoma Spinigerum} A roundworm that can cause gnathostomiasis after people eat raw or undercooked freshwater animals. Migrating larvae may produce recurring painful skin swellings and, rarely, eye or brain disease.

\medterm{Gnathostomiasis} A parasitic infection caused by Gnathostoma larvae acquired from undercooked freshwater fish or other hosts, producing migratory skin swellings or deeper neurologic or ocular disease.

\medterm{Gnawing Pain} A dull, persistent pain often described as an ache or a feeling that something is eating away at the body. It can arise from muscles, nerves, internal organs, or inflammation. Its location, duration, triggers, and accompanying symptoms help identify the cause.

\medterm{GNE Myopathy} A rare inherited muscle disease that usually begins in early adulthood and slowly weakens the muscles of the lower legs, causing foot drop, then the hands and arms. The thigh muscles are often spared until late disease. There is no cure, but rehabilitation helps maintain function.

\synonyms
Nonaka Myopathy; Distal Myopathy with Rimmed Vacuoles (DMRV); Hereditary Inclusion Body Myopathy (HIBM)

\medterm{GnRH} Gonadotropin-releasing hormone, a hormone from the brain that signals the pituitary gland to release hormones controlling ovulation, sperm production, and sex-hormone production.

\synonyms
Gonadotropin-Releasing Hormone

\medterm{GnRH Agonist} A medicine that first increases and then suppresses signals from the pituitary gland to the reproductive organs. It lowers sex-hormone production and is used in selected fertility treatments, endometriosis, fibroids, and hormone-sensitive cancers.

\medterm{GnRH Antagonist} A medicine that directly blocks the brain-to-pituitary signal controlling reproductive hormones. It lowers these hormones quickly without the temporary hormone rise that can occur when a GnRH agonist is started.

\medterm{GnRH Antagonists} Medicines that block a brain hormone signal and quickly reduce pituitary release of reproductive hormones. They are used in some fertility treatments and hormone-sensitive cancers, with effects depending on the specific medicine.

\medterm{Goal Attainment Scaling} A way to measure progress toward personal healthcare goals. The goal is defined in advance with levels ranging from much less than expected to much better than expected, then scored after treatment.

\medterm{Goal Setting} Choosing clear, realistic outcomes and deciding how and when to work toward them. In healthcare, goals can reflect a person’s priorities, include a way to measure progress, and change as needs or circumstances change.

\medterm{Goals-Of-Care Discussion} A clinical conversation that aligns treatment recommendations with a patient's values,
priorities, prognosis, acceptable burdens, and preferred outcomes. It may include
decisions about hospitalization, resuscitation, artificial nutrition, and
comfort-focused care.

\medterm{Goblet Cell} A mucus-secreting epithelial cell, abundant in the respiratory and intestinal tracts, that protects surfaces with a mucin-rich layer.

\medterm{Goblet Cell Adenocarcinoma} A distinctive appendiceal adenocarcinoma composed of goblet-like mucinous cells, formerly called goblet cell carcinoid. It can spread within the peritoneum or to the ovaries and is managed according to grade, stage, and dissemination.

\medterm{Goblet Cells} Cells lining the intestines and airways that make and release mucus. Mucus moistens and protects these surfaces and helps trap particles, germs, and other substances for removal.

\medterm{Goeckerman Regimen} A dermatological therapy used for extensive or recalcitrant plaque psoriasis combining the daily topical application of crude coal tar with exposure to ultraviolet B (UVB) phototherapy.
\synonyms
Goeckerman Therapy; Coal Tar and UVB Regimen

\medterm{Goiter} Enlargement of the thyroid gland, either throughout the gland or as one or more lumps. It may occur with normal, overactive, or underactive thyroid function and can result from iodine deficiency, immune disease, or a tumor.

\synonyms
Thyromegaly; Struma; Thyroid Enlargement

\medterm{Goitre} Enlargement of the thyroid gland. It may occur with normal, low, or high thyroid hormone levels and can result from iodine deficiency, autoimmune disease, nodules, or inflammation.

\medterm{Goitrogenic Foods} Foods or food components that can interfere with thyroid hormone production or iodine use under some conditions; risk depends on intake, preparation, iodine status, and thyroid health.

\medterm{Gold} A metal used in selected dental restorations, implants, laboratory applications, and historical medicinal preparations; exposure can cause contact allergy or toxicity depending on the compound and dose.

\medterm{Gold Alloy} A mixture of gold and other metals used for some dental restorations. It is durable, fits closely, and is usually well tolerated by body tissues, but it costs more and has a distinctive appearance.

\medterm{Gold Alloys} Metal mixtures containing gold and other elements, used in selected dental restorations or medical devices; properties and biocompatibility depend on composition.

\medterm{Goldenhar Syndrome} A birth condition, also called oculo-auriculo-vertebral spectrum, in which one side of the face, ear, eye, or spine develops differently. Features may include facial asymmetry, a small or misshapen ear, hearing loss, eyelid or eye-surface growths, and spinal differences.

\medterm{Golden Hour} The golden hour is the early period after a serious injury or medical emergency when rapid assessment and treatment may improve the chance of survival and recovery.

\medterm{Goldmann Tonometry} A slit-lamp examination that estimates intraocular pressure by measuring the force needed to flatten a defined area of the cornea.

\medterm{Gold Therapy} Treatment with gold-containing medicines, formerly used mainly for rheumatoid arthritis. Its use is now limited because benefit is slow and serious effects can include skin or mouth inflammation, kidney injury, and reduced blood-cell counts.

\medterm{Golfer Elbow} An overuse injury of the tendons on the inner side of the elbow. Repeated gripping or wrist movement can cause pain and tenderness, usually worsened by lifting, bending the wrist, or making a fist.

\medterm{Golfer's Elbow} Painful irritation or degeneration of tendons attached to the inner side of the elbow. Repeated gripping or wrist movement can cause tenderness, reduced strength, and pain during lifting or bending.

\medterm{Golgi Apparatus} A structure inside cells that modifies, sorts, and packages proteins and fats made by the cell, sending them to other parts of the cell or releasing them outside it.

\medterm{Golgi Cisternae} Flattened membrane-bound sacs that form the Golgi apparatus and modify, sort, and package proteins and lipids for transport within or outside the cell.

\medterm{Golgi Tendon Organ} A sensory structure where a muscle joins its tendon. It detects tension in the muscle and tendon and helps the nervous system control muscle force and protect against excessive strain.

\medterm{Gomori's Stain} Gomori's stain is a group of laboratory stains used to highlight structures such as connective tissue, carbohydrates, fungi, or muscle components in tissue samples.

\medterm{Gomphosis} A fixed joint in which a tooth is held in its socket by a strong ligament.

\medterm{Gonad} A testis or ovary, an organ producing reproductive cells and sex hormones. Gonadal function influences puberty, fertility, menstrual cycles, sexual development, bone health, and other body processes.

\medterm{Gonadal Agenesis} Congenital absence or severe underdevelopment of a gonad, resulting in impaired sex-hormone production and infertility or differences in sexual development.

\medterm{Gonadal Dysgenesis} Incomplete or abnormal development of the ovaries or testes, often resulting from
chromosomal abnormalities, gene variants, or disordered sex development. It may cause
atypical genital development, delayed or absent puberty, infertility, and abnormal
gonadal hormone production.

\medterm{Gonadal Vein} A vein draining a testis or ovary. The right gonadal vein usually drains directly into the inferior vena cava, while the left usually drains into the left renal vein.

\medterm{Gonad Disorder} A disease affecting testes or ovaries and causing abnormal reproductive-cell production, hormone secretion, development, fertility, or sexual function.

\medterm{Gonadoblastoma} A rare tumor that develops in an abnormal or incompletely developed ovary or testicle. It may produce sex hormones and has a risk of becoming a malignant germ-cell tumor, especially when a Y-chromosome change is present.

\medterm{Gonadorelin} Synthetic gonadotropin-releasing hormone that stimulates pituitary release of luteinizing hormone and follicle-stimulating hormone when given intermittently.

\medterm{Gonadotrophin} A hormone that acts on the ovaries or testes. The pituitary hormones follicle-stimulating hormone and luteinizing hormone regulate egg development, ovulation, sperm production, and sex-hormone release.

\medterm{Gonadotrophin-Releasing Factors} Signals from the brain that control release of the pituitary hormones acting on the ovaries or testes. The main regulator is gonadotropin-releasing hormone, released in pulses from the hypothalamus.

\medterm{Gonadotrophins} Hormones that act on the ovaries or testes. Follicle-stimulating hormone helps develop eggs and sperm, while luteinizing hormone triggers ovulation and supports sex-hormone production.

\medterm{Gonadotrophs} Pituitary cells that release luteinizing hormone and follicle-stimulating hormone after stimulation by GnRH, regulating the ovaries, testes, fertility, and sex-hormone production.

\medterm{Gonadotropin} A hormone, chiefly luteinizing hormone or follicle-stimulating hormone, that acts on the ovaries or testes to regulate gamete production and sex-hormone secretion.

\medterm{Gonadotropin Hormone} A hormone that acts on the ovaries or testes. The pituitary gonadotropins follicle-stimulating hormone and luteinizing hormone regulate reproductive development, ovulation, sperm production, and sex-hormone production.

\medterm{Gonadotropin-Releasing Hormone} A hormone made in the brain that signals the pituitary gland to release two hormones, FSH and LH. These hormones control the ovaries and testes, including egg development, ovulation, sperm production, and sex-hormone production.

\synonyms
GnRH

\medterm{Gonadotropins} Hormones that act on the gonads. In humans, the pituitary gonadotropins are follicle-stimulating hormone and luteinizing hormone, which regulate gamete production and sex-hormone secretion.

\medterm{Gonads} The reproductive organs: testes in males and ovaries in females. They produce sperm or eggs and make sex hormones such as testosterone, estrogen, and progesterone. These hormones influence puberty, fertility, menstrual cycles, pregnancy, and bone health.

\medterm{Gondoic Acid} A 20-carbon monounsaturated fatty acid found in some seed oils and incorporated into dietary and tissue lipids.

\medterm{Goniometer} An instrument measuring joint angles and range of motion, used to assess stiffness, document impairment, and track progress during rehabilitation.

\medterm{Gonioscopy} An eye examination using a special contact lens to view the angle where the iris meets the cornea. It helps determine whether glaucoma is open-angle or angle-closure and guides treatment.

\medterm{Gonococcal} Relating to Neisseria gonorrhoeae, the bacterium that causes gonorrhea. Gonococcal infection commonly affects the genitals, rectum, or throat and may spread to joints or blood.

\medterm{Gonococcal Arthritis} A disseminated infection of joints caused by Neisseria gonorrhoeae, often with fever, migrating joint or tendon pain, skin pustules, and inflammation around tendons.

\medterm{Gonococcal Pharyngitis} Infection of the throat by Neisseria gonorrhoeae, usually acquired through sexual contact. It may
cause sore throat or no symptoms.

\medterm{Gonococcal Urethritis} Infection and inflammation of the urethra caused by Neisseria gonorrhoeae. It may cause burning urination or discharge, but some people have no symptoms; untreated infection can spread.

\medterm{Gonococcus} A spherical bacterium of the genus Neisseria. Neisseria gonorrhoeae, the gonococcus, causes the sexually transmitted infection gonorrhea.

\medterm{Gonorrhea} A sexually transmitted infection caused by the bacterium Neisseria gonorrhoeae. It may cause discharge, painful urination, pelvic or rectal symptoms, or no symptoms. Untreated infection can lead to infertility, joint infection, or infection in a newborn.

\medterm{Gonorrhoea} Alternate spelling; see \textit{Gonorrhea}.

\medterm{Goodell Sign} A clinical sign of early pregnancy characterized by marked softening of the uterine cervix (produced by increased pelvic vascularity and edema), identifiable upon digital pelvic bimanual examination around the sixth to eighth week of gestation.

\medterm{Goodpasture's Syndrome} An autoimmune disease in which antibodies attack the lungs and kidney filters, causing coughing blood, kidney inflammation, blood in urine, and sometimes rapid kidney failure.

\medterm{Goodpasture Syndrome} An autoimmune disease in which antibodies attack the lung air sacs and kidney filters. It can cause coughing blood, breathlessness, blood in the urine, and rapidly worsening kidney failure.

\medterm{Good Samaritan} A person who voluntarily helps someone in an emergency, especially by providing immediate aid before professional responders arrive.

\medterm{Good Samaritan Law} A law that may protect a person who provides reasonable emergency assistance from certain civil claims, subject to jurisdiction-specific conditions and exceptions such as gross negligence.

\medterm{Goose Bump} A small raised area around a hair follicle caused by contraction of tiny arrector pili muscles, commonly in response to cold, fear, or strong emotion.

\medterm{Gooseflesh} Temporary raised bumps around hair follicles caused by contraction of tiny skin muscles, often in response to cold, fear, or strong emotion. It is usually harmless.

\medterm{Goose-Skin} Goose skin is the temporary raising of hair follicles, usually caused by cold, fear, or strong emotion.

\medterm{Gorlin's Syndrome} An inherited disorder that increases the risk of multiple basal-cell skin cancers, jaw cysts, and abnormalities of the bones, eyes, and nervous system.

\medterm{Gorlin Syndrome} Alternate name; see \textit{Gorlin's Syndrome}.

\medterm{Goserelin} A gonadotropin-releasing hormone agonist that initially stimulates and then suppresses pituitary gonadotropins; it is used in hormone-sensitive prostate or breast cancer and selected gynecologic conditions.

\medterm{Gottron Papules} Raised red, purple, or scaly patches over the knuckles and other finger joints. They are a characteristic skin finding of dermatomyositis, an inflammatory disease that can also cause muscle weakness.

\medterm{Gottron Sign} Symmetrical, violaceous, macular or papular erythematous patches and scaling located over the extensor aspects of the joints (particularly the elbows, patellae, and medial malleoli), serving as a characteristic cutaneous hallmark of dermatomyositis.
\synonyms
Gottron's Sign; Dermatomyositis Extensor Erythema

\medterm{Gougerot-Carteaud Syndrome} A skin disorder causing brown, slightly raised spots that join into net-like patches, usually on the upper chest, back, or neck of young people.

\medterm{Gout} A type of inflammatory arthritis caused by a buildup of urate, a substance made when the body breaks down purines. When urate levels stay high, needle-shaped crystals can form in and around joints. These crystals cause sudden attacks, or flares, of severe pain, swelling, warmth, redness, and tenderness, often in the big toe but also in other joints. Repeated flares can damage joints. Long-standing gout may cause firm crystal deposits called tophi and uric acid kidney stones.

\synonyms
Gouty Arthritis; Monosodium Urate Arthropathy

\medterm{Gout Flare} A sudden attack of gout caused by inflammation around urate crystals in a joint. The joint becomes very painful, swollen, warm, and red, often beginning in the big toe. Infection can cause similar symptoms.

\medterm{Gout Nodulosis} A condition in which many firm urate deposits, called tophi, form under the skin or around joints because of long-standing gout. They may occur at the ears, fingers, toes, elbows, or Achilles tendons and can restrict movement or break through the skin.

\medterm{Gout Suppressants} Medicines that lower uric-acid production or increase its excretion to prevent recurrent gout attacks and urate crystal deposition.

\medterm{Gouty Tophus} A nodular, chalky subcutaneous or periarticular deposit of crystallized monosodium urate monohydrate crystals surrounded by chronic foreign-body granulomatous inflammation, occurring in chronic uncontrolled hyperuricemia on the helix of the ear, olecranon bursa, Achilles tendon, and distal joints.

\synonyms
Tophus; Tophi

\medterm{Gowers Sign} A maneuver in which a person uses the hands to push up from the floor or thighs because of proximal muscle weakness, classically seen in muscular dystrophy.

\medterm{Gower Syndrome} A form of reflex syncope, also called vasovagal or vasodepressor syncope, in which a trigger causes a fall in blood pressure and sometimes heart rate, briefly reducing brain perfusion.

\medterm{G-Protein Coupled Receptors} Proteins on the cell surface that receive signals from hormones, medicines, and other substances and pass those signals into the cell. They help control heart rate, vision, digestion, mood, and inflammation.

\medterm{Graafian Follicle} A mature, fluid-filled sac in an ovary that contains an egg ready for release. It enlarges during the menstrual cycle and normally opens at ovulation; after release, the remaining tissue forms the corpus luteum.

\medterm{GRACE Score} A risk score used in people with an acute coronary syndrome to estimate the chance of death or another heart attack during the hospital stay and over the next six months. It uses age, vital signs, kidney function, cardiac arrest, ECG changes, heart-failure signs, and heart-injury blood tests to guide treatment intensity.

\synonyms
Global Registry of Acute Coronary Events Score; GRACE Risk Model

\medterm{Gracilis Muscle} A long, slender strap-like muscle located on the medial aspect of the thigh originating from the inferior pubic ramus and body of the pubis, descending to insert into the pes anserinus on the proximal medial tibia; acts to adduct the thigh and flex and internally rotate the leg.

\medterm{Grade} A measure of how abnormal cells look under a microscope and, in many cancers, how quickly the growth is likely to behave. Lower grades usually look more like normal cells; higher grades usually look more abnormal. The grading system differs by disease.

\medterm{Graded Potentials} Small changes in an electrically active cell caused by a stimulus. Their size depends on the stimulus strength, and they can either increase or decrease the chance that the cell will send a nerve signal.

\medterm{Grade IV Astrocytoma} An older term for glioblastoma, an aggressive brain tumor that grows into nearby brain tissue. It may cause headache, seizures, weakness, speech or vision problems, or changes in thinking. Treatment often combines surgery, radiation, and medicines.

\synonyms
Glioblastoma; GBM

\medterm{Gradenigo Syndrome} A classic clinical triad consisting of suppurative otitis media (otorrhea), retro-orbital or temporal pain in the sensory distribution of the trigeminal nerve (cranial nerve V), and ipsilateral abducens nerve palsy (cranial nerve VI) caused by petrous apicitis.
\synonyms
Gradenigo's Syndrome; Petrous Apicitis Triad

\medterm{Graefe's Sign} A delayed downward movement of the upper eyelid during downward gaze, classically associated with thyroid eye disease but not specific to it.

\medterm{Graft} Tissue or a synthetic material transferred to another site to replace, repair, reinforce, or bypass a damaged body structure.

\medterm{Grafting} The surgical transfer of tissue or a synthetic substitute to repair, replace, or support a body structure. An autograft comes from the same person, whereas an allograft comes from another person of the same species.

\medterm{Graft Versus Host} Alternate name; see \textit{Graft Versus Host Disease}.

\medterm{Graft Versus Host Disease} A complication of a stem-cell or bone-marrow transplant in which immune cells from the donor attack the recipient’s tissues. It commonly affects the skin, liver, and digestive tract and may be short-term or long-lasting; prevention and treatment use immune-suppressing medicines.
\synonyms
GVHD; Transfusion-Associated or Allogeneic GVHD

\medterm{Graft-Versus-Host Disease} Alternate spelling; see \textit{Graft Versus Host Disease}.

\medterm{Graft-Versus-Host Reaction} An immune reaction in which donor immune cells attack the recipient's tissues after transplantation, potentially affecting skin, intestines, liver, lungs, or other organs.

\medterm{Gram} A metric unit used to measure mass, equal to one-thousandth of a kilogram. Grams are commonly used for medicines, laboratory samples, food, and body substances; one thousand milligrams make one gram.

\synonyms
Metric Gram; Unit G

\medterm{Gramicidin} Gramicidin is a topical antibiotic that forms channels in bacterial membranes; it is not generally used systemically because it can be toxic.

\medterm{Gram Mass Unit} A metric unit of mass equal to one-thousandth of a kilogram; medication doses and laboratory measurements are commonly expressed in grams or smaller metric units.

\medterm{Gram Metric Unit} A metric unit of mass equal to one-thousandth of a kilogram or approximately 15.432 grains.

\medterm{Gram-Negative} Describing bacteria with a thin peptidoglycan cell wall and an outer membrane containing lipopolysaccharide; they stain pink or red with the Gram stain.

\medterm{Gram-Negative Bacilli} Rod-shaped bacteria that appear pink or red with a Gram stain because they have a thin peptidoglycan layer and an outer membrane. They include harmless body organisms and bacteria that can cause urinary, respiratory, intestinal, bloodstream, wound, or other infections.

\medterm{Gram-Negative Bacteria} Bacteria that do not retain crystal violet during Gram staining and therefore appear pink or red.
They have a thin peptidoglycan layer and an outer membrane that may contain
lipopolysaccharide. Examples include Escherichia coli, Klebsiella, Pseudomonas, and
Neisseria.

\medterm{Gram-Negative Cocci And Coccobacilli} Gram-negative bacteria that are either round or short rod-shaped. Cocci are round, while coccobacilli are short rods that may look almost round. They include bacteria that can cause respiratory, genital, bloodstream, intestinal, and other infections.

\medterm{Gram-Negative Meningitis} Meningitis caused by gram-negative bacteria, including Escherichia coli, Klebsiella, and Neisseria species. It can progress rapidly to sepsis, neurologic injury, and death.

\medterm{Gram-Positive} Describing bacteria with a thick peptidoglycan cell wall that retains crystal violet during Gram staining and therefore appears purple.

\medterm{Gram-Positive Bacilli} Rod-shaped bacteria that retain the purple part of the Gram stain because of their thick peptidoglycan cell wall. Some are harmless colonizers, while others can cause wound, intestinal, respiratory, bloodstream, or nervous-system infections.

\medterm{Gram-Positive Bacteria} Bacteria that retain crystal violet in the Gram stain (appear purple/blue),
characterized by a thick peptidoglycan cell wall and no outer membrane.

\medterm{Gram-Positive Cocci} Round bacteria that appear purple with a Gram stain because they have a thick peptidoglycan cell wall. Important groups include staphylococci, streptococci, and enterococci, which may normally live on the body or cause skin, throat, lung, urinary, heart-valve, or bloodstream infections.

\medterm{Gram's Stain} Gram's stain is a laboratory method that classifies bacteria by how their cell walls retain the stain, helping guide identification and initial antibiotic choices.

\medterm{Gram Stain} A laboratory test that colors bacteria to group them as Gram-positive or Gram-negative. It helps identify possible causes of infection quickly and can guide initial treatment while more specific tests are performed.

\medterm{Gram Stain Of Skin Lesion} Microscopic examination of a skin lesion specimen stained with Gram stain to identify bacteria. Useful for diagnosing bacterial skin infections and guiding antibiotic therapy. Results available within hours.

\medterm{Gram Stain Of Tissue Biopsy} Microscopic examination of tissue biopsy specimen stained with Gram stain to identify bacteria. Used to diagnose tissue infections and guide antimicrobial therapy. May be performed alongside culture and histopathology.

\synonyms
Tissue Gram Stain; Histochemical Bacterial Stain

\medterm{Grand Mal} Alternate historical term; see \textit{Tonic-Clonic Seizure}.

\medterm{Grand Mal Seizure} Alternate name; see \textit{Tonic-clonic seizure}.

\medterm{Grand Multipara} A person who has given birth five or more times beyond the gestational threshold used to count parity; definitions can vary by obstetric system.

\medterm{Grand Rounds} An educational medical conference in which clinicians present a case, diagnosis, management, or clinical topic to a professional audience.

\medterm{Granisetron} Granisetron is a selective 5-HT3-receptor antagonist used to prevent chemotherapy-, radiotherapy-, or surgery-related nausea and vomiting; headache, constipation, and QT prolongation can occur.

\medterm{Granisetron Hydrochloride} The hydrochloride form of granisetron, a medicine that blocks serotonin 5-HT3 receptors to prevent nausea and vomiting from chemotherapy, radiation, or surgery. Headache, constipation, and heart-rhythm changes may occur.

\medterm{Granular Casts} Small grainy particles formed inside kidney tubules and seen in urine. They may occur when kidney cells are injured or break down and can be found in acute kidney injury or some chronic kidney diseases.

\medterm{Granular Cell Tumor Of The Breast} An uncommon tumor showing Schwann-cell differentiation that can occur in the breast and mimic carcinoma clinically or on imaging. Histologic examination distinguishes it from malignancy; most are benign.

\medterm{Granular Leukocyte} A granulocyte, a white blood cell containing cytoplasmic granules; the main types are neutrophils, eosinophils, and basophils.

\medterm{Granular Parakeratosis} A rare skin condition causing brown or reddish, sometimes itchy patches or papules, often in skin folds such as the armpits. It reflects abnormal keratinization and may be difficult to distinguish from other rashes.

\medterm{Granulation} The formation of new vascular connective tissue during wound healing, producing moist red or pink granulation tissue made of capillaries, fibroblasts, and extracellular matrix.

\medterm{Granulation Tissue} New tissue that grows in a healing wound and contains small blood vessels, connective tissue, and immune cells. Healthy granulation tissue is usually pink or red and moist.

\medterm{Granulicatella} Granulicatella is a group of bacteria that may live in the mouth and can rarely cause bloodstream infection or endocarditis.

\medterm{Granulicatella Adiacens} A mouth bacterium that can rarely cause serious infection, especially infection of heart valves or the bloodstream. It may be difficult to grow in culture, and treatment depends on susceptibility testing.

\medterm{Granulicatella Elegans} A bacterium normally found in the mouth that can rarely cause bloodstream infection or infection of a heart valve. It may grow poorly in routine culture, so repeated positive cultures and compatible symptoms are important when judging its significance.

\medterm{Granulocyte} A type of white blood cell containing visible granules. Neutrophils, eosinophils, and basophils are granulocytes that help fight infection, trigger inflammation, and respond to allergic or parasitic disease.

\medterm{Granulocyte Colony-Stimulating Factor} A hematopoietic glycoprotein growth factor that stimulates the bone marrow to produce and differentiate granulocyte precursors into mature neutrophils and mobilize them into circulating blood; recombinant forms (filgrastim, pegfilgrastim) are used to treat and prevent chemotherapy-induced neutropenia.

\synonyms
G-CSF

\medterm{Granulocyte Count} A granulocyte count measures the number of neutrophils, eosinophils, and basophils in the blood and helps assess infection, inflammation, and bone-marrow function.

\medterm{Granulocytic Sarcoma} A mass of immature myeloid blood-forming cells outside the bone marrow, also called myeloid sarcoma. It may occur with acute myeloid leukemia or precede it.

\medterm{Granulocytopenia} An abnormally low number of granulocytes, especially neutrophils, which can increase susceptibility to bacterial and fungal infection.

\medterm{Granuloma} A small area of organized inflammation formed when immune cells surround material they cannot easily remove, such as certain germs, foreign material, or substances released by the body.

\medterm{Granuloma Annulare} A benign, noninfectious inflammatory skin condition that causes raised red, pink, or skin-colored bumps arranged in rings, often with clearer skin in the center. It commonly affects the hands, feet, or elbows and may cause little or no discomfort. Widespread disease may be associated with diabetes or thyroid disease, but the cause is often unknown.

\medterm{Granuloma Inguinale} A chronic sexually transmitted infection causing painless, slowly enlarging ulcers in the genital or groin region. It is caused by bacteria and spreads through sexual contact; antibiotics are used in treatment.

\medterm{Granulomatosis} A disease in which long-lasting inflammation forms granulomas, small organized collections of immune cells. Depending on the cause, granulomas may affect the lungs, skin, lymph nodes, intestines, or other organs.

\medterm{Granulomatosis With Polyangiitis} An autoimmune disease that inflames and damages small and medium-sized blood vessels, often affecting the nose, sinuses, lungs, and kidneys. It may cause sinus problems, cough, coughing blood, or kidney injury.
\synonyms
GPA; Wegener Granulomatosis

\medterm{Granulomatous Amebic Encephalitis} A rare, severe brain infection caused by free-living amoebae, usually in people with weakened immunity. It develops over days or weeks and may cause headache, confusion, seizures, weakness, or coma.

\medterm{Granulomatous Colitis} Inflammation of the colon in which granulomas are found on tissue examination, often associated with Crohn disease but also possible with infections or other inflammatory conditions.

\medterm{Granulomatous Inflammation} A pattern of long-lasting inflammation in which immune cells form small clusters around material they cannot easily remove. It can occur with infections, foreign materials, autoimmune disease, or other inflammatory conditions.

\medterm{Granulomatous Mastitis} A chronic inflammatory breast disease characterized by noncaseating granulomas after other causes, especially infection and sarcoidosis, have been excluded. It may present as a painful mass, abscesses, or draining sinuses.

\medterm{Granulosa Cell} A granulosa cell is a cell surrounding an ovarian egg that supports its development and produces hormones such as estrogen and inhibin.

\medterm{Granulosa Cell Tumor} A tumor arising from hormone-producing cells in the ovary. It may release estrogen, causing irregular bleeding or early puberty, and can recur years after treatment even when it initially appears slow-growing.

\medterm{Granulosa Tumor} Alternate form; see \textit{Granulosa Cell Tumor}.

\medterm{Granzyme} A protease released by cytotoxic T lymphocytes and natural killer cells that enters target cells through perforin-associated pores and promotes programmed cell death.

\medterm{Grasp Reflex} An automatic closing of a baby’s fingers when the palm is touched. It is normal in early infancy and usually fades as voluntary hand control develops; persistence or asymmetry can occur with nervous-system problems.

\medterm{Grass Allergy} Hypersensitivity reaction to grass pollen, causing allergic rhinitis (hay fever) with symptoms including sneezing, nasal congestion, itchy eyes, and throat irritation. Seasonal, typically spring through fall. Treatment includes antihistamines, nasal corticosteroids, and allergen avoidance.

\synonyms
Grass Pollen Allergy; Gramineae Pollinosis

\medterm{Grass And Weed Killer Poisoning} Illness caused by swallowing, inhaling, or getting herbicide on the skin or eyes. Effects range from irritation, vomiting, and breathing problems to severe poisoning, depending on the product and amount.

\synonyms
Herbicide Intoxication; Weedicide Poisoning

\medterm{Graves Disease} An autoimmune disease that causes an overactive thyroid (hyperthyroidism). The immune system makes antibodies that act like thyroid-stimulating hormone and make the thyroid produce too much thyroid hormone. Symptoms may include weight loss, a fast or irregular heartbeat, sweating, heat intolerance, tremor, nervousness, trouble sleeping, frequent bowel movements, muscle weakness, or an enlarged thyroid (goiter). Some people develop Graves eye disease, which can cause dry or gritty eyes, puffiness, bulging eyes, light sensitivity, or double vision. Rarely, the skin over the shins becomes thickened and red.

\medterm{Graves Ophthalmopathy} An autoimmune inflammatory disorder of the retro-orbital tissues and extraocular muscles associated with Graves disease, driven by autoantibodies against the TSH receptor leading to fibroblastic proliferation, glycosaminoglycan deposition, proptosis, upper lid retraction, and diplopia.
\synonyms
Graves Orbitopathy; Thyroid Eye Disease

\medterm{Grave Wax} Alternate name; see \textit{Adipocere}.

\medterm{Gravida} The number of times a person has been pregnant, regardless of pregnancy outcome; the current pregnancy is included in the count.

\medterm{Gravidity} Gravidity is the number of times a person has been pregnant, including the current pregnancy and pregnancies that did not result in a live birth.

\medterm{Gray} A unit used to measure the amount of ionizing radiation absorbed by a material or body tissue. One gray equals one joule of radiation energy absorbed per kilogram. It is used in radiation treatment and safety assessment.

\synonyms
Absorbed Dose Unit Gy; Radiation Unit Gray

\medterm{Gray Baby Syndrome} A serious form of chloramphenicol toxicity in newborns, caused by immature drug metabolism and characterized by abdominal distension, vomiting, cyanosis or gray skin color, cardiovascular collapse, and potentially death.

\medterm{Gray Matter} Areas of the brain and spinal cord containing many nerve-cell bodies and connections. Gray matter processes information and helps control movement, sensation, memory, speech, and other functions.

\medterm{Greater Curvature} The long, rounded outer border of the stomach, extending from its upper opening to its outlet. It provides attachment for the greater omentum and receives blood from several stomach arteries.

\medterm{Greater Omentum} A fold of fatty peritoneal tissue that hangs from the stomach and partly covers the abdominal organs.

\medterm{Greater Trochanter} The large bony prominence on the outer upper thigh, where several hip muscles attach. It can be felt through the skin and is a common site of pain from tendon problems, bursitis, or injury.

\medterm{Greater Trochanteric Pain Syndrome} Pain over the lateral hip centered on the greater trochanter, commonly caused by tendinopathy or bursitis of the gluteal tendons. Presents with lateral hip pain worse with walking, climbing stairs, or lying on the affected side. Treatment includes physical therapy and activity modification.

\synonyms
GTPS; Trochanteric Bursitis; Gluteal Tendinopathy

\medterm{Greater Vestibular Glands} Paired glands near the vaginal opening that secrete mucus into the vestibule, contributing to lubrication. They are also called Bartholin glands.

\synonyms
Bartholin Glands; Vulvovaginal Glands

\medterm{Great Saphenous Vein} The longest superficial vein of the lower limb, ascending from the medial foot to drain into the femoral vein at the saphenofemoral junction.

\medterm{Transposition of the Great Vessels} A cyanotic congenital heart defect resulting from embryological failure of the aorticopulmonary septum to spiral, such that the aorta arises directly from the right ventricle and the pulmonary artery originates from the left ventricle, creating two parallel, independent circulations requiring a mixing communication for survival.

\synonyms
TGA; D-Transposition of the Great Arteries; d-TGA

\medterm{Green Blindness} A form of red-green color-vision deficiency in which distinguishing green from certain other colors is difficult, usually because of an inherited cone-photopigment difference.

\medterm{Green Nail Syndrome} Green discoloration of one or more nails, usually caused by colonization or infection with Pseudomonas aeruginosa, often in an onycholytic or chronically moist nail.

\medterm{Greenstick Fracture} An incomplete fracture in which a young bone bends and cracks rather than breaking all the way through.

\medterm{Green Stool} Stool with a green color that may result from diet, food coloring, rapid intestinal transit, medicines, or infection.

\medterm{Grenz Zone} A narrow, uninvolved zone of normal collagenous dermis located immediately beneath the epidermis, separating the epidermal basement membrane from a dense underlying dermal cellular infiltrate, classically seen in lepromatous leprosy, granuloma faciale, and cutaneous lymphomas.
\synonyms
Subepidermal Grenz Zone; Clear Dermal Zone

\medterm{Grey Ankle Phenomenon} A gray or dusky discoloration around the ankle associated with chronic venous insufficiency and hemosiderin deposition. It occurs with other signs of venous disease.

\medterm{Grey Matter} Brain and spinal-cord tissue containing many nerve-cell bodies, connections, and supporting cells. It processes information and helps control movement, sensation, thinking, memory, and other functions.

\medterm{Grey Turner Sign} Bruising on the sides of the belly between the ribs and the hips, which is a sign of internal bleeding from severe pancreas inflammation or an injured blood vessel.

\medterm{Grey Turner's Sign} Bruising or blue-purple discoloration of the skin at the sides of the lower back. It can occur when bleeding deep inside the abdomen spreads through the tissues, including in severe pancreatitis.

\medterm{Grid Implant} A flexible electrode array placed on or near the brain or another tissue surface to record electrical
activity or deliver stimulation.

\medterm{Grief} The natural emotional response to loss, particularly the death of a loved one, but also following other significant losses such as divorce, job loss, or declining health. It encompasses a range of feelings including sadness, anger, guilt, and numbness, often occurring in waves rather than linear stages.

\medterm{Grief and Bereavement} The natural emotional response to loss, encompassing the process of coping with the death of a loved one or other significant losses. Includes grief (emotional response), loss (the event), and bereavement (the period of mourning). May involve physical, emotional, social, and spiritual dimensions.

\synonyms
Bereavement; Mourning; Complicated Grief

\medterm{Grimontia Hollisae} A salt-loving bacterium from seawater and seafood, formerly called \textit{Vibrio hollisae}. Eating contaminated raw or undercooked seafood can cause watery diarrhea, abdominal cramps, nausea, and fever; severe infection is uncommon but more likely in people with serious underlying illness.

\medterm{Griping} Cramping or spasmodic abdominal pain, often caused by contractions, gas, inflammation, or obstruction of a hollow digestive organ.

\medterm{Grip Strength} The force generated when the hand closes around an object, commonly measured with a dynamometer to assess hand function, muscle strength, nutrition, or rehabilitation progress.

\medterm{Griscelli Syndrome} A rare inherited disorder causing unusually light or silvery hair and skin pigmentation. Some forms also cause neurologic problems or severe immune overactivation, which can lead to fever, enlarged organs, low blood counts, and life-threatening illness.

\medterm{Griseofulvin} Griseofulvin is an oral antifungal that deposits in newly formed keratin and disrupts dermatophyte mitosis, treating infections of skin, hair, and nails; treatment is prolonged and hepatic or drug-interaction risks apply.

\medterm{Groin} The region where the lower abdomen meets the upper thigh, containing the inguinal canal and related muscles, vessels, nerves, and lymph nodes.

\medterm{Groin Lump} A palpable mass in the inguinal region, which may represent enlarged lymph nodes, a hernia, cyst, tumor, or another condition.

\synonyms
Inguinal Mass; Inguinal Lymphadenopathy

\medterm{Groin Pain} Pain in the inguinal region, which may originate from musculoskeletal, genitourinary, gastrointestinal, or vascular causes. Common causes include muscle strain, hernia, enlarged lymph nodes, and hip disease.

\synonyms
Inguinal Pain; Pubalgia

\medterm{Grommet Tube} A miniature, spool-shaped synthetic ventilation tube surgically inserted through an incision in the tympanic membrane (myringotomy) to equalize middle ear pressure and facilitate fluid drainage in chronic otitis media with effusion.

\synonyms
Tympanostomy Tube; Ventilation Tube

\medterm{Groove} A long, narrow depression or channel on a body structure, such as a groove on a bone for a tendon or a groove between folds of the brain. Its exact function depends on its location.

\medterm{Gross Examination} Inspection of a tissue or other specimen with the unaided eye before it is examined under a microscope. The pathologist records its size, shape, color, and visible abnormalities to help guide diagnosis.

\medterm{Gross Hematuria} Visible blood in the urine, appearing pink, red, or brown. Causes include urinary tract infections, kidney stones, bladder or kidney cancer, trauma, and disease of the kidney filters.

\synonyms
Macroscopic Hematuria; Visible Hematuria

\medterm{Gross Motor Control} The ability to use large muscle groups for coordinated movements such as walking, running, jumping, and maintaining posture. Develops progressively from infancy through childhood and is assessed as part of developmental milestones.

\synonyms
Large Motor Coordination; Fundamental Motor Skills

\medterm{Gross Motor Skills} Movements that use the large muscles for posture, balance, walking, running, jumping, and climbing. Children develop these skills at different rates.

\medterm{Ground-Glass Opacity} A hazy area on a chest CT scan in which underlying lung vessels remain partly visible. It is a radiologic finding, not a diagnosis, and may result from infection, inflammation, fluid, bleeding, scarring, or sometimes a tumor; its meaning depends on the pattern and clinical context.

\medterm{Ground Itch} An itchy skin rash caused by hookworm larvae entering the skin, usually after bare-foot contact with contaminated soil. It produces small, red, intensely itchy bumps.

\medterm{Ground Substance} The gel-like, noncellular part of connective tissue surrounding cells and fibers. It contains water, proteoglycans, and other molecules that support tissues, allow nutrient movement, and influence inflammation and repair.

\medterm{Group A Strep} A bacterium that commonly causes sore throat and skin infections. It can occasionally cause scarlet fever, pneumonia, severe soft-tissue infection, bloodstream infection, or immune complications affecting the heart or kidneys.

\medterm{Group A Streptococcal Infections} Infections caused by Streptococcus pyogenes, including pharyngitis, scarlet fever, impetigo, cellulitis, pneumonia, necrotizing soft-tissue infection, and bloodstream infection. Immune-mediated complications can include rheumatic fever and poststreptococcal glomerulonephritis.

\medterm{Group A Streptococcus} Streptococcus pyogenes, a bacterium that can cause pharyngitis, skin infections, invasive disease, and immune-mediated complications.

\medterm{Group B Strep} A bacterium that can live harmlessly in the bowel or genital tract but can cause serious infection in newborns during or soon after birth. It can also cause infection in pregnant people and vulnerable adults. Pregnancy screening and care during labor can reduce newborn risk.
\synonyms
Group B Streptococcus; Streptococcus Agalactiae

\medterm{Group B Strep Disease} Illness caused by group B Streptococcus bacteria, which can cause serious infections in newborns and infections in pregnant people or adults with medical risks.

\medterm{Group B Streptococcal Disease} Infection caused by group B Streptococcus bacteria, which can be serious in newborns and vulnerable adults.

\medterm{Group B Streptococcal Infection In Pregnancy} Maternal colonization or infection with Streptococcus agalactiae, which can be transmitted to the newborn during labor and cause serious early-onset infection. Screening and intrapartum management reduce neonatal risk.

\medterm{Group B Streptococcal Septicemia Of The Newborn} Serious bloodstream infection caused by Group B Streptococcus in neonates, typically acquired during birth. It can cause fever, lethargy, poor feeding, and breathing difficulty; treatment includes antibiotics and supportive care.

\synonyms
Early-Onset GBS Sepsis; Neonatal Streptococcus Agalactiae Sepsis

\medterm{Group B Streptococcus} A bacterium that may live harmlessly in the bowel or genital tract but can cause serious infection in newborns during or soon after birth. It can also cause infection in pregnant people and vulnerable adults; pregnancy screening helps reduce newborn risk.

\medterm{Group B Streptococcus in Pregnancy} Asymptomatic colonization of the maternal rectovaginal tract with \textit{Streptococcus agalactiae}. Routine culture screening is performed between 36 and 37 weeks of gestation so that colonized individuals can receive intrapartum intravenous antibiotic prophylaxis to prevent neonatal early-onset sepsis, pneumonia, and meningitis.

\synonyms
Maternal GBS Colonization; Intrapartum GBS Prophylaxis; Maternal Group B Strep

\synonyms
Maternal GBS Colonization; Intrapartum GBS Prophylaxis

\medterm{Group Practice} A healthcare organization in which several clinicians share facilities, staff, records, or administrative services. Patients may receive coordinated care from different professionals within the same practice.

\medterm{Group Therapy} A form of psychotherapy where a small group of individuals with similar issues meet to discuss, interact, and develop coping strategies under the guidance of a trained therapist. Used for depression, anxiety, substance use disorders, grief, and chronic illness. Benefits include peer support and shared experiences.

\medterm{Grover Disease} A transient or recurrent acantholytic dermatosis causing itchy red papules or papulovesicles, usually on the trunk of middle-aged or older adults; heat, sweating, and occlusion may aggravate it.

\medterm{Growing Pains} Recurrent, bilateral leg pain in children ages 3-12, typically occurring in the late
afternoon or evening and resolving by morning. Pain is deep, aching in the calves,
shins, and thighs (not joints). No limping, fever, swelling, or functional impairment.

\medterm{Growth} The process of increasing in physical size and development, including physical, cognitive, and emotional maturation. In medicine, growth is monitored using growth charts and anthropometric measurements to assess health and development in children.

\medterm{Growth Plate} A layer of cartilage near the end of a growing bone where the bone lengthens. It gradually closes after puberty.

\medterm{Growth Chart} A graphical representation of growth (weight, height/length, head circumference, BMI)
plotted against age, standardized by population (WHO growth standards 0-5 years; CDC
growth charts 2-20 years; percentiles/z-scores).

\medterm{Growth Cone} A specialized, moving tip at the end of a growing nerve fiber. It senses chemical and physical signals and guides the axon toward its target during nervous-system development or repair.

\medterm{Growth Disorder} A condition in which a child or adult grows unusually slowly, quickly, or unevenly because of genetic, hormonal, nutritional, chronic, or other medical causes.

\medterm{Growth Factor} A protein that sends signals controlling cell growth, division, specialization, survival, or repair. Different growth factors act on different tissues and are involved in wound healing, blood-vessel formation, immunity, and cancer.

\medterm{Growth Failure} Growth that is slower than expected or weight gain that is inadequate for age. It can result from poor nutrition, chronic illness, hormone problems, genetic conditions, or difficulty absorbing nutrients.

\medterm{Growth Faltering} A pattern of inadequate weight gain or growth in a child. It is the preferred modern term for many cases previously called failure to thrive and may result from feeding difficulty, inadequate intake, poor absorption, chronic disease, or other causes.

\medterm{Growth Hormone} A hormone released by the anterior pituitary that promotes growth partly through insulin-like growth factor 1 and also affects protein, fat, and carbohydrate metabolism.

\medterm{Growth Hormone Deficiency} Insufficient secretion of growth hormone from the anterior pituitary, causing short
stature in children and metabolic effects in adults. In children: delayed growth, short
stature, delayed bone age, and increased fat mass. In adults: decreased muscle mass,
increased fat mass, reduced bone density, fatigue, and dyslipidemia.

\medterm{Growth Hormone Physiology} Growth hormone is released in pulses by the pituitary gland, especially during sleep. It acts directly on tissues and also stimulates the liver to produce insulin-like growth factor 1, which supports bone and tissue growth.

\medterm{Growth Hormone-Releasing Hormone} A 44-amino-acid hypothalamic peptide hormone synthesized by neurons of the arcuate nucleus that travels via the hypophyseal portal system to bind G-protein coupled receptors on anterior pituitary somatotropes, stimulating growth hormone synthesis and pulsatile secretion.

\synonyms
GHRH; Somatoliberin

\medterm{Growth Hormone Stimulation Test} A test for suspected growth-hormone deficiency. A medicine or other controlled stimulus is given, and blood samples are taken over time to see whether the pituitary gland releases enough growth hormone. Results depend on the protocol and laboratory.

\synonyms
GH Provocation Test; Insulin Tolerance Test for GH

\medterm{Growth Hormone Suppression Test} A test for excess growth hormone, usually performed after drinking a glucose solution. In healthy people, glucose lowers growth-hormone levels; failure to suppress them supports a diagnosis such as acromegaly.

\synonyms
GH Suppression Test; Glucose Loading GH Assay

\medterm{Growth Hormone Test} Laboratory measurement of growth hormone levels in the blood, used to evaluate GH deficiency or excess. Random GH levels have limited utility; stimulation or suppression tests are more informative.

\synonyms
Serum Growth Hormone; Somatotropin Assay

\medterm{Growth Medium} A nutrient preparation supporting microbial multiplication in vitro, supplying carbon,
nitrogen, minerals, and growth factors in liquid or solid form.

\medterm{Growth Plate Fracture} A break involving the cartilage growth area of a child’s or adolescent’s bone. It can disturb future growth if displaced or untreated; imaging and specialist management may be used.

\synonyms
Physeal Fractures; Salter-Harris Fractures

\medterm{Growth Plate Fractures} Injuries involving the cartilage growth plates at the ends of growing children's bones, with possible effects on future bone growth.

\medterm{Growth Retardation} Older wording for growth that is slower than expected for age or developmental stage. Possible causes include inadequate nutrition, chronic illness, hormone or genetic disorders, poor absorption, and normal family variation; growth is assessed over time rather than by one measurement.

\medterm{Growth Scan} An ultrasound examination used to assess fetal size, growth pattern, amniotic fluid, and sometimes
blood flow. Measurements are compared with gestational-age standards to identify possible growth restriction or excessive
growth.

\medterm{Gtp Binding} GTP binding is the attachment of guanosine triphosphate to a protein, often switching that protein into an active state for cell signaling.

\medterm{Gtt.} An abbreviation from the Latin \textit{guttae}, meaning drops. It refers to a liquid dose or the number of drops.

\medterm{Guaifenesin} An expectorant used to loosen respiratory mucus and make coughing more productive. It provides
symptomatic relief in some respiratory illnesses; dosing depends on the formulation and patient.

\medterm{Guanfacine} Guanfacine is a centrally acting alpha-2A adrenergic agonist used for hypertension and, in extended-release form, ADHD; sedation, hypotension, bradycardia, and rebound hypertension after abrupt withdrawal are important risks.

\medterm{Guanine} Guanine is a purine nitrogenous base in DNA and RNA; it pairs with cytosine and is present in nucleotides such as GMP and GTP that support energy transfer and signaling.

\medterm{Guanosinetriphosphatase} An enzyme that hydrolyzes guanosine triphosphate to guanosine diphosphate and phosphate, regulating signaling and other cellular processes.

\medterm{Guanosine Triphosphate} A guanine nucleotide that provides energy and acts as a molecular switch in protein synthesis, cytoskeletal regulation, and G-protein signaling.

\medterm{Guanylate Cyclase} Guanylate cyclase is an enzyme that converts GTP to cyclic GMP, a second messenger involved in nitric-oxide signaling, smooth-muscle relaxation, phototransduction, and intestinal fluid secretion.

\medterm{Guarana} A plant whose seeds contain caffeine and other stimulants. It is used in energy products and supplements; excessive intake may cause anxiety, insomnia, rapid heartbeat, high blood pressure, or interactions with medicines.

\medterm{Guardianship} A legal arrangement in which an appointed person makes decisions for someone unable to make them independently.

\synonyms
Legal Guardianship; Conservatorship

\medterm{Guarding} Involuntary tightening of the abdominal muscles when the abdomen is examined, often reflecting irritation or inflammation of the abdominal lining.

\medterm{Guevedoces} A term used in the Dominican Republic for individuals with 46,XY differences in sex development caused by 5-alpha-reductase type 2 deficiency; they may have female-appearing or ambiguous external genitalia at birth and increased virilization at puberty.

\medterm{Guide Device} A tool used to direct, position, or stabilize another instrument during a procedure, such as a guidewire, drill guide, or needle guide. It helps ensure accurate placement and reduces injury to surrounding tissue.

\medterm{Guided Imagery} Guided imagery is a relaxation technique that uses imagined sights, sounds, or other sensations to promote calm and sometimes help manage pain or stress.

\medterm{Guided Tissue Regeneration} A dental procedure placing a barrier membrane to encourage regrowth of bone, root covering, and supporting ligament around a tooth. It may repair selected periodontal defects after inflammation is controlled.

\medterm{Guillain-Barre Syndrome} A rare disorder in which the immune system attacks peripheral nerves, often after an infection. Weakness and tingling usually begin in the legs and may spread upward; severe cases affect breathing or swallowing.

\synonyms
GBS; Acute Inflammatory Demyelinating Polyneuropathy (AIDP)

\medterm{Guillain-Barré Syndrome} An uncommon immune-mediated disorder in which peripheral nerves are attacked, usually causing rapidly progressive weakness and reduced reflexes. It can impair breathing or swallowing and is monitored closely in severe cases.

\medterm{Guinea Worm} A parasitic infection caused by Dracunculus medinensis, acquired by drinking water containing infected copepods and characterized by emergence of a long worm through the skin.

\medterm{Gulf War Syndrome} A term used for a group of persistent symptoms reported by some Gulf War veterans, such as fatigue, pain, cognitive difficulties, gastrointestinal symptoms, and respiratory problems. It is also called Gulf War illness or chronic multisymptom illness; causes are complex and treatment focuses on individual symptoms and function.

\medterm{Gullet} The esophagus, a muscular tube carrying swallowed food and liquid from the throat to the stomach. It moves contents by coordinated squeezing and is protected from injury by its moist inner lining.

\synonyms
Esophagus; Food Pipe

\medterm{Gum} The soft tissue, also called gingiva, that surrounds and protects the teeth and covers the jawbone. Healthy gums are usually firm and do not bleed with gentle brushing.

\medterm{Gum Biopsy} Histological examination of gingival tissue obtained by biopsy, used to diagnose gingival diseases, infections, or malignancies. May involve excisional or incisional biopsy techniques.

\synonyms
Gingival Biopsy; Periodontal Tissue Sampling

\medterm{Gum Disease} Inflammation or infection affecting the gums and tissues supporting the teeth. Early disease may cause bleeding and swelling, while advanced disease can destroy bone, loosen teeth, and cause tooth loss.

\synonyms
Periodontal Disease; Periodontitis

\medterm{Gumma} A soft area of tissue damage caused by late untreated syphilis. It can develop in the skin, bones, liver, or other organs and may cause a lump, pain, or organ damage.

\synonyms
Syphilitic Gumma; Tertiary Syphilis Granuloma

\medterm{Gummata} More than one gumma: soft inflammatory masses caused by late untreated syphilis. They can affect skin, bone, liver, or other organs and may damage tissue.

\medterm{Gummy Smile} A smile showing more gum than usual. It may result from lip movement, tooth eruption, jaw shape, or gum position; treatment is optional and depends on the cause and the person’s goals.

\medterm{Gum Problems} Gum problems include inflammation, bleeding, recession, ulcers, infection, trauma, and growths. Causes include plaque, tobacco, dry mouth, medicines, nutritional or immune disease, and poorly fitting appliances.

\medterm{Gums} The firm pink tissue, also called gingiva, surrounding and supporting the teeth. Healthy gums usually do not bleed with gentle brushing; swelling or bleeding can indicate gum disease.

\synonyms
Gingiva; Gingival Mucosa

\medterm{Gunn Sign} An ophthalmologic sign of hypertensive retinopathy in which an arteriosclerotic, rigid retinal arteriole crosses and compresses an underlying retinal venule, producing visible tapering and apparent deflection of the venule on both sides of the crossing.
\synonyms
Gunn's Sign; Arteriovenous Nicking

\medterm{Gunshot Wound} An injury caused by a firearm projectile. Damage may include bleeding, tissue destruction, fractures, nerve injury, organ damage, and infection.


\synonyms
Firearm Injury; Ballistic Trauma; Bullet Wound

\medterm{Gunther Disease} A severe autosomal recessive disorder of heme biosynthesis caused by a deficiency of uroporphyrinogen III synthase, leading to massive accumulation of isomer I porphyrins in tissues, manifesting with extreme cutaneous photosensitivity, blistering, mutilating scar formation, erythrodontia, and hemolytic anemia.
\synonyms
Congenital Erythropoietic Porphyria; CEP

\medterm{Gustation} The sense of taste. Special cells in taste buds detect sweet, salty, sour, bitter, and savory flavors. Most taste buds are on the tongue, with some in the soft palate and throat. Signals travel through several head nerves to the brain. Taste can change with infections, medicines, smoking, aging, or damage to the mouth or nerves.

\medterm{Gustatory} Relating to taste or the sense of taste, including taste buds, taste nerves, and the brain pathways that interpret flavors.

\medterm{Gustatory Lacrimation} Excessive tearing triggered by eating or smelling food, usually after an injury to the facial nerve or recovery from facial paralysis. The nerve signals for saliva and tears become incorrectly linked.

\medterm{Gustatory Sweating} Sweating triggered by eating or thinking about food, usually over the cheek or temple after injury or surgery involving the parotid region; it is also called Frey syndrome.

\medterm{Gut} Alternate informal term; see \textit{Digestive Tract}.

\medterm{Gut-Brain Axis} The two-way communication between the digestive tract and brain through nerves, hormones, immune signals, and substances made by gut microorganisms. It influences digestion, appetite, stress responses, and mood, but its clinical importance is still being studied.

\medterm{Gut Epithelium} The layer of cells lining the intestines. It absorbs nutrients and water, forms a barrier against harmful substances, and interacts with immune cells and gut microorganisms.

\medterm{Guthrie Test} An older name for a newborn blood-spot screening test, originally used to detect phenylketonuria. Modern screening covers more conditions; an abnormal result is followed by confirmatory testing.

\medterm{Gut Microbiome} The complex community of trillions of microorganisms, including bacteria, viruses, and
fungi, residing in the gastrointestinal tract. The gut microbiome plays a crucial role
in digestion, immune function, vitamin synthesis, and protection against pathogens.

\medterm{Guttae} Small wart-like bumps on the inner surface of the cornea, the clear front window of the eye. Numerous or severe guttae can impair fluid removal, causing corneal swelling, blurred vision, glare, or eye pain.

\medterm{Guttate Psoriasis} A form of psoriasis causing many small, scaly, drop-shaped spots, often after a streptococcal throat infection. It commonly affects children or young adults and may clear, recur, or later become plaque psoriasis.

\medterm{Guttural} Relating to the throat. A guttural sound is produced mainly in the throat or back of the mouth rather than at the front of the tongue and lips.

\medterm{Guyon Canal Syndrome} An entrapment neuropathy of the ulnar nerve as it passes through Guyon canal (the fibro-osseous ulnar tunnel at the wrist between the pisiform and hook of the hamate), causing sensory loss over the little finger and medial half of the ring finger and motor weakness of the hypothenar and interosseous muscles.
\synonyms
Ulnar Tunnel Syndrome; Handlebar Palsy

\medterm{Gymnast Wrist} A chronic overuse repetitive stress injury of the distal radial growth plate (physis) occurring in skeletally immature gymnasts and athletes who subject the upper extremities to intense axial loading and hyperextension. Radiographs typically demonstrate distal radial physeal widening, irregular metaphyseal margins, and cystic changes.

\synonyms
Distal Radial Physeal Stress Syndrome; Gymnast's Epiphysitis

\medterm{Gynaecologist} A medical doctor who specializes in the health of the female reproductive system
(vagina, cervix, uterus, fallopian tubes, ovaries) and breasts.

\medterm{Gynaecology} The medical specialty concerned with the female reproductive system, including the ovaries, fallopian tubes, uterus, cervix, vagina, and vulva, as well as related hormonal health.

\medterm{Gynaecomastia} Alternate name; see \textit{Gynecomastia}.

\medterm{Gynatresia} Absence of or blockage in an opening of the female reproductive tract, present from birth or acquired later. If menstrual blood cannot leave the body, it may cause absent periods, pelvic pain, or enlargement of the uterus or vagina.

\medterm{Gynecic} Relating to the female reproductive system or to gynecology.

\medterm{Gynecoid Pelvis} A commonly described pelvic shape with a rounded inlet, relatively broad dimensions, and a spacious outlet; it is associated with typical female pelvic anatomy but individual variation is substantial.

\medterm{Gynecological Examination} A gynecological examination assesses the female reproductive organs and may include inspection, a pelvic examination, and screening tests when appropriate.

\medterm{Gynecologist} A gynecologist is a physician who provides care for the female reproductive system, including preventive care, diagnosis, and treatment.

\medterm{Gynecology} The medical specialty concerned with the health of the female reproductive system
(uterus, ovaries, fallopian tubes, cervix, vagina, and vulva).

\medterm{Gynecomastia} A benign enlargement of male breast tissue caused by proliferation of glandular ductal components and stroma secondary to an increased ratio of estrogen to androgen activity. It commonly occurs as a physiological finding in neonates, pubertal adolescents, and older men. Pathological causes include primary or secondary hypogonadism (such as Klinefelter syndrome), liver cirrhosis, chronic kidney disease, estrogen-producing tumors, and medications such as spironolactone, antiandrogens, and anabolic steroids.

\synonyms
Gynaecomastia; Male Breast Enlargement

\medterm{Gynecomastia, Male} Alternate name; see \textit{Gynecomastia}.

\medterm{Gyrate Atrophy} An inherited retinal disorder causing gradually worsening peripheral and night vision loss. It results from abnormal ornithine metabolism and may be managed with dietary treatment, supplements, and regular eye monitoring.

\medterm{Gyrus} A raised fold on the surface of the brain, separated from neighboring folds by grooves called sulci. These folds increase the cerebral cortex’s surface area within the skull and contain areas responsible for different functions.
